main.tex

% !TeX program = xelatex
% =====================================================================
%  An Atlas of Human Cognition
%  Master book file
% =====================================================================

\documentclass[11pt,oneside]{book}

\usepackage[
  paperwidth=6in, paperheight=9in,
  inner=0.85in, outer=0.7in, top=0.9in, bottom=0.95in,
  headsep=12pt, footskip=24pt
]{geometry}

% ---- fonts (XeLaTeX) ----
\usepackage{fontspec}
\usepackage{unicode-math}
\setmainfont{TeX Gyre Pagella}
\setsansfont{TeX Gyre Heros}[Scale=0.92]
\setmonofont{Latin Modern Mono}[Scale=MatchLowercase, AutoFakeBold=2.2]
\setmathfont{Latin Modern Math}
\defaultfontfeatures{Ligatures=TeX}

\usepackage{microtype}
\usepackage{xcolor}
\definecolor{ink}{HTML}{1A1A1A}
\definecolor{slate}{HTML}{5A6472}
\definecolor{accent}{HTML}{8A2A2E}
\definecolor{rulegray}{HTML}{C7CCD1}

% ---- unicode symbols ----
\usepackage{newunicodechar}
\newunicodechar{→}{\ensuremath{\rightarrow}}
\newunicodechar{≈}{\ensuremath{\approx}}
\newunicodechar{×}{\ensuremath{\times}}

\setlength{\emergencystretch}{3em}

% ---- links ----
\usepackage{xurl}
\usepackage[hidelinks]{hyperref}
\hypersetup{
  pdftitle={An Atlas of Human Cognition},
  pdfauthor={Your Name},
  unicode=true
}

% ---- headings ----
\usepackage{titlesec}

\titleformat{name=\chapter,numberless}[display]
  {\normalfont}
  {}
  {0pt}
  {\Huge\bfseries\raggedright}

\titlespacing*{name=\chapter,numberless}{0pt}{4pt}{14pt}

% chapter dek / subtitle
\newcommand{\chapterdek}[1]{%
  \vspace{-6pt}%
  {\noindent\color{slate}\itshape\large #1\par}%
  \vspace{8pt}{\color{rulegray}\hrule height 0.5pt}\vspace{16pt}}

% reusable chapter command
\newcommand{\atlaschapter}[2]{%
  \chapter*{#1}%
  \markboth{#1}{#1}%
  \addcontentsline{toc}{chapter}{#1}%
  \chapterdek{#2}%
}

% section style
\titleformat{\section}
  {\normalfont\large\bfseries\color{ink}}{}{0pt}{}

\titlespacing*{\section}{0pt}{1.4\baselineskip}{0.5\baselineskip}

% aside block
\usepackage{quoting}
\newcommand{\aside}[1]{%
  \begin{quoting}[leftmargin=1.4em,rightmargin=1.4em,vskip=6pt]
  {\itshape\color{slate} #1}\par
  \end{quoting}}

% ---- running heads / folios ----
\usepackage{fancyhdr}
\pagestyle{fancy}
\fancyhf{}
\fancyhead[C]{\small\itshape\nouppercase{\leftmark}}
\fancyfoot[C]{\thepage}
\renewcommand{\headrulewidth}{0pt}

\fancypagestyle{plain}{%
  \fancyhf{}%
  \fancyfoot[C]{\thepage}%
  \renewcommand{\headrulewidth}{0pt}%
}

% ---- prose ----
\linespread{1.06}
\setlength{\parindent}{1.2em}
\setcounter{secnumdepth}{0}

\begin{document}

\frontmatter

\begin{titlepage}
\centering
\vspace*{2in}

{\Huge\bfseries An Atlas of Human Cognition\par}

\vspace{0.5in}

{\Large Joshua G. Stern\par}

\vfill

{\small Copyright \textcopyright{} 2026 Joshua G. Stern\par}

\vspace{0.5\baselineskip}

{\small\itshape The writing process for this book made extensive use of proprietary frontier large language models.\par}

\end{titlepage}

\tableofcontents

\mainmatter

\include{chapters/ch1-memory}
\include{chapters/ch2-working-memory}
\include{chapters/ch3 pain}
\include{chapters/ch4}
\include{chapters/ch5}
\include{chapters/ch6-perception}
\include{chapters/ch7}
\include{chapters/ch8}
\include{chapters/ch9}
\include{chapters/ch10}
\include{chapters/ch11}
\include{chapters/ch12}
\include{chapters/ch13-self and other}
\include{chapters/ch14-the minds eye}

\backmatter

% Optional:
% \chapter*{Acknowledgments}
% \addcontentsline{toc}{chapter}{Acknowledgments}
% Your acknowledgments here.

\end{document}



ch1-memory.tex
\atlaschapter{Memory}{Why the thing you call ``my memory'' is not one thing, and not a recording, and why the experts can't agree on how to draw its map.}

Ask someone where a memory is and they will usually point, vaguely, at
the back of their own head, as though there were a small dark room in
there with shelves. The metaphor is ancient and almost irresistible.
Memory is a \emph{store}; remembering is \emph{retrieval}; forgetting is
a misfiled folder or a faded page. We talk about having ``a good
memory'' the way we talk about having a big hard drive --- as if the
whole apparatus were a single container of variable size, and the only
interesting question were how much fits in it.

Much of that picture is wrong. Memory is not one container
but several, running on different machinery and obeying different rules.
It is not a recording that plays back what was filed; it is a
reconstruction, rebuilt a little differently each time you ask for it.
And the act of remembering, far from leaving the stored item untouched,
can partially rewrite it. The folk metaphor gets memory backwards: it
imagines a passive archive, when what we actually carry is an opinionated
editor.

This chapter is about how science took that single, intuitive ``memory''
apart. When you line up the field's major reference works, the catalogs and
ontologies that different research traditions use to organize cognition,
they agree that memory matters and disagree about almost everything
else: how many memories there are, what contains what, even whether
``memory'' names a single faculty at all. One catalog hangs everything
beneath a single grand heading called \textsc{Memory}. Another has no
such heading anywhere and treats short-term and long-term remembering as
two unrelated abilities that happen to share a syllable. A third refuses
to call memory a category of \emph{anything} and instead wires it in as
a component, a box with cables running to other boxes.

\section{The man who lived in the present}

The case that broke the single-store picture was that of a man
from Connecticut named Henry Molaison, known for half a century in the
literature only by his initials, H.M. As a young man he suffered
seizures so severe and so frequent that in 1953, aged twenty-seven, he
agreed to an operation that a surgeon named William Scoville himself
called frankly experimental. Scoville removed a chunk of tissue from
deep in both temporal lobes, including most of a pair of
seahorse-shaped structures called the hippocampi. The seizures
improved. Something else went catastrophically wrong.

H.M. could no longer form new memories of facts or events. He could
hold a thought for as long as he kept rehearsing it --- a phone number,
a sentence --- but the moment his attention moved on, it was gone, with
no trace that it had ever been there. He met the same researchers
hundreds of times across decades and greeted each visit as the first.
He could not learn that his father had died, and grieved the news
freshly whenever it surfaced. His intelligence was intact; his
personality, his vocabulary, his childhood memories, all intact. What
had been destroyed was specifically the machinery that lays down new
conscious memories, and it had been destroyed cleanly, leaving the rest
standing. He lived, as one of his examiners put it, in a permanent
present tense.

The neuropsychologist Brenda
Milner asked H.M. to trace the outline of a star while looking only at
his hand and the star reflected in a mirror --- a maddening task,
because the mirror reverses every correction your hand tries to make.
People are clumsy at it on the first day and smoother by the third. So
was H.M. Day after day his tracing improved, his hand learning the
trick of the reversed feedback --- and day after day he insisted,
sincerely, that he had never attempted anything of the kind in his life.
His body was accumulating a skill that his conscious mind had no access
to whatsoever.

There it was, in a single person: not one memory but at least two,
dissociated so sharply that one could be obliterated while the other ran
on untouched. The kind of memory you can declare --- \emph{I know that,
I remember when} --- and the kind you can only demonstrate ---
\emph{watch, I can do this}. Researchers came to call them
\emph{declarative} and \emph{nondeclarative}, or explicit and implicit.
H.M. had lost the first and kept the second, which meant they could not
be the same system wearing two hats. They lived in different tissue.

\aside{The hippocampus, H.M. proved, is not where memories are kept. He
had lost his and could still recall his youth perfectly well. It is
where new declarative memories are \emph{made} --- a factory, not a
warehouse. Damage the factory and the existing stock survives; you
simply can't manufacture anything new.}

H.M. died in 2008, having contributed immensely to the science of memory
without ever being able to remember that he had done so. The lesson his case taught is the foundation everything else in
this chapter is built on: ``memory,'' the single word, names a committee,
not an individual.

\section{Two ways to remember a croissant}

If declarative memory is the kind you can put into words, it splits
again as soon as you look closely --- and noticing the split was the
life's work of an Estonian-Canadian psychologist named Endel Tulving,
who died in 2023 and whom many of his colleagues considered the most
important memory researcher of the twentieth century. One of them
remarked that there is memory before Tulving and memory after Tulving.

The distinction he drew, in 1972, sounds almost too simple. There is the memory of
\emph{facts} --- that Paris is the capital of France, that a year has
twelve months, that croissants are flaky --- which Tulving called
\emph{semantic} memory. And there is the memory of \emph{episodes} ---
the specific croissant you ate on a specific morning, the table you sat
at, the person across from you, the quality of the light. That he called
\emph{episodic} memory. Knowing where
people generally leave their cars is semantic; remembering where you
left yours this morning is episodic. The first is knowledge that has
come loose from the occasion of learning it. The second is bolted to a
particular moment of your own life.

What makes episodic memory genuinely peculiar --- and what Tulving spent
decades arguing --- is that retrieving one is a kind of \emph{mental
time travel}. To remember this morning's croissant is not merely to look
up a fact about your past; it is to re-inhabit the moment, to project
yourself backward and experience a faint replay from the inside.
Tulving gave this self-locating, time-traveling quality of awareness a
name, \emph{autonoetic} --- self-knowing --- consciousness, and
contrasted it with the flatter, merely \emph{noetic} ``knowing'' that
accompanies a bare semantic fact. You don't time-travel to recall the
capital of France. You just know it. The same reflective machinery, he
suspected, is what lets us project \emph{forward} as well: the capacity
to remember a specific past and the capacity to imagine a specific
future may be two uses of one faculty, which is why people with certain
kinds of amnesia struggle not only to recall their lives but to picture
their tomorrows.

Stack these distinctions up and you get the standard tree. Memory at the
top; under it, declarative and nondeclarative; under declarative,
episodic and semantic; under nondeclarative, the family of skills and
habits and reflexes that H.M.'s hand could still learn. It is a clean,
satisfying taxonomy, a Linnaean classification of the kinds of
remembering. It is also only one of two completely different ways to
carve the territory.

\section{A taxonomy, or a machine?}

The tree just described answers a particular question: \emph{what kinds
of things do we remember?} It sorts memory the way a naturalist sorts
beetles, by type. Episodic and semantic are \emph{kinds of} declarative
memory the way a beetle is a kind of insect. This is the carving you
find in the catalogs built by psychologists and clinicians and
lexicographers --- people whose job is to name and classify constructs.
Ask such a catalog ``what is short-term memory?'' and it answers ``a
kind of memory,'' and points upward to the parent heading.

But there is a second question you can ask, which sounds similar and is
not: \emph{what is the machinery that does the remembering, and how is
it wired together?} This question doesn't sort memory into kinds. It
lays out a process and a set of components. Information comes in and is
\emph{encoded}. Over hours and days it is \emph{consolidated} --- made
durable, gradually shifted from the fragile hippocampal factory toward
more permanent storage spread across the cortex. Later it is
\emph{retrieved}, which, as we'll see, is itself an act of partial
reconstruction. And alongside this pipeline sits a separate small
workspace --- \emph{working memory} --- where you hold and manipulate a
few items right now, the mental desktop on which a phone number or the
first half of this sentence is briefly laid out.

In this second picture, the parts of memory are not \emph{kinds} of
anything. They are \emph{components} --- stages in a process, modules in
an architecture, boxes with cables. This is the carving you find in the
catalogs built by people who model cognition as a machine: the
cognitive architectures that try to write running computer simulations
of a mind, and the clinical frameworks that try to tie remembering to
specific brain circuits. To them, working memory is not ``a kind of
memory.'' It is a \emph{place} --- a buffer wired to the rest of the
system --- and the long-term store is another place, and the
hippocampus is the indexer that files things between them.

The two paradigms nest memory in
\emph{opposite directions}. The taxonomist's tree puts one big
\textsc{Memory} at the root and makes everything else a subtype
\emph{beneath} it: unity first, divisions below. The architect's diagram
has no such root; it has a working-memory box and a long-term box and a
hippocampal indexer, each a \emph{part} of a larger cognitive system,
with no overarching ``Memory'' faculty sitting above them at all. And a
third tradition --- the psychometric one, descended from the study of
how people differ in measurable ability --- does something stranger
still: it declines to posit any general ``memory'' and instead treats
short-term and long-term remembering as two coordinate \emph{abilities},
siblings hanging off general intelligence, with the very split that the
taxonomist tucks near the bottom of the tree hoisted right to the top.

\aside{Ask the field's catalogs the simplest possible structural
question --- \emph{what contains what?} --- and they answer
incompatibly. One says memory is a category you subdivide. One says it
is a mechanism you wire. One says it isn't a single thing to be
contained at all. None of them is wrong. They are answering different
questions and mistaking each other for rivals.}

A \emph{kind-of} tree and a \emph{part-of} wiring diagram are both
legitimate, and they cannot be merged into one because they are not the
same kind of map. You cannot ask whether ``episodic memory'' sits above
or below ``the hippocampus,'' any more than you can ask whether
\emph{sonata} sits above or below \emph{violin}. One is a kind of thing;
the other is a piece of the apparatus. The reason the catalogs nest
memory in opposite directions is that some of them are classifying
\emph{what we remember} and others are diagramming \emph{the organ that
remembers}, and those are different subjects that happen to share a
word.

Almost every famous ``debate'' about memory turns
out to have this shape: not a contest over facts, but two communities
carving the same territory along different axes for different and
perfectly good reasons. The psychometricians want something they can
\emph{measure} across a thousand people, so they cut memory into
factors that vary between individuals. The architects want something
they can \emph{run} as a simulation, so they cut it into modules that
pass data. The clinicians want something they can \emph{lesion or treat},
so they cut it into circuits. Each cut is real. None is the whole shape.

\section{Memory is not a recording}

So much for the structure. Now for the most counterintuitive thing the
science has to say about the \emph{content} --- the thing that the
hard-drive metaphor gets most dangerously wrong.

A memory is not played back. It is rebuilt.

The first person to demonstrate this rigorously was a British
psychologist, Frederic Bartlett, who in the 1930s had people read an
eerie, culturally unfamiliar Native American folk tale and then retell
it, again and again, over days and weeks. They did not retain it and
slowly lose pieces, the way a recording degrades. They \emph{remodeled}
it. The strange elements quietly went missing; the gaps filled with
ordinary, expected detail; the whole thing drifted toward a tidier,
more familiar story that fit the teller's own expectations. People were
not recalling the tale. They were reconstructing something plausible
from a few surviving fragments plus a large helping of what \emph{ought}
to have been there. Bartlett's word for the expectation-shaped templates
doing the filling was \emph{schema}, and the insight has never been
overturned: remembering is an act of construction, and construction
runs on assumption.

Half a century later, the psychologist Elizabeth Loftus turned this from
a curiosity into something with consequences in courtrooms. In a famous
study, she showed people a film of a car accident and then asked about
it using subtly different words. People asked how fast the cars were
going when they \emph{smashed} into each other gave higher estimates
than those asked when the cars \emph{hit} --- and a week later, the
``smashed'' group was markedly more likely to remember broken glass at
the scene. There had been no broken glass. The leading word, slipped in
\emph{after} the event, had become part of the memory of the event. In
later work Loftus and others went further and implanted whole episodes
that had never happened --- convincing a sizeable fraction of subjects
that they had, as children, been lost in a shopping mall, or met a
costumed character at a theme park that had never employed one --- not
by hypnosis or trickery, but simply by suggesting the false episode
confidently and letting the mind's own reconstructive machinery flesh
it out. The subjects were not lying. They genuinely remembered, often
in vivid and emotional detail, things that had not occurred.

The unsettling implication is that the \emph{vividness} of a memory, its
richness and emotional force and the absolute conviction that
accompanies it, is no guarantee whatsoever of its accuracy. A confident,
detailed, deeply felt memory can be substantially or entirely
manufactured. The cognitive scientist Daniel Schacter catalogued the
characteristic ways memory misleads us --- its forgettings, its
distortions, its suggestibility, its habit of misattributing where a
piece of information came from --- and made the case that these are not
bugs to be engineered away but the inevitable price of a system built
for usefulness rather than fidelity. A memory optimized to help you act
sensibly in the future has little reason to preserve the past with a
court reporter's accuracy, and ours does not.

\section{Editing the trace}

If memory is reconstructed, you might ask what, physically, is being
reconstructed \emph{from}. What is the trace --- the thing in the tissue
that a remembering reaches for? Researchers have a old, slightly
romantic word for it: the \emph{engram}, the physical imprint a
particular memory leaves in the brain. For most of the twentieth century
the engram was almost a ghost. Karl Lashley
spent thirty years cutting into the brains of trained rats trying to
find where a specific memory lived, and concluded, more or less in
despair, that it didn't seem to live anywhere in particular --- it was
maddeningly distributed, resistant to being pinned to a spot.

In the last fifteen years that ghost has been caught --- in mice, and
with a technology that sounds like science fiction. Using a technique called
optogenetics, researchers at MIT and elsewhere learned to tag the
specific population of neurons that switched on while a mouse formed a
particular memory, and to install in just those cells a molecular switch
that makes them fire in response to a pulse of light. They could then,
in effect, find a memory and \emph{press play} on it. Shine the light,
reactivate that exact set of cells, and the mouse behaves as though
recalling the original event --- freezing in fear, for instance, at a
place that is presently perfectly safe, because the cells encoding a
frightening place have been switched back on.

Next, they reactivated
the cells encoding a safe, neutral location at the same moment they
delivered a small shock in a \emph{different} location. The mouse
afterward feared the safe place --- a place where nothing bad had ever
happened to it --- because a fear had been welded onto an artificially
reactivated memory of it. They had implanted, at the level of identified
cells, a false memory. The engram, far from being a fixed inscription,
turned out to be something you could not only locate and trigger but
\emph{edit}. (More recent work suggests the trace is even more fluid than
that first wave of experiments implied: the exact cells making up an
engram appear to shift over hours and days after learning, neurons
joining and dropping out, so that even the physical memory is less a
fixed photograph than a slowly reshuffled committee.) Such
\emph{reconsolidation} is one of the most provocative ideas in modern
memory science. For a long time the
assumption was that a memory, once consolidated --- once made durable in
the days after learning --- was essentially set, like cured concrete. In
2000 a team studying fear in rats showed otherwise. When you
\emph{retrieve} a consolidated memory, it appears to become briefly
fragile again, editable, as though pulling the file out of the cabinet
unlocks it --- and during that window, before it is re-stored, it can be
weakened or altered. Retrieval is not a read-only operation. To remember
is, sometimes, to overwrite.

The clinical dream this set off is obvious and seductive: if a traumatic
memory becomes briefly editable each time it is recalled, perhaps you
could have a patient recall the trauma under a drug --- propranolol, a
common blood-pressure medication that blunts the stress chemistry that
normally stamps emotional memories in hard --- and let the memory
re-store in a softened, less harrowing form. The memory of the event
would remain; its capacity to ambush and overwhelm would be dialed down.
Early trials were exciting. The honest current picture is more sober.
Reconsolidation itself is a real and well-replicated laboratory
phenomenon. Whether propranolol meaningfully relieves post-traumatic
stress disorder in actual patients is genuinely unsettled: some careful
trials show a benefit, others find the drug does no better than a
placebo --- and the placebo, the mere structured act of revisiting the
memory in a safe setting, often helps a great deal on its own.
Reviews of the evidence disagree with each other. It is a frontier,
not a cure, and a useful reminder that a beautiful mechanism in a mouse
is a long way from a treatment in a person.

\section{Why the maps disagree}

Return, finally, to the puzzle we started with: the field's reference
works, lined up side by side, nest memory in incompatible directions and
can't agree on how many memories there even are. Having walked through
what is actually known, the disagreement should look different now ---
less like a mess to be cleaned up and more like an honest reflection of
the subject.

Memory genuinely is several things. H.M.'s mirror-tracing hand proved
that the declarable and the merely demonstrable run on separate
machinery. Tulving's croissant proved that within the declarable, the
fact and the episode are different in kind, the second carrying a
strange time-traveling awareness the first lacks. So any catalog that
splits memory into multiple systems is right to. But memory genuinely
\emph{is} also a single integrated process --- encode, consolidate,
retrieve --- with shared components and a common currency, so any
catalog that models it as one wired machine is also right to. And it
genuinely \emph{does} vary between people along measurable axes, so the
psychometrician who carves it into abilities is right too. The taxonomy,
the architecture, and the ability-profile are three faithful maps of one
territory, drawn for three purposes, and they nest memory in different
directions for the same reason a subway map and a street map and a
geological survey of the same city refuse to overlay. They are not in
conflict. They are in different projections.

What the disagreement is really recording, then, is not confusion about
memory but the irreducible richness of it --- the fact that a thing this
central to being a person cannot be captured by a single cut. The
intuitive picture we started with, the dark room with shelves, fails not
because it is too simple but because it is the \emph{wrong genre} of
thing: a static container standing in for a process that is plural,
reconstructive, and quietly self-revising every time you use it. You do
not retrieve your past from storage. You rebuild it, from fragments and
expectations, in the present, and in the rebuilding you change it a
little --- which is why the people we have been longest and remember
most often are, in a sense, the ones we have edited most.





ch2-working-memory.tex
\atlaschapter{Working Memory}{The mental desktop: what it holds, how little that is, and why fifty years of research still can't agree on whether the thing is a place or a process.}

Right now, as you read this sentence, something is keeping the beginning
of it alive long enough for the end to make sense. The same something is
holding, somewhere in the background, whatever you were thinking about
before you started reading --- a task, a name, a turn you need to take
in twenty minutes --- without letting it vanish entirely but also without
letting it crowd out the words in front of you. And if someone were to
call your name at this moment, a small part of your attention would
redirect toward the sound while the rest tried to keep its grip on the
sentence.

All of that happens continuously, effortlessly enough that you
ordinarily don't notice it. The name for the cognitive machinery doing
it is \emph{working memory}: the system that holds a small amount of
information in an active, immediately usable state while you do
something with or around it. Not storage, not the long-term filing
cabinet where your childhood memories and the name of the capital of
France are quietly kept. The immediate, lit-up workspace where current
thinking happens.

Working memory is probably the most-studied concept in cognitive
psychology and one of the most consequential. Individual differences in
it predict performance on a striking range of tasks --- reading
comprehension, reasoning, managing distractions, learning a new
language --- and its decline is one of the earliest and most disruptive
features of aging and dementia. A large industry of commercial software
is built on the promise of training it.

It is also, underneath the consensus on its importance, a concept that
the field's major theories describe in genuinely incompatible terms. Not
about what it does --- on that there is broad agreement --- but about
what it \emph{is}. Whether it is a place, a set of specialized storage
buffers where particular kinds of information sit. Or a process, an
attentional spotlight that temporarily illuminates things already in
long-term memory. Or something that resists being called either. The
models agree on the data. They disagree about the data's meaning.

\section{A stress test for the single store}

The story begins with a puzzle noticed in the early 1970s. The received
picture of short-term memory at the time was simple: a single, limited
container that held incoming information briefly before passing it along
to long-term storage. One container, one function. Two British
psychologists, Alan Baddeley and Graham Hitch, decided to stress-test
it. Their reasoning was direct: if short-term memory is one undivided
thing, loading it heavily should impair any other task that also needs
it.

So they gave people a string of digits to hold in mind --- six or seven
items, enough to fill the supposed container --- and simultaneously
asked them to do something cognitively demanding: comprehend a sentence,
reason through a logical problem, or track a concept through a paragraph
of text. If the single-store model was right, performance on the
secondary tasks should have collapsed. The container was already full.

It didn't collapse. People could do both things at once, imperfectly but
far better than the single-store model predicted. Holding a digit string
didn't gut language comprehension. The two tasks were drawing on
different resources that could, within limits, run in parallel. Short-term
memory, Baddeley and Hitch concluded, was not one container but a
system with multiple components --- and they called the whole system
\emph{working memory}.

\section{The three-room model}

The model they built from that starting point has proved extraordinarily
durable. In its original 1974 form it had three parts; a fourth was
added in 2000. The architecture is worth understanding because it
captures something genuinely real about how different kinds of
information interfere with each other.

The first component is the \emph{phonological loop} --- the inner voice.
It is a short-term store for sound-based information, particularly
speech and language, coupled with a silent rehearsal process: the
subvocal muttering that lets you hold a phone number by repeating it
to yourself. The loop's capacity is not measured in items but in time:
roughly two seconds of speech. Words that take longer to say are harder
to hold, and Welsh speakers, whose words tend to be longer than English
ones, have shorter digit spans in their native language than English
speakers do. If you suppress the rehearsal mechanism --- by making
someone repeat a meaningless syllable aloud, ``the the the the,'' while
trying to hold a list --- immediate verbal recall collapses, because the
mental rehearsal and the mouth are competing for the same articulatory
machinery.

The second component is the \emph{visuospatial sketchpad} --- the inner
eye. It handles visual and spatial information: the shape of an object,
the layout of a room, the sequence of moves in a chess position. It is
measurably separate from the phonological loop. You can load both
simultaneously with much less interference than loading either one
twice, because verbal and visuospatial material draw on different
machinery. A person trying to hold a mental image while repeating digits
aloud suffers far less than someone trying to hold two mental images or
repeat two strings of digits, because the first pairing splits the load
between systems that do not substantially compete.

The third, most important, and most underspecified component is the
\emph{central executive}: the attentional supervisor that coordinates
the other two, manages the switching of attention between tasks,
suppresses irrelevant information, and decides what gets processed
and when. Baddeley described it as the most important part of the
model and the one he understood least. It is in some ways a placeholder
--- a label for the cognitive work that cannot be assigned to either
of the specialized stores.

In 2000 Baddeley added a fourth part, the \emph{episodic buffer}, to
solve a problem the original model couldn't: how does information from
the two separate stores combine into a single coherent experience? When
you understand a sentence, you bind the sounds of the words (phonological
loop), any images the sentence evokes (visuospatial sketchpad), and
everything you already know about the topic (long-term memory) into a
single unified meaning. The episodic buffer is Baddeley's name for the
temporary workspace where that binding happens --- a limited-capacity
stage on which the other systems' outputs are assembled into scenes.

\aside{The clearest evidence that the loop and sketchpad are genuinely
separate systems, not one system wearing two hats, comes from brain
damage. Some patients lose verbal short-term memory while their
visuospatial ability remains intact; others show the reverse. The
systems can be knocked out independently because they live in different
tissue.}

The model is elegant and has generated fifty years of productive
research. It is also, as every student of cognition eventually discovers,
a partial answer to the question it seems to settle.

\section{How much fits}

Before pressing on the model's foundations, it is worth dwelling on the
question of capacity --- of how little the workspace actually holds ---
because the answer is more interesting than it first appears.

The number most people half-remember from a psychology class is seven.
In 1956 the psychologist George Miller published a paper with the title
``The Magical Number Seven, Plus or Minus Two,'' one of the most
cited papers in the history of psychology. Miller had noticed that
across a range of tasks --- remembering a list, making absolute
judgments about stimuli, categorizing sounds --- human performance
seemed to top out at somewhere between five and nine items. He
proposed seven as the central tendency and called it the channel
capacity of short-term processing.

What Miller did not say, but later had to clarify, is that each
``item'' in his count could be almost arbitrarily complex, because the
mind has a trick called \emph{chunking}: recoding several units of
information into one. A phone number is ten digits, but held as three
groups it becomes three chunks. A chess position is thirty-two pieces,
but for a master it resolves into a handful of familiar patterns, each
holding as a single chunk. Expertise is, in large part, the
accumulation of ever-richer chunks --- which is why the master can
glance at a board mid-game and reconstruct it from memory while a
beginner can't recall even five pieces. The capacity didn't change.
What counts as one unit did.

Later research narrowed the number further. The psychologist Nelson
Cowan, reviewing a wide body of experiments in 2001, argued that the
true limit of working memory --- when rehearsal is blocked and chunking
via long-term knowledge is controlled --- is closer to four items, plus
or minus one. Not seven. The two estimates aren't actually in
contradiction: seven chunks are possible when language and expertise
are available to help pack them; four is what remains when all those
props are removed and you are left with the bare attentional capacity.

Four. Stripped down, that is the size of the mental desktop. It is a
sobering number for a species that finds its own intelligence
impressive.

\aside{There is a reason phone numbers are broken into groups, why
postal codes mix letters and digits, why a social security number comes
in three chunks rather than nine consecutive digits. These conventions
are all accidental empirical discoveries: formats that people found
easier to hold and pass along, without knowing they were calibrating
to the four-item limit of the attentional focus.}

\section{A place, or a process?}

Here is where the field divides, and where it is worth spending some
time, because the divide is not cosmetic.

Baddeley's model treats working memory as a \emph{place}, or rather a
set of places: dedicated storage buffers with their own architecture,
their own neural tissue, their own capacity limits. The phonological
loop is a place where verbal information sits while it waits to be
used. The sketchpad is a place where images sit. The episodic buffer is
a place where bound representations sit. These are distinct from
long-term memory and from the rest of cognition; they are a
specialised system.

The psychologist Nelson Cowan had a different reading of the evidence.
In his \emph{embedded-processes} model, there is no separate working
memory system. What we call working memory is simply the part of
long-term memory that is currently \emph{activated} --- called up from
storage and held in a heightened state of availability --- with a
subset of that activated material sitting inside a focus of attention.
Working memory is not a room adjacent to long-term memory. It is a
spotlight shining on part of it. The same information that lives in
long-term memory is temporarily lit; when the spotlight moves, it goes
dark again. There is no separate buffer, no transfer of content from
one container to another. There is one system, with different regions
of it more or less active at any moment.

Randall Engle at Georgia Tech pushed the attentional reading further
still. On his account, what working memory \emph{capacity} actually
measures is not the size of any buffer but the ability to control
attention --- to maintain goal-relevant information in the face of
interference and distraction, to hold on to a thought when something
is trying to knock it loose. The individual differences that make
working memory interesting as a predictor of intelligence and academic
performance are, in his view, differences in this attentional control,
not differences in the size of a storage tank.

\aside{The three positions are not disputes about the data. They are
disputes about what the data is a picture \emph{of}. Baddeley sees
specialized buffers. Cowan sees activated long-term memory with an
attentional focus. Engle sees attentional control, full stop. All
three frameworks make predictions that the experiments broadly
support. The data does not pick a winner. It is, in this respect, a
clean example of the same thing that runs through every chapter of
this atlas: identical evidence, different questions behind it, and
therefore different answers that cannot be adjudicated by adding
more data of the same kind.}

It is tempting to ask which view is right, and the honest answer
is that the field does not know. What is clear is that the three
models are not simply the same claim in different words. Baddeley's
buffers predict patterns of interference and dissociation that
Cowan's activated-long-term-memory framework accounts for differently,
and what each predicts about aging, individual differences, and clinical
populations varies in subtle but real ways. The debate is alive and
the experiments continue.

\section{What the neurons do}

Brain imaging arrived as a hoped-for arbiter. If working memory is a
dedicated system with its own neural tissue, it should light up in
ways that distinguish it cleanly from long-term memory and from
attentional control. The results have been informative without being
decisive.

The most consistent finding, established first in monkeys and then
confirmed in humans, is this: during the delay in a working-memory
task, the prefrontal cortex keeps firing. Show a monkey a spot of light,
then turn the light off and impose a delay of several seconds, then ask
the monkey to look at where the light was. Individual neurons in the
dorsolateral prefrontal cortex remain active throughout the dark
interval, as if the light were still there --- bridging the gap
between seeing and responding, holding the location in mind through
the period of sensory absence. This \emph{delay-period activity}
was identified by Patricia Goldman-Rakic at Yale and has been
replicated across dozens of laboratories. It is the closest thing
neuroscience has found to a neural signature of the ``holding''
function that working memory performs.

The question the finding leaves open is exactly the question the
behavioral debate left open. Is the prefrontal cortex holding the
content itself --- storing the location of the dot like a temporary
visual buffer? Or is it doing the attentional work of keeping that
representation from fading, the control process that Engle's framework
emphasizes? More recent imaging work suggests that the actual
representation of what is being held --- the sensory content of the
memory --- often lives not in the prefrontal cortex but in the same
posterior sensory areas that originally processed the stimulus. The
prefrontal cortex may be less the storage cabinet and more the
thing that keeps the cabinet from closing. That interpretation favors
Cowan and Engle over Baddeley, but the debate at the neural level
mirrors the debate at the behavioral level, and it is not settled.

\section{Can it be expanded}

Working memory's four-item limit has made it a target for intervention.
If it predicts intelligence, academic success, and resilience to
distraction, and if it declines with age and is disrupted in conditions
from ADHD to schizophrenia, the practical question writes itself: can
you train it?

The answer, arrived at after two decades of often over-hyped research
and extensive meta-analysis, is nuanced in a way that is worth
stating precisely. You can train working memory tasks. Give people
several weeks of practice on the \emph{n-back} task --- in which you
must track whether the current item in a sequence matches the one that
appeared some fixed number of items ago --- and they get better at the
n-back task. Give them the Cogmed program, perhaps the most
commercially successful working memory training system, and they improve
on other unpracticed working memory tests. This is called \emph{near
transfer}: improvement on tasks closely related to the training task.

What the meta-analyses have, with increasing consistency, failed to
find is \emph{far transfer}: improvement on the general reasoning,
reading comprehension, and fluid intelligence that working memory
capacity was supposed to predict in the first place. Training the task
does not seem to train the underlying capacity. People get better at
the drills without getting better at the thing the drills were supposed
to tap.

This finding cuts interestingly across the theoretical divide. If
working memory were a fixed-capacity buffer, you would not expect
training to expand it; the limits are architectural, not learned. If
it is attentional control, you might expect that control to be
improvable with practice. The fact that training produces near but
not far transfer has been interpreted as evidence against the
architectural buffer view (because even the near transfer is
unexpected if the capacity is fixed) and against the attentional
control view (because if that control is trainable, why doesn't the
training generalize?). The data is equally inconvenient for both
theories, which is perhaps the most honest summary of where things
stand.

What does seem to expand working memory's effective reach, without
expanding its raw capacity, is the thing it has always depended on:
knowledge. Because the four-item limit applies to chunks, and because
what counts as a chunk is determined by what you already know, the
person with the richer long-term memory can hold more in working
memory without holding more items. The chess master and the novice
have the same attentional spotlight. The master is shining it on
bigger objects.

\section{Why the models disagree}

The previous chapter argued that the field's incompatible maps of
memory are not a failure but a reflection of genuinely different
questions being asked of the same territory. Working memory is the
clearest possible illustration of that argument.

Baddeley's multicomponent model is a description of the \emph{storage
architecture} --- the structure of the system's components and the
evidence that they are separable. That is a productive question to
ask, and the model that emerged from it has organized five decades of
experiment.

Cowan's embedded-processes model is a description of the
\emph{dynamic relationship between working memory and long-term
memory} --- the question of whether the two are one system or two.
That is also a productive question, and the evidence that working
memory is simply activated long-term memory plus an attentional focus
is substantial.

Engle's attentional control framework is a description of
\emph{individual differences} --- of what varies between people and
what that variance predicts. That is a third productive question, with
a substantial body of evidence behind it too.

None of these is a wrong question. None of the frameworks is
empirically false. They are organizing different aspects of a thing
that is genuinely complex, and the disagreement between the maps ---
which the field's reference ontologies faithfully reproduce, nesting
working memory as a kind, then as a component, then as an ability,
in incompatible hierarchies --- is evidence of that complexity rather
than a failure to resolve it.

What they do collectively say, with more confidence than any single
model, is this: the mental desktop is real, it is limited, it is
different from long-term storage, and it is deeply entangled with
attention. The exact architecture of that entanglement is where the
science is still being done. The four items you can hold right now are
not a simple fact about a simple system. They are the visible surface
of one of the more genuinely open questions in cognitive science.







ch3 pain.tex
\atlaschapter{Pain}{Why pain is not a signal from the body but a verdict by the brain, why that verdict can be wrong in both directions, and what happens when the system gets stuck.}

In 1944, an American military surgeon named Henry Beecher was treating
wounded soldiers at the Allied landings in Italy and noticed something
that troubled him enough that he wrote it down. Many of the men arriving
from the front with severe injuries --- broken bones, deep lacerations,
gunshot wounds serious enough to require amputation --- were not in
distress. Asked about pain, a significant number said they felt little
or none, and declined morphine. They were conscious. They were not in
shock. They simply did not hurt, or did not hurt in proportion to what
had been done to them.

Beecher was a careful observer, and the finding nagged at him. After the
war he compared the same wounds in civilian surgical patients and found
the reverse: civilians with injuries of equivalent severity almost
universally wanted morphine and reported severe pain. The difference, he
concluded, had to be in the meaning of the injury. For the soldier, a
wound meant survival, evacuation, the war over for him. For the
civilian, the same wound was a catastrophe interrupting a normal life.
Same tissue damage, opposite pain. The nerve fibers were identical. The
brain processing their signals was not.

Beecher's observation is the entry point into one of the most
counterintuitive facts in cognitive science: pain and injury are not the
same thing, and the relationship between them is governed not by the
damage done to the body but by what the brain makes of that damage.
Pain is not a signal transmitted from the body to the brain, the way a
telephone transmits a voice. It is a decision made by the brain, using
signals from the body as evidence, about whether the situation is
threatening enough to hurt. And like any decision, it can be
correct, mistaken, delayed, revised, or stuck.

\section{The telegraph model and why it fails}

For most of history, pain was pictured as a telegraph. An injury fires a
signal along dedicated wires to a pain center in the brain, and the
brain rings its alarm. The louder the damage, the louder the alarm.
The model was proposed explicitly by Descartes in the seventeenth
century, who imagined threads running from the skin to the brain that,
when pulled by heat or injury, rang a bell. Different century, same
picture: a straight line from body to experience, damage as the cause
and pain as the effect, locked in a one-to-one relationship.

The facts have never cooperated. Soldiers at Anzio were one kind of
counterexample. Phantom limbs were another: people with amputated limbs
frequently experience pain in the absent one, sometimes severe and
chronic, in tissue that no longer exists. The telegraph has no wire to
send. The signal has no origin. And yet the pain is entirely real
--- as real as any other pain, as distressing, as debilitating, and
resistant to treatments aimed at the missing body part, because the body
part is not where the trouble is.

The telegraph model also cannot explain why the same injury hurts
differently depending on attention. An athlete who breaks a finger during
a game may not notice until the final whistle; a person who breaks the
same finger stepping off a curb notices immediately and cannot think
about anything else. The finger is the same. The damage is the same.
The difference is whether the brain, occupied with the urgencies of
competition, has flagged the injury as requiring immediate processing or
can afford to defer it. Pain, whatever it is, is interruptible in a way
that a simple sensory signal should not be.

\section{The gate}

The decisive break from the telegraph model came in 1965, in a paper
published in \emph{Science} by two researchers --- Ronald Melzack, a
psychologist at McGill, and Patrick Wall, a neuroscientist at MIT ---
that is now considered one of the foundational papers of modern
pain research. They called their account the gate control theory, and
the word \emph{gate} turned out to be an inspired choice: it evoked
an image clear enough that even people who did not follow the
physiology could grasp the main point.

The main point was this. In the dorsal horn of the spinal cord ---
the region where sensory signals from the body arrive before being
relayed upward to the brain --- there is a modulatory mechanism that
can open or close the channel through which pain signals pass. Two
types of nerve fibers bring information from the site of an injury: thin
fibers that carry pain signals, and thicker fibers that carry touch,
pressure, and vibration. Crucially, activity in the large fibers
\emph{inhibits} the transmission cells that relay pain upward.
When you rub a bruise, you are not distracting yourself from the pain;
you are literally closing the gate, flooding the spinal cord with
large-fiber touch signals that suppress the thin-fiber pain signals
competing for the same relay. The gate control theory gave that
homely experience a physiological mechanism for the first time.

What made the theory genuinely revolutionary was the second part.
Melzack and Wall proposed that the gate was not only modulated by
signals coming up from the body; it was also modulated by signals
coming \emph{down} from the brain. The brain could, in other words,
turn the gate up or down based on its own state --- based on attention,
emotion, expectation, memory of prior pain, and the meaning of the
current situation. This was not a concession to psychology as a soft
addendum to the hard physiology. It was a physiological claim: the
nervous system has a descending pathway that can amplify or dampen
incoming pain signals at the spinal gate before they ever reach
consciousness. The soldier's brain, having assessed the situation as
one of escape and survival, was not failing to transmit the wound;
it was actively closing the gate.

\aside{The gate control theory was not perfectly accurate in its
details --- the specific circuit Melzack and Wall proposed in 1965
required subsequent revision --- but its conceptual core has been
confirmed and extended by fifty years of research. What they got right
was the architecture: pain is modulated at multiple levels of the
nervous system, and the brain's top-down influence on incoming signals
is as real as the signals themselves.}

The practical consequences followed immediately. If pain can be gated
by large-fiber activity, then applying mild electrical stimulation to the
skin --- transcutaneous electrical nerve stimulation, or TENS --- should
recruit large fibers and close the gate, reducing pain without drugs.
It does, and TENS devices are now standard rehabilitation equipment.
If pain is modulated by attention and emotion, then psychological
interventions are not complementary to pain treatment; they are pain
treatment, acting on the same neural machinery as analgesic drugs,
through a different point of entry.

\section{Pain from a limb that isn't there}

Phantom limb pain is the clearest possible demonstration of the gate
control theory's central implication: that pain is constructed by the
brain, not transmitted from the body. It also introduced a piece of
neuroscience that, once understood, reframes the experience of pain
entirely.

When a limb is amputated, the region of the brain's sensory cortex that
used to process sensation from that limb does not go quiet. The brain
maintains a detailed spatial map of the body --- the somatosensory
cortex, where each body region occupies a strip of tissue proportional
roughly to that region's sensitivity --- and after an amputation, the
territory formerly assigned to the missing limb begins to be colonized
by neighboring regions. In the cortical map, the hand sits adjacent
to the face. After a hand amputation, the face's cortical territory
expands into the vacated hand region, so that touch on the cheek
activates neurons that used to represent the missing fingers.

The neuroscientist V.~S. Ramachandran demonstrated this elegantly in
the 1990s by touching different spots on the faces of arm amputees and
asking what they felt. Many reported feeling the touch simultaneously in
their phantom hand, with different spots on the face corresponding to
different phantom fingers --- a miniature ghost of the missing hand
mapped across the cheek. The cortical representation had literally
moved into new real estate, and sensation from the face's skin was
activating it. The phantom limb was not an illusion or a psychosomatic
complaint. It was the brain's body map running without the body part
it was mapped to.

From this Ramachandran reasoned that phantom pain might arise from a
mismatch in the map: the motor cortex was sending movement commands to
the phantom limb, but with no limb there to execute them and no sensory
feedback to confirm they had succeeded, the brain was receiving a
signal consistent with a limb frozen in an agonising contraction it
could not resolve. His intervention was almost absurdly simple: a
cardboard box with a vertical mirror down the centre. The patient
placed the intact limb on one side and the residual stump on the other,
and looked into the mirror from the intact side, so that the reflection
appeared to show two intact limbs. They then moved the real hand while
watching the mirror. The brain, receiving visual feedback of a hand
moving normally in the position of the phantom, sometimes experienced
the phantom limb as unclenching and relaxing --- and with the
unclenching came relief from the pain. A mirror was pharmacologically
active, because what it was acting on was not a limb but a brain's
representation of one.

Mirror therapy's evidence base remains uneven --- results vary widely
between patients, and the mechanisms are still debated --- but as a
demonstration of the principle it is hard to improve on. The therapy
works, when it works, by deceiving the brain's body map with visual
information, and the deception has the same neural currency as any
other pain intervention.

\section{The brain's pharmacy}

If the gate control theory established that the brain modulates pain
from above, the study of placebo analgesia established just how
concretely it does so.

Give a person a sugar pill and tell them it is a powerful painkiller.
Their pain often decreases measurably. This is not a finding about
gullibility or self-deception. In 1978 researchers showed that
the pain reduction produced by a placebo could be blocked by naloxone
--- a drug that occupies opioid receptors and prevents anything else
from binding to them. If naloxone blocked the placebo's effect, then
the placebo was working \emph{through opioid receptors}. The brain,
in response to the expectation of relief, was releasing its own
endogenous opioids --- the same class of molecules as morphine ---
into its own descending pain-suppression pathway. The placebo was not
fooling the patient into thinking the pain was less. It was triggering
the brain to chemically reduce it.

The brain, it turns out, manufactures its own analgesic pharmacy and
can dispense it in response to expectation and conditioning. This is
why placebo effects in pain trials are often large --- not because
patients are pretending or measuring incorrectly, but because
expectation genuinely closes the gate, via the same descending
pathways that Melzack and Wall identified. A patient who believes a
treatment will work is partly right regardless of whether the treatment
has any intrinsic pharmacological activity, because the belief itself
activates a real analgesic system.

The mirror of this is the \emph{nocebo} effect: tell a person that a
harmless procedure will cause pain, and it often does. Nocebo
hyperalgesia --- pain produced by the expectation of pain --- involves
measurably different neurochemistry from placebo analgesia, including
the activation of the stress hormone axis and the release of
cholecystokinin, a peptide that is, among other things, an endogenous
anti-opioid: it blocks the very system that placebo analgesia
activates. Belief that something will hurt can pharmacologically
disable the system that would otherwise dampen the pain. The mind does
not override the body in these cases; the mind \emph{is} part of the
mechanism.

\aside{The nocebo effect has consequences that go well beyond the
laboratory. Patients who are told a procedure will be very painful
and are given a full list of side effects experience more pain and
more side effects than patients given the same procedure and the
same information framed differently. The informed consent process,
in other words, is pharmacologically active. Negative expectation
is not inert.}

\section{When the system gets stuck}

Acute pain is pain working as designed. A tissue injury produces
signals, the gate opens, the brain constructs an experience of pain
proportional to its assessment of the threat, the person attends to
and protects the injured area, healing occurs, and the pain resolves.
This is the system doing its job.

Chronic pain is what happens when the system gets stuck --- and
understanding \emph{how} it gets stuck has required giving up the
telegraph model entirely.

In many chronic pain conditions, there is no detectable ongoing injury
or tissue damage that would account for the pain's persistence.
The wound has healed. The inflammation has resolved. Scans show nothing.
And yet the pain is as real and as severe as ever, sometimes worse.
The old model had no vocabulary for this beyond ``it's in your head,''
a phrase that has caused considerable harm to people with legitimate,
disabling conditions. The better account requires understanding what
happens to the nervous system when pain goes on too long.

The phenomenon is called \emph{central sensitization}. When pain
signals are sustained for long enough, the spinal cord and brain
undergo changes that amplify the system's gain. Neurons in the dorsal
horn become more excitable, responding to inputs they would previously
have ignored. The brain's descending inhibitory pathway --- the one
that closes the gate --- becomes less effective. Pain begins to arise
from stimuli that are not painful: a light touch, a mild temperature
change, a vibration. The threshold has dropped so far that the system
is now generating pain in response to ordinary life, and generating it
not because the body is damaged but because the pain-processing
machinery itself has been recalibrated.

This is not a metaphor. Central sensitization is measurable in
neurophysiology, and its effects on the brain are visible on scans.
People with chronic pain show reduced gray matter in regions of the
prefrontal cortex involved in pain modulation, and altered connectivity
between brain regions that normally communicate well. The brain of
someone in chronic pain is structurally different from the brain of
someone who is not. It has been changed by the pain it has been
processing.

One of the most reliable predictors of whether acute pain will become
chronic is a psychological variable called \emph{pain catastrophizing}:
the tendency to ruminate on pain, magnify its threat, and feel helpless
to control it. This is not a character flaw or a weakness of will. It
is a measurable cognitive style with its own neural correlates, and it
predicts chronicity better than the severity of the original injury.
The reason is not obscure given the mechanism: catastrophizing keeps
attention locked on pain signals, sustains the emotional arousal that
opens the gate, and impairs the top-down inhibitory pathways that would
otherwise help close it. The psychology and the physiology are not
parallel stories; they are the same story in different languages.

The implication for treatment has been slow to reach clinical practice.
If chronic pain is often a disorder of the pain-processing system rather
than a disorder of the body part that originally hurt, then treatments
aimed exclusively at the body --- more imaging, more surgery, more
injections --- are aimed at the wrong target. The treatments with the
strongest evidence for many chronic pain conditions are psychological:
cognitive behavioral therapy that addresses catastrophizing and
fear-avoidance, mindfulness approaches that change the relationship
between the patient and the sensation, graded exposure that
de-catastrophizes movement. These act on the central sensitization
directly. They are not alternatives to ``real'' pain treatment. They
are the closest thing to one.

\section{What the maps disagree about}

Pain sits at an unusual intersection in the landscape of cognitive
science, which is why the field's reference atlases handle it so
differently. A neurophysiologist mapping pain describes nociception:
the specialized sensory apparatus that detects tissue damage, the
thin fibers and their thresholds, the spinal relay. A clinical
psychologist mapping pain describes suffering and catastrophizing and
the conditions under which acute pain becomes chronic. A
phenomenologist mapping pain describes the experience itself, which
has properties --- its urgency, its demandingness, its tendency to
monopolize attention --- that set it apart from any other sensation.
These are not descriptions of the same thing at different levels of
resolution. They are partially overlapping accounts of different
aspects of a phenomenon that runs from molecular receptor to
conscious dread, and no single vocabulary covers the full span.

The four-rung ladder that careful taxonomists draw through this
territory --- nociception, pain, imagined or anticipated pain, the
suffering that surrounds it --- is not pedantry. Each rung marks a
real discontinuity. You can have nociception without pain: the soldier
at Anzio, the athlete who finishes the game on a fractured ankle. You
can have pain without nociception: the phantom limb, the centrally
sensitized patient whose body has healed. You can have the anticipation
and dread of pain without current tissue damage or even current pain,
and that anticipation can, via nocebo mechanisms, produce the very pain
it anticipates. The rungs are distinct.

What they share, which no rung of the ladder fully captures, is that
pain is the one experience the body has that is inescapably a cognitive
event from start to finish. Nociception is peripheral and relatively
mechanical. Pain is what the brain does with nociception, and that
doing involves expectation, memory, attention, meaning, and belief in
proportions that vary between people and moments and that no simple
model has yet captured whole. The disagreement in the maps is a
faithful record of how much of the terrain remains genuinely open.





ch4.tex
\atlaschapter{Attention}{Attention is not the same as perception. It is the selection of what gets perceived --- and what it selects, and how, has turned out to be genuinely difficult to explain.}

Here is a simple test. Read this sentence to its end, then stop and ask
yourself: what sounds were in the room while you read it? You almost
certainly cannot say. The sounds were there. Your auditory system
processed them, at least partially, at the level of the cochlea and
the early auditory cortex. But they did not reach awareness. Something was
selecting them out before they could cross into conscious experience ---
amplifying some signals, suppressing others, before the question of
awareness even arose. That selection process is attention.

Attention is the process by which the mind selects, from the continuous
flood of sensory information arriving every moment, what gets amplified
and processed in depth, and what gets passed over. The flood does not stop
when you attend to something; the world goes on producing light and
sound and movement at the same rate whether you are watching or not.
What changes is what gets amplified and what gets passed over. Attention
is, in this sense, less like a spotlight and more like an editor --- and,
as with any editor, what gets cut can matter as much as what makes it
through.

Previous chapters have touched on attention repeatedly without naming it.
The gate control theory of pain showed how attentional state opens or
closes the spinal gate: the athlete who doesn't notice a fracture until
the final whistle is not suppressing pain through willpower but through
attentional allocation. Working memory's central executive is, in large
part, an attentional control system. Memory's reconstructive quality
depends on what was attended at encoding: details outside the
attentional focus leave fainter traces, which is partly why eyewitness
memory is unreliable for peripheral events. Attention runs through the
whole of cognition. But it is not peripheral infrastructure. It is one
of the most actively contested topics in the field, with three
well-developed theoretical frameworks that agree on the data and
disagree, sharply, about what the data means.

\section{The cocktail party problem}

The formal study of attention began with a noise problem. In the years
after the Second World War, British air traffic controllers were
receiving voice messages from multiple aircraft simultaneously through
a single loudspeaker in the control tower. The voices overlapped and
tangled. Controllers were failing to separate them. Someone needed to
understand how people extracted a single signal from a mixture of
competing signals --- not as an engineering problem but as a cognitive
one.

In 1953, a Cambridge researcher named Colin Cherry gave the problem its
enduring name. Cherry called it the cocktail party problem: the puzzle
of how a person at a noisy party manages to follow one conversation
while ignoring several others. He studied it in the laboratory by
presenting two different spoken messages simultaneously, one to each
ear through headphones, and asking participants to \emph{shadow} one of
them --- repeat it aloud, word by word, in real time. Participants
managed this well. What they could not do was report much of anything
from the message they had been ignoring. They noticed the speaker's
voice changed gender. They noticed whether the unattended ear received
speech or pure tone. They did not notice when the unattended message
switched to a foreign language, or when a single word was repeated
dozens of times, or when the speech was played backwards. The meaning
of the ignored message was apparently not reaching them.

Almost. Because Cherry also noted something that didn't fit his findings
cleanly: people sometimes heard their own name spoken in the unattended
ear, even when they missed everything else in it. A message you have
tuned out can apparently still search for you in it.

This wrinkle mattered. A system that filtered the unattended channel so
completely that its language went unregistered should not know when that
channel contained your name --- because recognising your name requires
processing its meaning. Somehow, without conscious awareness, the
unattended stream was being processed deeply enough to extract identity
information. The full implications of that fact drove decades of
argument.

\section{The gorilla in the room}

Cherry's finding was about what we \emph{hear} when distracted. In
1999, two psychologists at Harvard, Daniel Simons and Christopher
Chabris, demonstrated something comparable about vision, and did it in
a way that has become one of the most replicated and referenced
demonstrations in the history of cognitive science.

They made a short video. In it, two teams of three players each ---
one team in white shirts, one in black --- moved around each other
passing basketballs. Participants were asked to watch the video and
count the number of passes made by the white team. It required
concentration: the players were moving, the two teams were interleaved,
and distractors were present. About halfway through the video, a person
wearing a full gorilla costume walked into the centre of the frame,
faced the camera, thumped their chest, and walked out. The gorilla was
on screen for nine seconds.

Roughly half of participants, asked afterwards whether they had noticed
anything unusual, said they had not. When told to watch the video again
--- same video, same gorilla --- they were frequently astonished. The
gorilla was not subtle. It was not a visual trick or an ambiguous
figure. It walked through the centre of the scene in plain view. The
participants had simply not seen it. Not because their visual system
had failed to register light from that region of the screen, but
because their attention was occupied counting white-shirt passes, and
the gorilla, wearing black, fell outside the attentional scope they
were maintaining.

Simons and Chabris called this \emph{inattentional blindness}: the
failure to consciously perceive a fully visible, fully lit stimulus
because attention was directed elsewhere. It is related to, but
distinct from, \emph{change blindness}, a phenomenon documented around
the same period: the remarkable tendency of people to miss substantial
changes to a scene --- a character's jacket changing colour between
film cuts, a different person completing a conversation when the
original was briefly hidden by a passing van --- when the change does
not trigger an attentional alert.

\aside{The gorilla experiment has since been replicated with thousands
of participants in multiple countries, under varied conditions, and in
virtual reality. Detection rates vary with the difficulty of the
counting task, the colour of the gorilla relative to the attended
team, and the gorilla's speed of movement. The core finding holds:
when attention is fully engaged on one task, clearly visible objects
performing unexpected actions can go entirely unnoticed. Radiologists
scanning lung images for nodules routinely miss a gorilla image
embedded in the scan at a rate of 83 per cent, in a study that was
deliberately designed to echo the original.}

Together, Cherry's cocktail party and Simons and Chabris's gorilla
establish something that most people find genuinely uncomfortable: the
rich, detailed perceptual world we seem to inhabit is largely a
construction. We consciously experience the small fraction of the
environment that attention has selected and amplified. The rest is
not experienced at all, or is experienced only as a vague, low-
resolution surround that carries little information. The sense that we
see everything in front of us is an illusion assembled from a narrow
attentional beam and a great deal of inferential filling-in.

\section{Three theories}

Cherry's name-in-the-unattended-ear problem triggered thirty years of
theoretical argument that divided into two camps and eventually
produced a third. The camps are worth understanding not because the
argument is settled --- it is not --- but because each position captures
something real, and the three together map the terrain more honestly
than any one alone.

The first account, proposed by the British psychologist Donald
Broadbent in 1958, held that attention is a \emph{filter} that
operates early, before the brain has processed the meaning of incoming
signals. Information is selected on the basis of physical
characteristics --- the location it comes from, the pitch of the voice
carrying it, its loudness --- and only selected information receives
full processing. The rest is blocked at the gate. Broadbent's model
was explicit and testable, and it accounted for Cherry's main
finding well. People didn't know the language of the unattended message
because the filter blocked it before linguistic processing occurred.

The problem was your name. The name-recognition problem pushed the
philosopher and psychologist Anne Treisman, working in the 1960s, to
propose an amended account. The filter, she suggested, was not a gate
that blocked signals entirely but a volume control that
\emph{attenuated} them --- turned them down rather than off. Unattended
signals were still processed, but at reduced strength. Most stimuli,
weakened below the threshold of awareness, failed to register
consciously. But certain stimuli had very low thresholds: your own
name, a cry of pain, any sudden loud sound. These could punch through
the attenuated signal and reach awareness even from the
unattended channel. The gorilla might fit here too: not blocked but
turned so far down that only the most salient objects break through.

Both filter theory and attenuation theory addressed the same question:
at what stage is selection made, and on what basis? The third approach,
developed by the psychologist Daniel Kahneman in 1973, asked a
different question entirely. Rather than asking where the gate is,
Kahneman asked: what is the budget? His proposal treated attention not
as a filter at all but as a limited \emph{resource} --- a pool of
cognitive fuel that can be allocated between tasks. Some tasks are
cheap, consuming little of the pool: breathing, walking on a
familiar route, reading words in your native language after years of
practice. Others are expensive: parsing a sentence in a second
language, tracking two sets of moving players while counting passes.
When the total demand on the pool exceeds its capacity, performance
suffers on all competing tasks proportionally.

On this account, the gorilla is missed not because a filter blocked
it but because the counting task consumed so much of the resource pool
that nothing was left to process unexpected events in peripheral
attention. The gorilla was expensive to process; the budget was spent;
the gorilla was effectively unaffordable.

\aside{The three theories answer three questions that sound the same
but are not. Broadbent asks: which signals get in? Treisman asks: how
far does an unattended signal get before it is turned down? Kahneman
asks: how much cognitive effort is available, and how is it spent?
Each question is real. Each theory generates predictions, experiments,
and findings. The data does not pick a winner because the data was
generated by different experiments designed to probe different aspects
of a phenomenon that is genuinely multi-faceted.}

\section{What the neurons do}

When brain imaging became possible, attention researchers hoped for
a verdict. Instead, they found evidence that complicated all three
accounts in illuminating ways.

The clearest finding came from a 1985 experiment by Jeffrey Moran and
Robert Desimone, recording from individual neurons in a visual area
of the macaque cortex called V4. They placed two stimuli simultaneously
inside a single neuron's receptive field --- the region of the visual
world to which that neuron responds --- and had the animal attend to
one of them. When the animal attended to a stimulus that drove the
neuron well (a red bar, say, for a neuron that preferred red bars),
the neuron fired strongly. When the animal attended to the other
stimulus instead, the neuron's response was dramatically suppressed,
even though exactly the same red bar was sitting in the same
receptive field. Attention was not merely brightening what it selected.
It was also darkening what it did not.

This result, endlessly replicated and extended over the decades since,
supports a picture of attention as a bias in a competition. At any
moment, the visual cortex contains many neurons responding to many
stimuli simultaneously, and these representations compete with each
other for dominance. Attention biases the competition: it tilts the
scales toward the attended stimulus and against its rivals. The winners
of the competition propagate their signals forward through the
visual hierarchy. The losers are suppressed, often before they reach
the regions associated with awareness.

This neural picture does not map cleanly onto any of the three
behavioural theories. It is not early selection in Broadbent's sense,
because the competition takes place within processing areas, not at
a gate before them. It is not simply Treisman's attenuation, because
the suppression is active and competitive rather than passive. And it
is not Kahneman's resource pool, though the competition-for-
representation mechanism could plausibly explain why capacity is
limited: there is only so much cortical territory, and attended
representations consume more of it. The neural data adds a fourth
vocabulary to describe what attention is doing, and the four
vocabularies do not translate into one another without friction.

\section{The illusion of the full picture}

Running through all of these accounts is an implication that the
research has made increasingly difficult to avoid: we believe ourselves
to see and hear the world, but what we actually experience is a small,
sharply processed sample of it, surrounded by a vague, expectation-
filled fog. The sense of inhabiting a rich, continuous perceptual
world is not a report of what the visual system is doing. It is a
story the mind tells itself, assembled from the attended region,
smoothed out by prior knowledge and expectation, and mistaken for
comprehensive coverage.

Change blindness experiments make this stark. In one classic sequence
of studies, people failed to notice when the person they were giving
directions to was replaced, mid-conversation, by a different person
--- the swap concealed behind a door momentarily passed between them.
Same height, different face, different jacket. Many participants
continued the conversation without apparent awareness that anything
had changed. In film editing studies, changes of camera angle, lighting,
and even major props between cuts went undetected so reliably that
film editors developed an intuitive art of hiding such changes by
timing them to moments of high action, when attention was occupied.

These effects are not about low intelligence or poor observation.
They reflect the basic architecture of the system. The visual system
did not evolve to record the world. It evolved to support action in
the world. For that purpose, what is currently relevant to the task
at hand matters enormously. What lies outside the task's attentional
scope can safely be backgrounded --- and for most of evolutionary
history, the things backgrounded were genuinely irrelevant. The system
is not broken when it misses the gorilla. It is doing exactly what it
was built to do: reserving its processing capacity for the task in
front of it. The mismatch is between what the system was built for
and what we now expect it to do in a complex, unpredictable modern
world where the gorilla is sometimes a car entering an intersection
or an alarm in a hospital ward.

\section{Why the maps disagree}

The field's taxonomies of attention reproduce its theoretical fault
lines without fully resolving them. Some catalogs treat attention as
a single faculty with subtypes --- sustained, selective, divided,
executive --- arranged in a hierarchy. Others treat it as a family
of loosely related processes that happen to share a name. Others wire
it in as a resource that other faculties draw on, like working memory's
central executive, with no stand-alone heading at all.

These are not transcription errors. They reflect genuine uncertainty
about whether attention is a \emph{thing} or a \emph{label} for
whatever happens when a limited-capacity system has to choose what to
process. Broadbent thought it was a mechanism --- a filter with a
definite physical location in the processing stream. Kahneman treated
it more as an economic fact about cognitive systems: limited capacity
requires allocation, and allocation is attention. The neuroscience
suggests it may be better described as a property of competition
between representations, not a separate mechanism at all.

What the three accounts agree on is more important than what they
dispute. They agree that the mind cannot process everything at once.
They agree that the selection of what gets processed is active,
consequential, and often wrong. And they agree that the feeling of
awareness --- the sense of looking directly at something, of hearing a
voice clearly, of having something fully in focus --- is not a reliable
guide to what attentional processes are actually selecting or
suppressing. The gorilla was in plain sight. Attention was elsewhere.
That is not an accident or a flaw in one experiment. It is, as far as
the evidence goes, how the system works.

What none of the three theories fully addresses is why any of this
produces the feeling of awareness at all --- why an amplified signal
becomes an experience rather than simply a bias in behaviour, why
attending to something feels like anything from the inside. That
question, which attention research keeps bumping into without being
able to answer, is where the next chapter begins.






ch5.tex

\atlaschapter{Consciousness}{The previous chapter explained what attention does. This one asks why it feels like anything.}

Right now, something is happening that is easy to overlook precisely
because it is so close. You are reading these words, and there is
something it is \emph{like} to read them. The black shapes on the
page are not merely being detected and decoded --- there is an
experience accompanying that decoding, an inner illumination, a
sense of meaning arriving. That experience is not automatic, in the
sense that a great deal of neural processing happens without any
accompanying inner light at all. Your visual system is doing
continuous work on the peripheral regions of the page --- adjusting
for contrast, tracking the edges of letters --- without any of it
rising to awareness. Somewhere in the processing chain, a threshold
is crossed, and the words you are attending to become not merely
processed but \emph{experienced}.

What is that threshold? What crosses it? And why does crossing it
produce experience at all, rather than simply producing more and
faster computation, happening in the dark?

These questions are not new. They have occupied philosophers for
centuries and neuroscientists for decades. What is new is that we now
have two serious scientific theories that try to answer them, a body
of empirical evidence that neither theory fully accounts for, and a
sharper understanding than we have ever had of exactly where the
difficulty lies. The chapter that follows is not a report of a solved
problem. It is a map of the most interesting unsolved one in
cognitive science.

\section{The easy problems and the hard one}

In 1995, the Australian philosopher David Chalmers drew a distinction
that has organized the debate ever since, between what he called the
easy problems of consciousness and the hard one. The easy problems are
not actually easy --- they include explaining how the brain integrates
information from different sensory systems, how it discriminates
between stimuli, how it controls attention, how it reports its own
internal states. These are the problems that neuroscience and cognitive
psychology are progressively solving, with increasing success. They
are called easy not because the answers are obvious but because we can
at least see what kind of answer would work: a mechanism that performs
the function. Explain the circuit, trace the pathway, map the
computation, and the problem is solved.

The hard problem is different in kind. Even granting a complete
functional explanation of how the brain processes information, selects
what to attend to, and reports its internal states, there is a further
question that the functional explanation does not touch: why is any of
this accompanied by experience? Why is there something it is \emph{like}
to see red, to feel pain, to follow a sentence? Why isn't the whole
machinery running in the dark --- processing, selecting, reporting, all
without an inner light?

Chalmers illustrated the gap with a thought experiment. Imagine a
being --- a philosophical zombie --- that is physically and functionally
identical to you in every measurable way: same neurons, same firing
patterns, same behaviour, same verbal reports about having rich inner
experiences. The zombie says ``I am in pain'' when injured and ``I see
red'' when looking at a tomato, because those are the outputs its
functional state produces. But there is, by stipulation, nobody home ---
no experience accompanying the outputs. The zombie's brain processes
information without anything it is like to be that brain.

The question Chalmers posed is whether such a being is even conceivable.
He argued yes, it is perfectly conceivable, and that this conceivability
reveals an explanatory gap: consciousness is not logically entailed by
the physical and functional facts, which means that a complete physical
and functional explanation of the brain would leave something out.
What it would leave out is experience itself.

\aside{Chalmers did not claim that consciousness is non-physical ---
only that the functional story is incomplete. Many neuroscientists
dispute this, arguing that the gap is an artefact of how we think
about the problem rather than a genuine feature of nature. The debate
is live and substantive. But the gap, whatever its ultimate status,
is a useful diagnostic: any theory of consciousness that cannot
explain why its proposed mechanism should produce experience, rather
than merely produce function, has not yet reached the hard problem.}

\section{The global broadcast}

The most empirically productive scientific theory of consciousness in
the past thirty years is the global neuronal workspace theory, developed
by the French cognitive neuroscientist Stanislas Dehaene and his
colleagues, building on an earlier framework proposed by the psychologist
Bernard Baars.

The starting observation is that the brain is not a single unified
processing system but a collection of highly specialised, largely
independent modules operating in parallel. Face recognition, motion
detection, phonological decoding, pain signalling, colour processing ---
each of these is handled by dedicated neural machinery that runs its
computations largely in isolation from the others and largely without
conscious access. The brain is doing an enormous amount at every moment.
Almost none of it reaches awareness.

Dehaene's proposal is that consciousness is what happens when information
escapes its local module and is \emph{broadcast} widely across a
long-range network of neurons with axons connecting the prefrontal and
parietal cortex to the rest of the brain. This network --- the global
workspace --- acts like a broadcasting system, taking information that
has been processed locally and amplifying it into a signal that becomes
available to many systems simultaneously: to memory, to language, to
reasoning, to voluntary action. Before broadcast, information is local
and functionally constrained; a stimulus can influence your motor
response without your being able to report it, describe it, or reason
about it. After broadcast, information is globally available. That
global availability is what Dehaene means by consciousness.

The broadcast is not a gradual process. It is, in his term, an
\emph{ignition}: a sudden, nonlinear, all-or-nothing explosion of
coordinated activity across the workspace network, triggered when
an incoming signal is strong enough to overcome the threshold. The
same stimulus, presented at slightly different strengths or under
slightly different attentional conditions, either ignites or it
doesn't --- and the difference in the resulting brain state is
dramatic. Below threshold, you see a faint, brief flicker of local
activity that dies out quickly. Above threshold, the entire workspace
lights up almost simultaneously in a cascade that persists for hundreds
of milliseconds. The difference between a stimulus you see and one you
don't is not a matter of degree. It is a phase transition.

The empirical support for this picture is substantial. Subliminal
priming experiments show that stimuli presented below the threshold
of awareness can still influence behaviour --- speeding responses to
related words, activating semantic associations --- without triggering
the late, widespread cortical activity that accompanies aware perception.
In the attentional blink paradigm, two targets are shown in rapid
succession; if the second appears within about half a second of the
first, it is often missed entirely, even though it was physically
presented --- the workspace, still broadcasting the first target, cannot
ignite again for the second. Neural recordings in such experiments show
the missed target producing early local activity in sensory areas but
failing to propagate into the broader network.

Return, with this in mind, to the gorilla from the previous chapter.
The counting task was keeping the global workspace fully occupied.
Gorilla-related signals were locally processed in visual areas ---
the patterns of light were detected, edges registered, motion tracked ---
but the workspace was not available to receive them. No ignition
occurred. Without ignition, no broadcast. Without broadcast, no
awareness. The gorilla was processed and not experienced. The
participants were not inattentive in some global sense; they were
paying fierce attention to the white shirts. They were unconscious
of the gorilla in the specific, mechanistic sense that the gorilla's
representation never achieved global availability.

\section{The model of attention}

Dehaene's theory is, at its core, an answer to the question of which
information becomes conscious and why. What it is less focused on is the
further question: why does broadcast feel like anything? Why does global
availability produce an inner light rather than simply producing flexible
behaviour? This is where the second major theory, proposed by the
Princeton neuroscientist Michael Graziano, offers something different.

Graziano begins with an analogy. The brain maintains a \emph{body
schema}: an internal model of the body's size, shape, and position in
space, which it uses to plan and control movement. The body schema is
not a photograph; it is a functional simplification, good enough for
the purposes of motor control, sometimes wrong in revealing ways.
Phantom limb pain, as the previous chapter described, arises partly
because the body schema continues to represent a limb that no longer
exists. The model and the body diverge; the model's distress is real
even though its object is gone. In body dysmorphic disorder, the model
misrepresents a body that is present. The point is that the brain
builds models of physical things --- simplified, useful, imperfect
--- in order to control them, and the models can persist and behave
independently of what they model.

Graziano's proposal is that the brain builds an analogous model of
attention itself --- an \emph{attention schema}. Attention is a
neural selection process: a competition among signals for the brain's
limited processing resources, which the previous chapter described
in terms of Moran and Desimone's biased-competition mechanism. That
process is complex, distributed, and largely opaque to introspection.
The attention schema is a simplified internal description of it, built
by the brain so that the control system for attention can regulate
itself without having to track every detail of what the underlying
competition is doing.

Awareness, on Graziano's account, is the content of that model. When
the brain's attention schema registers ``attention is currently applied
to X,'' the system concludes that it is \emph{aware of} X. This is not
a report of the underlying neural competition; it is a report of the
simplified model. And because the model is a simplification, it is
sometimes wrong. It describes attention as a unified, present-tense
experience rather than as the complex, distributed, temporally
smeared process it actually is. It makes consciousness feel immediate,
total, and transparent when the underlying mechanism is none of those
things.

\aside{The body schema analogy does real work here. Consider that
when a stroke damages the brain's body schema for the left side,
patients sometimes deny that the paralysed left arm belongs to them
at all --- it is present, they can see it, but the model says it is
not theirs. If damage to the attention schema produced an analogous
failure, one would expect deficits in the \emph{claim} to have
awareness --- the patient would process information without
generating the report of experience. Some neurological conditions
may involve exactly this. The prediction is testable, which is one
of Graziano's theory's strengths.}

The gorilla case looks different, and richer, from this angle.
Dehaene explains why the gorilla was not broadcast. Graziano explains
why the participants were \emph{confident} they had seen everything.
Their attention schema was tracking the white-shirted players ---
that was what the attention-control system needed to model. The schema
said ``attention is fully deployed on the scene.'' From the inside,
that felt like comprehensive awareness. The model was wrong: attention
was not comprehensively deployed on the scene in any sense that
included the gorilla. But models, as the phantom limb shows, do not
update automatically when the thing they model changes. The
participants' confidence in their own awareness was a report of the
model, not of the underlying neural reality.

\section{Where they converge and where they part}

Dehaene and Graziano are not rivals in the way rival theories usually
are. They are addressing adjacent questions, and the two accounts may
be parts of a single, more complete picture.

What GNW explains well is the \emph{mechanism of access}: which
information becomes globally available, why the transition is nonlinear,
and why subliminal stimuli can influence behaviour without being
reportable. The empirical evidence for the role of the prefrontal--
parietal network in conscious access is substantial and growing. If you
want to know why some processing reaches awareness and some does not,
GNW is the most concrete and testable answer currently available.

What AST explains well is the \emph{phenomenology of awareness}: why
conscious experience has the character it does, why it feels like a
unified, present-tense illumination rather than a competition among
representations, and why the brain's reports about its own awareness
are systematically oversimplified. The claim that the brain builds
a model of its own attentional state, and that the model's content is
what we call awareness, gives the sense of inner experience a
functional role --- the model is there because a system that can
model its own attention can control that attention better. Awareness
is not an epiphenomenon on this account; it is part of the engineering.

A natural synthesis is available: the global broadcast might be what
triggers the updating of the attention schema. When information ignites
the workspace, the attention-control system registers the ignition and
updates its model: ``attention is now applied to X.'' The model's
content is the experience of being aware of X. On this reading, GNW
tells you when the light comes on; AST tells you why there is a light
at all.

Neither theory, however, fully answers Chalmers's hard problem. Dehaene
is explicit that he considers the hard problem either dissolved or
misposed --- that explaining the functional and neural correlates of
consciousness is all that science can or needs to do, and that asking
why the mechanism produces experience rather than dark function is a
confusion. Graziano's position is similar: awareness \emph{is} the
model's content; there is nothing further to explain. But critics of
both theories point out that the question does not simply vanish
when you have a good functional account. Why should the attention
schema's content be \emph{experienced}? Why shouldn't the model update
silently, influencing behaviour without producing an inner light? The
gap between the mechanism and the experience remains, at least
apparently, and closing it may require ideas that neither theory yet
provides.

\section{Blindsight and the proof of dissociation}

The sharpest empirical evidence that attention and awareness are
genuinely distinct comes from a neurological condition called
blindsight, first documented systematically by the British neuroscientist
Lawrence Weiskrantz in the 1970s.

The primary visual cortex --- the first cortical area to receive signals
from the eyes --- is necessary for normal visual awareness. Patients with
damage to it lose conscious vision in the corresponding part of the
visual field: they report seeing nothing there. But Weiskrantz's patient,
known as DB, could not only detect the presence of a stimulus in his
blind field when pressed to guess; he could tell you something about
its orientation, its motion, its rough location --- all at rates well
above chance, while sincerely insisting that he was merely guessing and
had seen nothing whatsoever. He was astonished when shown his own
performance data. His visual system was processing the stimuli. His
awareness of that processing was absent.

A subsequent experiment by Robert Kentridge and colleagues went further.
They showed that spatial attention could be directed to locations in
DB's blind field, improving his detection performance in that region
--- a reliable signature of attentional allocation --- even when neither
the cue directing attention nor the target appeared to produce any
awareness. The attentional process was operating. The awareness was not.

This is the cleanest possible demonstration that attention is a neural
selection process and awareness is something additional, something that
the selection process usually but not necessarily produces. The
mechanism and the experience are dissociable. Damage the right
circuitry and you get the mechanism without the experience --- vision
without sight, detection without awareness, response without report.

Blindsight fits both theories. For Dehaene, it makes sense because
damage to V1 disrupts the feedforward signal that would otherwise
trigger ignition in the global workspace; the local processing in
surviving visual areas is insufficient to reach the threshold, so no
broadcast occurs, and no awareness follows. For Graziano, it makes
sense because the attention schema is presumably built on the output of
awareness-generating circuits; without that output, the schema has
nothing to model, and no claim of awareness arises. The theories make
the same broad prediction --- no awareness without the relevant
circuitry --- and the prediction is confirmed.

\section{Why the maps disagree}

Consciousness sits at an unusual position in the field's catalogs.
Some taxonomies treat it as a property that other cognitive processes
can have more or less of --- perception can be conscious or
unconscious, memory encoding can be explicit or implicit, attention
can operate with or without accompanying awareness. On this view,
consciousness is not a faculty alongside memory and attention but a
modifier that applies to them. Other frameworks treat it as a level
of processing that information either reaches or doesn't. Others
decline to categorise it at all, treating the question of what
consciousness is as genuinely open and inappropriate to
classify prematurely.

These are not clerical differences. They reflect a substantive
disagreement about whether consciousness is a \emph{thing} --- a
distinct cognitive system or state, the kind of entity that merits
its own heading in an ontology --- or a \emph{property} of other
things, or a \emph{threshold} that processing either crosses or
doesn't. Dehaene's theory implies something close to the threshold
view: consciousness is what happens when processing crosses the
ignition threshold into the global workspace. Graziano's theory
implies something closer to the property view: awareness is a
quality that the attention schema attributes to a state, and that
attribution can in principle be present or absent somewhat
independently of the processing. Chalmers's hard problem implies
that whatever we say about the mechanisms, the experiential dimension
may resist full capture by any functional categorisation.

What the field agrees on, which is itself more than was known fifty
years ago, is that consciousness is not a simple on/off switch
controlled by arousal. You can be awake and unaware of something
right in front of you, as the gorilla experiment shows. You can be
asleep and aware --- as in lucid dreaming. You can have processing
without awareness, as blindsight demonstrates, and perhaps awareness
without full attentional deployment, as some peripheral perception
studies suggest. The word covers a family of related but separable
phenomena, and any taxonomy that treats them as one thing is trading
accuracy for tidiness.

The most important thing this chapter can leave with the reader is not
a theory but a distinction. When you read these words and find them
meaningful, two things are happening: a neural selection process is
amplifying the signals from the page, and you are aware of the result.
The first is attention. The second is consciousness. They usually come
together, and they can come apart, and the question of why the first
produces the second --- why the selected signals feel like anything at
all --- is, at the moment of writing, genuinely open. The gorilla
was not seen. Its signals were processed. Those two facts, together,
are one of the most important things modern neuroscience has learned
about the mind.






ch6-perception.tex

\atlaschapter{Perception}{The eye is not a camera and the brain is not a screen. What we call seeing is a continuous act of inference, guided by expectation and corrected by surprise.}

You are looking at a flat page. The light reflected from it enters your
eyes and lands on the retina as a two-dimensional pattern. That
pattern is genuinely flat: it contains no depth, no three-dimensional
structure, nothing that unambiguously indicates the size or distance
of the objects in the room around you. It is blurry at the edges,
interrupted by a blind spot where the optic nerve exits --- a gap in
the image that covers a region of the visual field large enough to
swallow your thumb held at arm's length --- and sharply resolved only
in a small central region a few degrees across. The signal that reaches
the visual cortex is impoverished in ways that most people, accustomed
to the vividness of their own experience, would find astonishing.

And yet you see a room. Objects at different distances. A page in your
hand. You never see the blind spot, even though it is there. You
experience a continuous, richly detailed, three-dimensional world.

The gap between the signal and the experience is not filled by better
optics. It is filled by the brain. What you see is not a faithful copy
of the world arriving through your eyes. It is the brain's best current
hypothesis about what world would have produced the pattern of light
currently hitting the retina. Perception is inference.

\section{Helmholtz's insight}

The German physicist and physiologist Hermann von Helmholtz stated this
clearly in the nineteenth century, which makes it all the more
remarkable that it took so long to become the organizing principle of
perceptual science. The visual system, he argued, performs what he
called \emph{unconscious inference}: it takes ambiguous, impoverished
sensory data and infers --- rapidly, automatically, below the threshold
of awareness --- the most probable cause of that data. What reaches
consciousness is not the data but the inference.

The word ``inference'' is deliberate. The relationship between the
retinal image and the objects in the world is not unique: any given
pattern of light on the retina is consistent with infinitely many
possible arrangements of objects, distances, and illumination
conditions. A small object close to your eye produces the same retinal
image as a large object far away. A brightly lit dark surface and a
dimly lit light surface can produce the same signal at the retina.
The brain's task is to determine the most probable cause of the signal
--- and to do so fast enough to be useful for action.

It solves this by using \emph{prior knowledge}: accumulated experience,
encoded in the structure of the visual system over both evolution and
individual development, about how the world tends to be. Faces are
almost always convex. Illumination tends to come from above. Objects
at the same apparent size but farther away are larger. These
regularities are built into the inference machinery as expectations,
and those expectations are what let the brain resolve the ambiguity
that the raw signal cannot.

The inference is usually so fast and so accurate that it seems like
direct perception. You do not feel yourself inferring the room. You
simply see it. The invisible machinery of inference is part of what the
previous chapter called the processing that happens without awareness.
The experience of seeing is the inference's output, not the inference
itself.

\section{The brain as hypothesis-tester}

A more precise and more testable version of Helmholtz's insight
emerged a century later, in a computational framework proposed by
Rajesh Rao and Dana Ballard in 1999 and now known as
\emph{predictive coding}. It has since been developed into a
broad theory of cortical function by the neuroscientist Karl Friston.

The core proposal is this. The visual cortex is not organized as a
one-way processing pipeline, taking raw signals from the retina and
progressively extracting edges, shapes, objects, and scenes as the
signal passes upward through a hierarchy of areas. Instead, it runs
in two directions simultaneously. The higher levels of the hierarchy
--- areas encoding objects, categories, contexts --- send
\emph{predictions} downward to lower levels, specifying what they
expect the lower areas to be detecting. The lower levels compare
incoming signals to those predictions, compute the discrepancy ---
the \emph{prediction error} --- and send only that error signal
upward for further processing.

The architecture has a striking consequence. What propagates up the
visual hierarchy is not the full sensory signal but only what was
not predicted: the surprise, the deviation from expectation. Everything
that fits the current hypothesis is silently confirmed; only what
violates it demands the brain's limited processing resources. The
brain does not so much receive the world as continuously guess at it,
using the incoming sensory data not as a picture to be read but as a
set of corrections to apply to the guess.

This is more than a metaphor. The anatomy of the visual cortex
matches the prediction. There are two sets of connections between each
level of the hierarchy: backward connections, which are numerous and
send signals downward toward the retina, and forward connections,
which send signals upward toward higher areas. On the predictive
coding account, the backward connections carry predictions; the
forward connections carry errors. The predictions flow down; the
corrections flow up. Perception is the running negotiation between
them.

\aside{When you enter a familiar room, most of what you see is
predicted so accurately that almost no error signal is generated. The
brain barely needs to update its model. When something is wrong ---
a piece of furniture moved, a stranger sitting where a friend
usually sits --- the prediction error is large and it reaches
awareness immediately. Surprise is not merely an emotion. It is
the mechanism by which new information enters perception.}

The framework connects the Consciousness chapter's account of
attention and ignition to the mechanics of seeing. The gorilla was
not predicted; but the workspace, occupied with the counting task,
was not available to amplify and propagate the error signal its
appearance generated. The surprise of the gorilla existed, briefly,
as a prediction error in visual areas --- and was then not amplified
into the global broadcast that would have made it conscious.

\section{When the prior wins}

The value of predictive coding as a theory lies partly in the
experiments that it enables. If perception is inference guided by
prior expectations, then illusions --- cases where perception is
systematically wrong --- are not failures of the visual system but
diagnostic readouts: they show which priors are operating and
how strongly they resist correction by sensory data.

The most instructive example is the hollow mask illusion. Take a
mask of a human face --- the rotating Charlie Chaplin mask is the
classic version --- and film it turning slowly through 360 degrees.
When the outside of the mask is facing you, you see what you expect:
a convex face. When the inside rotates toward you, something strange
happens. The hollow concave surface appears to pop outward and present
itself as convex, a face gazing back at you. The mask appears to be
looking at you from the front even as it visibly rotates away. The
illusion is powerful, persistent, and entirely resistant to the
knowledge that it is an illusion. You can stare at it knowing exactly
what it is and what is happening, and you cannot see it correctly.

What is doing this? The prior. Faces are almost invariably convex in
human experience. The visual system has encoded this as a near-absolute
expectation: if a face-shaped object is present, it is convex.
When the binocular depth information signals concavity --- as it
correctly does, because we have two eyes for exactly this purpose ---
the face-prior is so strong that it overrides the sensory evidence.
The prior wins. What you see is the brain's expectation, not the
world's actual geometry.

There is a remarkable clinical observation attached to this. People
with schizophrenia are, as a group, largely immune to the hollow mask
illusion. They see the mask correctly as concave when its inside faces
them, where neurologically typical observers cannot. This sounds like
a perceptual advantage, and in a narrow technical sense it is.
The explanation, on the predictive coding account, is that schizophrenia
involves a weakening of top-down priors relative to bottom-up sensory
signals --- the brain's predictions are less able to override incoming
evidence. In the case of the hollow mask, weaker priors means the
binocular depth signal wins, and the mask is seen veridically. In
most perceptual situations, however, strong priors are what allow
perception to work at all: the blind spot is filled, ambiguities are
resolved, the world is made coherent faster than sensory data alone
could justify. The ability that lets schizophrenia patients see the
hollow mask correctly is the same weakened prior system that in
other contexts contributes to the disorder's characteristic
perceptual and inferential anomalies.

A more quotidian example of priors at work was the 2015 internet
dispute over a photograph of a dress. Some people saw it as white
and gold; others saw exactly the same photograph as blue and black.
Neither group was wrong about what they saw, and neither group was
seeing it incorrectly given their priors. The relevant prior is about
the colour of the illumination: if the brain assumes the dress is
lit by warm light (sunlight or tungsten), it subtracts that warmth
from the surface colour and infers white and gold. If it assumes
the dress is in shadow and lit by the cooler ambient sky, it infers
blue and black. The photograph was ambiguous enough that people's
different assumptions about illumination led to genuinely different
percepts of the same image. The dress case is trivial in itself
and revealing in what it proves: that the percept depends on the
prior, that different priors produce different percepts, and that
neither group is hallucinating.

\section{Two visual systems}

So far the chapter has treated vision as a single system inferring
a single world. The story is more complicated. Vision turns out to
contain two largely parallel processing streams with different
purposes, different neural substrates, and different relationships
to consciousness.

The distinction was initially framed in 1982 by Leslie Ungerleider
and Mortimer Mishkin as the difference between ``what'' and ``where'':
one stream running from primary visual cortex downward into the
temporal lobe, specialised for identifying objects, and another
running upward into the parietal lobe, specialised for locating them
in space. A decade later, the neuropsychologists Melvyn Goodale and
David Milner refined the distinction in a way that has proved more
fruitful: the two streams do not divide identification from location
but rather \emph{perception} from \emph{action}. The ventral stream
builds representations of objects for recognition, categorisation,
and conscious report. The dorsal stream guides the real-time control
of movements directed at objects, using a continuously updated
egocentric representation that is largely unavailable to awareness.

The evidence for this reformulation came principally from a patient
known as DF. In the early 1990s, DF suffered severe brain damage from
carbon monoxide poisoning, which destroyed a region of her lateral
occipital cortex in the ventral stream. She developed profound
\emph{visual form agnosia}: she could not identify objects by
their shape, could not tell you which of two blocks in front of her
was wider, could not indicate by gesturing with her hand how tall a
cylinder was. When shown a card inserted into a slot at different
orientations, she could not rotate her hand to match the slot's
angle. Her visual system, tested through conscious report, appeared
to have lost access to the basic geometry of objects.

Then Goodale and Milner asked her to post the card through the slot.
She did so without difficulty, rotating her hand accurately to match
the orientation of the slot as she moved. When asked to reach out
and pick up a block, her hand opened to precisely the right width
as it approached, calibrated to the block's size throughout the
movement. The information she could not report was guiding her
actions in real time. The ventral stream, which supports recognition
and report, was damaged. The dorsal stream, which supports action
guidance, was largely intact. The two systems had been surgically
separated by the pattern of her injury.

\aside{The reverse dissociation exists as well, providing a double
confirmation. Patients with \emph{optic ataxia} --- damage to the
dorsal stream in the posterior parietal cortex --- can describe
objects accurately and recognise them normally, but cannot reach
to them correctly. Their hand misses the target or approaches at
the wrong angle, even while they can see and name exactly what they
are reaching for. Recognition intact, action guidance impaired:
the mirror image of DF.}

What the two patients together establish is not merely a curious
anatomical fact but a conceptual revision. ``Seeing'' is not one
thing. The ventral stream constructs the conscious, reportable,
categorised world of perception: the world of objects with names,
colours, and meanings. The dorsal stream constructs a parallel,
largely unconscious representation of the same scene, continuously
updated for the moment-to-moment control of movement, which is
discarded almost immediately after use rather than stored. The two
streams can be damaged independently, run in parallel, and often
operate without either being aware of the other.

\section{The controlled hallucination}

The neuroscientist Anil Seth has proposed a phrase for what these
findings add up to: \emph{controlled hallucination}. Perception, he
argues, is what hallucination would be if the sensory signals were
allowed to constrain it. A hallucination is a perceptual experience
generated entirely by the brain's internal models, uncorrected by
incoming data. Ordinary perception is the same process, but with the
world continuously supplying prediction errors that keep the model
tethered to reality. The only difference between seeing a room and
hallucinating a room is whether the sensory corrections are arriving
and being applied.

This is deliberately provocative, and its value lies in what it
implies about the perceptual experience of people for whom the
correction mechanism is disrupted. In psychosis, the balance between
prior and sensory evidence is disturbed --- the priors become too
insistent, or the sensory errors are too weakly weighted, or both
--- and the brain's model of reality begins to drift from what the
world is actually sending. Hallucinations and delusions are not
random malfunctions. They are the predictive machinery running with
inadequate correction. The hollow mask illusion is a laboratory
miniature of this: everyone, in that case, is briefly hallucinating
a convex face, and the people least susceptible to the hallucination
are the ones whose correction mechanism is most active.

The phrase also reconnects perception to the broader themes of the
atlas. Memory, as an earlier chapter showed, is reconstructive: what
we retrieve is rebuilt from fragments and expectations rather than
played back from a faithful record. Attention, another chapter showed,
determines what receives the prediction errors that update the model.
Consciousness, the chapter before this one, argued that awareness is
what happens when predictions and their corrections ignite into global
broadcast. Perception is the chapter where all of those claims have
their most immediate sensory grounding. The visual world is rebuilt
fifty times a second, from a sparse and ambiguous signal, using
expectations encoded over a lifetime and corrected by whatever the
world is currently sending back.

\section{Why the maps disagree}

Perception occupies an unusually contested position in the field's
reference atlases. Some taxonomies split it into modalities ---
visual, auditory, tactile, olfactory --- treating each sense as a
separate faculty with its own entry. Others organise it around stages
--- sensation, feature detection, object recognition, scene
understanding --- treating it as a pipeline from the raw to the
interpreted. Others subordinate it entirely to attention, arguing
that perception just is what attention selects from incoming signals.
And the predictive coding framework blurs the distinction between
perception and memory almost completely: the ``predictions'' that
the brain sends downward through the visual hierarchy are, in a
real sense, memories --- learned regularities about how the world
tends to look --- and perception is what happens when those memories
meet the world's corrections.

What the field cannot agree on, and what the maps' disagreement
faithfully records, is where inference ends and perception begins.
If perception is inference, then it shades continuously into
cognition, belief, and expectation. If it is a separate, modular
process that merely feeds cognition, then a cleaner boundary can
be drawn. The predictive coding framework implies the former.
Classical accounts of feedforward, modular perception imply the
latter. The maps embed that theoretical choice without always
acknowledging it.

What the evidence does settle, for the purposes of this atlas, is the
point with which the chapter opened: the eye is not a camera and the
brain is not a screen. The richly detailed world you seem to see
right now --- the room, the page, the objects at their various
distances --- is not a copy of anything. It is the brain's most
recently corrected guess, built from expectations accumulated over
a lifetime, constrained by the signals arriving from the world,
and assembled so quickly and so reliably that it almost never
feels like what it is.







ch7.tex

\atlaschapter{Reading Other Minds}{Underneath every conversation is a continuous act of mind-reading. What happens in the brain when you do it, why it breaks down, and why the obvious explanation for how it works turned out to be wrong.}

You have spent your entire waking life reading other minds without
knowing that is what you were doing. Every conversation involves
forming a model of what the other person believes, what they want,
what they already know that you don't need to explain. You calibrate
your words to their knowledge. You follow what they mean rather than
just what they say, which is often not the same thing. When someone
asks ``can you pass the salt?'' you do not answer yes and then sit
still; you understand it as a request, which means you have modelled
the gap between what was literally said and what was intended. You
understand irony, understatement, implication. You notice when someone
is pretending to agree while disagreeing. You predict, with modest
accuracy, what a friend will want for lunch, what a colleague will
object to in a meeting, what a stranger reading this sentence is
currently thinking.

All of this is so continuous, so fast, and so accurate in its usual
operation that it does not feel like computation. It feels like just
talking to people, just knowing what they mean. The machinery is
invisible.

This chapter is about the machinery. It turns out to be more
anatomically distinct, more developmentally timed, and harder to
explain mechanistically than almost anyone expected. It also turns
out to be woven through a controversy --- involving one of the
most celebrated and then most criticized findings in modern
neuroscience --- that illustrates, in compressed form, how the
field navigates the distance between a genuine discovery and the
theory built on top of it.

\section{The task that turned out to be hard}

In 1978, the psychologists David Premack and Guy Woodruff asked whether
chimpanzees have a ``theory of mind'' --- whether they could attribute
mental states to others and use those attributed states to predict
behaviour. The phrase they coined became the name of a field.

Having a theory of mind means understanding that other people have
beliefs, desires, and intentions that may differ from your own ---
that the world as Sally sees it is not necessarily the world as it is.
It is the capacity to represent a representation: to hold in mind not
just the state of the world but someone else's model of the state
of the world, which may be wrong. Logicians call this a second-order
belief. Developmental psychologists call it the ability to pass the
false-belief task. The rest of us call it knowing that someone can be
mistaken.

The first clean test was designed by Heinz Wimmer and Josef Perner
in 1983. A child watches a short story: Maxi puts his chocolate in
a blue cupboard and goes out to play. While Maxi is away, his mother
moves the chocolate to a green cupboard. Maxi comes back. Where will
Maxi look for his chocolate? The correct answer is the blue cupboard,
because that is where Maxi \emph{believes} the chocolate is, and Maxi
does not know it has been moved. But a child who cannot model Maxi's
false belief will answer with the green cupboard, because that is
where the chocolate \emph{is}. The child answers from their own
perspective rather than from Maxi's.

Wimmer and Perner found a sharp developmental boundary. Children under
about four years of age reliably gave the wrong answer --- the one from
their own perspective. Children above four increasingly gave the right
one. Something changes around the fourth birthday that enables the
child to track another person's beliefs separately from their own.

Two years later, Simon Baron-Cohen, Alan Leslie, and Uta Frith adapted
the task into what became the Sally-Anne test, with two dolls and a
marble. They were not primarily interested in development. They were
testing a hypothesis about autism: that autistic children might
specifically lack the capacity to attribute false beliefs. Their
results were striking. Eighty-five per cent of typically developing
children and eighty-six per cent of children with Down syndrome
--- the latter included as a control for general developmental delay ---
passed the test. Only twenty per cent of the autistic children
did. The pattern suggested that the false-belief capacity was
not simply delayed in autism but specifically impaired, and that
the impairment was not a consequence of general intellectual
disability.

\aside{The Sally-Anne finding launched decades of research, but the
interpretation has evolved considerably. Many autistic people do pass
false-belief tasks, especially when given adequate time or explicit
verbal instruction. The capacity appears intact in a significant
proportion of the autistic population but accessed differently ---
less automatically, more effortfully --- and the picture turns out
to be substantially more complicated than ``autistic people lack
theory of mind.'' More on this below.}

What the false-belief paradigm established, which has held up
through all the subsequent revision, is that representing someone
else's false belief is a specific cognitive achievement with a
developmental timetable, not a simple byproduct of general
intelligence. Something happens at around four that makes it
possible, and whatever that something is, it can fail to happen
or happen differently.

\section{The brain's mentalizing network}

When neuroimaging arrived as a tool, researchers hoped to locate the
hardware. They found it, more cleanly than expected.

Reading other minds consistently engages a set of regions that
have come to be called the \emph{mentalizing network}, or the
theory-of-mind network. Its core nodes are the
temporoparietal junction (TPJ), a patch of cortex at the
meeting of the temporal and parietal lobes; the medial prefrontal
cortex (mPFC), a midline region toward the front of the brain
involved in self-referential and social thinking; and the posterior
superior temporal sulcus, associated with tracking the movements
and intentions of agents. The network activates consistently across
an impressive range of tasks: reading stories that require inferring
what a character believes; interpreting indirect speech acts; detecting
faux pas (statements that reveal the speaker doesn't know something
they should know about the listener); judging whether an action was
intended or accidental. What the tasks share is the need to model
someone else's mental state, and where they overlap in the brain
is the mentalizing network.

Rebecca Saxe at MIT demonstrated that the right TPJ in particular
responds more strongly to stories about beliefs than to stories
requiring equivalent reasoning about physical states. The same
complexity, the same narrative structure, the same logical demands
--- but if the story requires you to infer what someone \emph{thinks}
rather than what a physical system \emph{does}, the right TPJ lights
up specifically. This is not simply a region for social information
in general. It is, as specifically as we can localize with fMRI,
a region for thinking about the content of other minds.

The network is dissociable from the circuits used for logical
reasoning in general, and from those used for making physical
predictions. This matters because it rules out a simple account
on which mentalizing is just hard reasoning of a particular kind.
Something in the brain treats the problem of other minds as a
distinct computational domain, not a special case of ordinary
inference.

\section{Mirror neurons: discovery and overshoot}

In the early 1990s, Giacomo Rizzolatti and Vittorio Gallese at the
University of Parma were recording the activity of individual neurons
in the premotor cortex of macaque monkeys --- specifically neurons
that fire when a monkey reaches for and grasps an object. They
noticed something unexpected. Before the formal experiment began,
Gallese reached out to pick up an object to show it to the monkey.
The neurons fired. Not because the monkey had moved, but because
the monkey had \emph{seen} the movement. The same cells that encoded
the motor action of grasping also fired in response to observing the
same action performed by someone else.

They called them mirror neurons. The cells seemed to create a bridge
between the self and other: the same neural code that represented
your own action represented another's action in your brain. By 1996,
when the finding was formally published with Rizzolatti as lead
author, its implications seemed profound. If the brain encodes
observed actions using the same cells as performed actions, perhaps
this was how we understand what others are doing --- by simulating
their movement in our own motor system, recognizing it from the
inside.

What happened next is a case study in how a genuine scientific
finding can be inflated by a single good metaphor. The V.~S.
Ramachandran who had introduced the atlas to phantom limbs called
mirror neurons ``the neurons that shaped civilization'' and proposed
that they were the neural basis of empathy, language acquisition,
cultural transmission, and imitation. The popular coverage amplified
these claims further. Mirror neurons became the explanation for
almost everything distinctively human.

The scepticism was not slow in arriving, but it was slower than
the hype. Several problems accumulated. First, mirror neurons had
been directly recorded in macaques at the single-cell level; in
humans, they had not --- for the obvious reason that recording
from individual neurons requires electrodes implanted in the brain,
a procedure not performed for research purposes. The evidence for
a human mirror system came from functional imaging, which measures
regional blood flow rather than the firing of individual cells
and cannot confirm that any given neuron has true mirror properties.
When rare single-neuron recordings were eventually obtained in humans
--- from patients undergoing depth electrode placement for epilepsy
surgery, in a 2010 study --- the findings were suggestive but messier
than the clean macaque results: some cells showed mirror-like
properties, but the selectivity and organisation were less clear.

Second, and more damaging to the grand theory, the ``broken mirror''
hypothesis of autism --- that autistic people's social difficulties
arise from mirror neuron dysfunction --- failed empirical test after
empirical test. The hypothesis predicted that autistic individuals
should show reduced activity in mirror system regions when observing
actions. Studies found either no difference or inconsistent and
small differences. A high-profile 2009 review concluded that autistic
people show no particular impairment in action understanding or
imitation under explicit instruction, which were supposed to be the
abilities that mirror neurons support. The broken mirror hypothesis was not supported by the data.

\aside{Rizzolatti himself, the original discoverer, ultimately stepped
back from the broader claims. He noted that mirror neurons might
be necessary for understanding actions from a first-person perspective
--- experiencing another's movement as if from inside --- but that
several mechanisms for understanding action probably coexist, and
that mirror neurons likely constitute one component of a larger
system rather than the whole explanation. A discovery scaled back
to its actual scope; a useful outcome of the period of overshoot.}

What remains, stripped of the excess claims, is genuine and
interesting. Mirror neurons exist in macaques. Some analog likely
exists in humans. The system probably contributes to recognizing
and predicting actions that are within one's own motor repertoire.
It is plausibly one input into the larger process of social cognition.
What it almost certainly does not do is explain empathy, language,
culture, or theory of mind by itself. The mentalizing network that
supports false-belief reasoning is largely distinct from the
motor regions where mirror neurons were discovered. Knowing what
someone is doing is a different computation from knowing what
they believe, and the brain appears to keep those computations
in different places.

\section{The autism revision}

The original Baron-Cohen and colleagues finding --- that autistic
children fail false-belief tasks at high rates --- was real and
replicable. The interpretation that autism is fundamentally a theory-
of-mind deficit became deeply influential, and the label
``mindblindness'' that Baron-Cohen later applied to it entered wide
circulation. The picture turned out to be considerably more
complicated.

Several lines of evidence have pushed back. Many autistic people
pass false-belief tasks, particularly when they are given adequate
time, explicit instruction, or when the task is presented in a
format that does not require rapid, automatic inference. This
suggests that the underlying capacity may be present but deployed
differently --- less reflexively, more deliberately, requiring
conscious effort where typically developing people achieve it
without trying. Performance varies enormously across the autistic
population in a way that makes any single characterization
unreliable.

More fundamentally, a different framing of the social difficulties
associated with autism has gained traction in recent years. The
researcher Damian Milton proposed what he called the
\emph{double empathy problem}: the observation that social
communication difficulties between autistic and non-autistic people
are bidirectional, not one-sided. Non-autistic people are just as
poor at reading autistic social signals as autistic people are at
reading non-autistic ones. The breakdown is at the interface between
two different social styles, not a deficit within one party. When
autistic people interact with other autistic people, the difficulties
are substantially reduced; the inference problems emerge primarily
at the cross-neurotype encounter.

This reframing matters because it changes the question. Instead of
asking ``what is wrong with autistic social cognition?'' it asks
``what happens when two minds that work differently try to model
each other?'' The difficulty then looks less like a disorder and more
like a communication mismatch, analogous to --- if not precisely
equivalent to --- what happens when people with genuinely different
cultural assumptions about directness, eye contact, or turn-taking
try to interact. Each group's model of the other is systematically
miscalibrated, not because either group lacks the capacity to model
minds but because the social signals don't translate.

This does not dissolve the original finding. There are aspects of
mentalizing that are genuinely different in autism, and the
neurological evidence for differences in how the mentalizing network
is recruited is substantial. But the nature and meaning of those
differences are more contested and more nuanced than the early
``theory of mind deficit'' framing implied, and the double empathy
account asks the field to take seriously the possibility that the
test itself was designed by and for one neurotype, and was measuring
something partly social rather than purely cognitive.

\section{Simulation or theory?}

Running through all of this is a question that the experimental data
has not yet settled: how does mind-reading actually work?

Two accounts have been in sustained competition since the late
1980s. The first, called \emph{theory theory}, holds that we
understand other minds by deploying a tacit folk-psychological
theory --- a set of general principles about how mental states
arise and how they connect to behaviour. ``She wants X. She believes
Y is the way to get X. So she will do Y.'' The inference from belief
and desire to action is rule-governed, and what the child acquires
around the fourth birthday is a more sophisticated version of the
rules, particularly the rule that beliefs can be false. On this
account, understanding other minds is continuous with understanding
other systems: we use general inferential machinery applied to a
particular domain.

The second account, \emph{simulation theory}, holds that we do not
theorize about other minds from the outside but simulate them from
the inside. To understand what someone will do, we put ourselves in
their position, run a pretend version of their situation through
our own mental machinery, and observe the output. The output of
the simulation is the prediction. If this is right, understanding
other minds is fundamentally self-referential: we understand others
by imagining ourselves as them, and the predictive accuracy depends
on how similar their mental machinery is to our own.

Mirror neurons were initially interpreted as evidence for simulation
theory --- they seemed to show the simulation happening at the
neural level, with the observer's motor system literally running the
observed action. But the mirror neuron setback did not settle
the debate, because simulation theory does not depend on mirror
neurons specifically. Simulation could be implemented by any system
that can run ``offline,'' in an as-if mode, without the simulation's
outputs reaching muscles and actually producing action.

The debate remains active, and probably the honest answer is that
both processes operate to some degree. Some mind-reading seems to
use something like simulation: imagining how you would feel in
someone's situation, projecting your own emotional response.
Other mind-reading seems to use something like theory: reasoning
from general principles about what people with certain beliefs
and desires tend to do. The mix likely varies by context, by
relationship, and by the specific mental state being attributed.
What neither account yet explains with full satisfaction is why
the machinery is anatomically distinct from general reasoning ---
why the brain treats the problem of other minds as a special
domain at all.

\section{Why the maps disagree}

Social cognition sits at the intersection of so many disciplines ---
developmental psychology, clinical science, cognitive neuroscience,
evolutionary biology, philosophy of mind --- that it is represented
in the field's atlases as almost different things depending on who
built the atlas.

A developmental atlas organises social cognition around the
emergence of theory of mind as a milestone: a capacity that appears
at a specific age, can be measured with a specific task, and whose
absence or delay is diagnostically significant. A cognitive
neuroscience atlas organises it around the mentalizing network:
a set of brain regions recruited by a family of tasks that share
the property of requiring mental state attribution. A clinical atlas
organises it around conditions in which social cognition breaks
down --- autism, schizophrenia, certain personality disorders, social
anxiety --- and around the functional consequences of those breakdowns.
An evolutionary atlas asks how mind-reading capacities differ across
species and what selective pressures might have produced them.

Each framing is legitimate. Each generates useful research. And each
draws the boundaries of ``social cognition'' differently, which is
why the catalogues don't nest consistently: what one calls a
cognitive capacity, another calls a developmental milestone, another
calls a clinical symptom, another calls an evolved module.

What all of them agree on, and what the chapter has tried to
establish, is that reading other minds is not a free bonus that
comes with having a mind. It is a specific, effortful, neurologically
distinct, developmentally acquired capacity that can be present in
degrees, can be deployed more or less automatically, can be
systematically disrupted, and can be miscalibrated in ways that
depend as much on the pair of minds attempting the exchange as on
either mind alone. The miracle is not that it sometimes fails. The
miracle is that the continuous, rapid, largely accurate reconstruction
of other people's inner lives that every conversation requires works
as well as it usually does.










ch8.tex

\atlaschapter{Language}{Hearing a sentence and understanding it feel like a single act. They are not --- and where the seam falls, between sound and meaning, is the oldest unsettled question in the science of language.}

Someone says your name across a room and, before you have decided to do
anything, you have already done a great deal. A pressure wave has
travelled through the air, vibrated your eardrum, been transduced into
neural firing by the cochlea, and arrived at your auditory cortex as a
smear of frequencies changing over a few hundred milliseconds. From that
smear --- which is, physically, just a structured disturbance of air ---
you have extracted a string of speech sounds, grouped them into a word,
matched the word to a meaning, recognised the meaning as your own name,
and felt the small jolt of being addressed. None of it felt like work.
It felt like simply hearing your name.

The experience of understanding speech is seamless in exactly the way
the experience of seeing a room is seamless, and for the same reason:
the machinery is invisible to the mind it serves. You are not aware of
decomposing the sound stream into phonemes, because by the time anything
reaches awareness the decomposition is already done. You do not feel
yourself retrieving the grammar that tells you ``the dog bit the man''
and ``the man bit the dog'' describe different events built from the
same five words. The seam between hearing a sentence and understanding
it is hidden so well that for most of human history nobody suspected
there was a seam at all.

There is a seam, and finding exactly where it falls --- where the
processing of sound ends and the processing of language begins --- has
turned out to be one of the genuinely hard problems in cognitive
science, hard enough that the field's reference works cannot agree on
how to draw the boundary, or whether language is a distinct faculty at
all or simply what general-purpose cognition does when the input happens
to be words. This chapter is about that seam, and about the two
discoveries --- one in a Paris hospital in 1861, one in a linguistics
seminar a century later --- that between them made language a science.

\section{The man who could say only one word}

The science of language begins, as the science of memory does, with a
single patient whose brain broke in an informative way. In 1861 a
fifty-one-year-old man named Louis Victor Leborgne was admitted to the
surgical ward of a hospital outside Paris, under the care of a young
physician and anthropologist named Paul Broca. Leborgne had been losing
the power of speech for two decades. By the time Broca met him, he could
produce essentially one syllable --- \emph{tan} --- which he repeated for
everything, usually twice, sometimes with a curse attached when
frustrated. The ward had taken to calling him Tan.

What made Tan extraordinary, and what Broca grasped immediately, was the
shape of what he had lost. He was not unintelligent. He understood what
was said to him, followed instructions, tracked the passage of time,
answered questions with gestures and the number of his single syllable.
His vocal apparatus was intact; he could move his tongue and lips. What
had been destroyed was specifically the capacity to assemble speech ---
to turn an intact understanding and an intact mouth into spoken words.
He died days after Broca examined him, and the autopsy revealed a lesion
in the left frontal lobe, toward the front and bottom, in a region now
called Broca's area. A second similar case followed within months.
Broca drew the conclusion that founded a field: a specific part of the
brain is responsible for the production of articulate speech, and
damage to it produces a specific and selective deficit.

The deficit has a name now --- Broca's aphasia --- and a characteristic
profile. Speech is effortful, halting, produced in short bursts with
visible struggle. It is \emph{agrammatic}: the small connective
machinery of language, the articles and prepositions and verb endings,
drops away, leaving a telegraphic residue of content words. A patient
asked about a weekend might manage ``son\ldots visit\ldots Sunday\ldots
good,'' meaning it and knowing it is not coming out right, often acutely
frustrated by the gap between the intact thought and the broken output.
Comprehension, by contrast, is largely preserved, at least for ordinary
conversation. The understanding is in there. The route from
understanding to speech is what the lesion has cut.

Thirteen years later, a young German neurologist named Carl Wernicke
described the mirror image, and in doing so proved that Broca had found
a part rather than the whole. Wernicke's patients had damage further
back, in the posterior part of the left temporal lobe, in a region now
called Wernicke's area. Their speech was the opposite of Tan's: fluent,
effortless, produced at a normal rate with normal melody and normal
grammatical scaffolding --- and almost entirely empty of meaning. The
words came easily and were the wrong words, or invented words, strung
into sentences that had the shape of sense without the content. A
patient might say, asked how they were feeling, something like ``I feel
very well, the thing is going around with the\ldots you know, the
\emph{flonнer},'' delivered with complete fluency and no apparent
awareness that anything was amiss. And crucially, these patients could
not understand speech either. The same region, damaged, took out both
comprehension and the meaningfulness of production.

\aside{The two aphasias are a dissociation of the same logical force as
H.M.'s mirror-tracing hand. If one patient loses fluent speech while
keeping comprehension, and another loses comprehension while keeping
fluent speech, then production and comprehension cannot be the same
system wearing two hats. They live in different tissue, and either can
be destroyed while the other runs on. ``Language,'' the single word,
names a committee --- as memory did, and attention did, and perception
did before it.}

\section{The classic model, and its cracks}

Out of Broca's and Wernicke's cases a model was built that dominated the
science of language for a century, and that still structures the way
aphasia is taught and diagnosed. Wernicke himself, and after him the
neurologist Ludwig Lichtheim, drew it as a diagram of boxes and arrows:
a centre at the back, in Wernicke's area, where the sound-images of
words and their meanings were stored; a centre at the front, in Broca's
area, where the motor programmes for producing words were assembled;
and a bundle of fibres connecting the two, the arcuate fasciculus,
carrying information from the comprehension centre to the production
centre.

The model made a prediction that turned out to be correct and that
sealed its prestige. If the connecting bundle were cut while both
centres survived, you should get a patient who could understand speech
(Wernicke's area intact) and produce fluent speech (Broca's area
intact) but could not \emph{repeat} what was said to them --- because
repetition requires routing the heard word from the back of the brain
to the front, and the route is severed. Such patients exist. The
condition is called conduction aphasia, and its existence, predicted
from the wiring diagram before it was confirmed in the clinic, looked
like proof that the diagram was essentially right.

It was not essentially right, or at least not in the clean way it was
taught. The cracks took a long time to surface because the original
evidence was so famous that few thought to re-examine it. Then, in 2007,
a team led by the neuroscientist Nina Dronkers did something nobody had
done: they retrieved the preserved brains of Broca's two original
patients, Leborgne among them, from a Paris museum and put them through
a modern MRI scanner. The lesion in Tan's brain was not a tidy hole in
Broca's area. It was far larger than Broca, working from the surface of
the specimen alone, had been able to see, extending deep into the
underlying white matter and neighbouring structures. The damage that
produced Tan's profound, decades-long muteness involved a great deal
more of the brain than the small patch of cortex that bears Broca's
name.

This is the recurring pattern of the atlas in another costume. The
nineteenth-century localisers had found something real --- damage here
disrupts production, damage there disrupts comprehension --- and had
then drawn a map cleaner than the territory, a map of discrete centres
joined by single cables. The modern picture replaces the centres with a
distributed network: language is supported by a large, interconnected
system spread across the left hemisphere (and, more than the classic
model allowed, the right), in which Broca's and Wernicke's areas are
important hubs rather than self-contained organs, and in which the white
matter bundles connecting regions matter as much as the regions
themselves. The boxes were never quite boxes. They were the densest
knots in a web.

\section{Two streams, again}

When the distributed picture was given a modern theoretical shape, it
arrived in a form that should feel familiar to any reader of the
Perception chapter. There, the visual system turned out to split into
two streams --- a ventral ``what'' stream running into the temporal
lobe and supporting conscious recognition, and a dorsal stream running
into the parietal lobe and supporting the real-time control of action,
the two dissociable by injury, as patient DF and the optic-ataxia
patients showed. Language, it turns out, is organised on strikingly
similar lines.

In an influential framework developed by the neuroscientists Gregory
Hickok and David Poeppel, the processing of speech divides into two
streams that diverge after the initial analysis of the sound. A
\emph{ventral stream}, running forward along the temporal lobe, maps
sound onto meaning --- it is the route by which an acoustic signal
becomes a word you understand. A \emph{dorsal stream}, running upward
and back toward the parietal lobe and then forward to the frontal motor
regions, maps sound onto articulation --- it is the route by which a
heard word becomes a pattern of movements your own mouth could make.
One stream is for comprehending; the other is for the sensorimotor work
of producing and repeating.

The parallel to vision is not a loose analogy. It is the same deep
design principle showing up in a second sense: the brain builds one
pathway to tell it \emph{what something is} and a separate pathway to
tell it \emph{what to do about it}, and the two can be damaged
independently. Conduction aphasia --- fluent, comprehending, unable to
repeat --- is, on this account, a dorsal-stream disorder: the
sound-to-articulation route is down while the sound-to-meaning route
survives. The classic syndromes of nineteenth-century neurology fall
out of the two-stream architecture as naturally as the visual syndromes
fell out of theirs.

\aside{The auditory atlas had quietly predicted this. The computational
cognitive architectures that contribute to the field's reference works
model audition, like vision, with separate \emph{what} and \emph{where}
buffers --- a ventral pathway for identity, a dorsal pathway for
location and action. The dual-stream model of language is what that
abstract architecture looks like when the auditory objects in question
are words. The same skeleton, three times over: vision, audition,
language.}

\section{Is speech special?}

We have been talking as though there were an obvious first step ---
``the initial analysis of the sound'' --- before the streams diverge.
This is exactly the step that hides the seam, and the question of what
happens in it is the deepest carving conflict in the whole science of
the mind: where does \emph{hearing} end and \emph{language} begin? Is
the perception of a speech sound an act of audition, the ear doing its
ordinary job on a particular kind of noise? Or is it already an act of
language, a process so specialised for speech that it runs on different
machinery from the perception of any other sound?

The case that speech is special begins with a strange and robust
phenomenon called \emph{categorical perception}. The physical
difference between the sounds ``ba'' and ``pa'' is, at bottom, a matter
of timing: how long after the lips part the vocal cords begin to buzz, a
quantity called voice onset time that varies smoothly and continuously.
You can synthesise a whole ladder of sounds with voice onset times
stepping evenly from a clear ``ba'' to a clear ``pa.'' What listeners
hear, though, is not a smooth ladder. They hear ``ba,'' ``ba,'' ``ba''
--- and then, at some boundary, abruptly ``pa,'' ``pa,'' ``pa,'' with
almost nothing in between. Two sounds straddling the boundary, though
physically no more different than two sounds well inside ``ba,'' sound
categorically distinct; two sounds inside a category, though physically
just as far apart, sound nearly identical. The continuous acoustic
world has been sorted into discrete linguistic bins before it reaches
awareness. Researchers at the Haskins Laboratories, who first
documented this in the 1950s, took it as evidence that speech is
perceived by a dedicated system --- one tuned not to raw sound but to
the units of language.

The most theatrical demonstration that speech perception is not simple
hearing is the \emph{McGurk effect}, and it has the rare property among
psychological phenomena that knowing about it provides no defence
against it. Take a video of a person's mouth saying ``ga.'' Dub onto it
the audio of the same voice saying ``ba.'' What you will hear --- not
deduce, not infer, but consciously hear --- is neither: it is ``da,'' a
compromise the brain manufactures by fusing the visual evidence of the
lips with the acoustic evidence of the ear. Close your eyes and it
reverts to ``ba.'' Open them and ``da'' returns. Your perception of a
speech sound is being constructed from the eyes as much as the ears,
which is not something one expects of a system that is merely
\emph{listening}. Speech perception reaches across the senses for
evidence about what was said.

\aside{The McGurk effect is the Perception chapter's argument arriving
in a new modality. There, perception was inference: the brain's best
hypothesis about what in the world produced the sensory signal, built
from priors and corrected by evidence. The hollow mask let a visual
prior override the senses; the McGurk effect lets visual evidence
override the ears. And the phoneme-restoration effect completes the
parallel --- replace a speech sound in a recorded sentence with a cough,
and listeners do not hear a gap. They hear the missing sound, supplied
by the sentence's meaning, and often cannot say where the cough fell.
Hearing speech, like seeing a room, is a controlled hallucination
tethered by what the world sends.}

So is speech special --- a perceptual act run by a language-specific
system --- or is it ordinary audition feeding a separate linguistic
stage downstream? The honest answer is that the field has argued this
for half a century without resolving it, and the argument is what the
reference atlases are quietly recording when they file the perception of
phonemes in contradictory places. The phenomena point both ways.
Categorical perception and the McGurk effect say speech is special.
But categorical-like perception turns up for some non-speech sounds, and
some non-human animals show speech-like categorical boundaries, which
say it is not so special after all. The seam between hearing and
language is real, and it is genuinely unclear which side of it the
phoneme sits on.

\section{The rules nobody taught you}

While neurologists were mapping where language lives in the brain, a
different revolution was redefining what language \emph{is} --- and it
came not from a clinic but from the structure of grammar itself. Its
author was the linguist Noam Chomsky, and the moment is often dated to
1959, when Chomsky published a review of a book by the behaviourist
psychologist B.~F. Skinner that argued language was a learned behaviour
like any other, built up from experience by reinforcement, a child
shaped into speech the way a pigeon is shaped into pecking. Chomsky's
review was a demolition, and it did more than refute Skinner; it helped
end the behaviourist era in psychology and inaugurate the cognitive one.

Chomsky's central argument was about a poverty. The speech a child
hears, he pointed out, is finite, fragmentary, and full of errors,
false starts, and interruptions. And yet every normally developing
child, from any such impoverished and idiosyncratic sample, converges in
a few years on a rich, rule-governed grammatical system --- one capable
of generating and understanding an unbounded number of sentences the
child has never heard before, including sentences no one has ever
uttered. Children make systematic errors that reveal they are applying
rules rather than imitating (``I goed,'' ``two foots''), errors no adult
modelled for them. They acquire grammatical distinctions that the
available evidence radically underdetermines. The input, Chomsky argued,
is too poor to account for the achievement. This is the
\emph{poverty of the stimulus}, and its conclusion is that the child
must be bringing something to the task --- a substantial innate
endowment, a \emph{universal grammar}, a set of constraints on possible
human languages that the child does not learn but is born with, and
that turns the meagre evidence into a full grammatical system.

On this view language is not a general-purpose skill, like chess or
arithmetic, that a clever animal could acquire with enough exposure. It
is a species-specific mental organ, as much a part of the standard human
endowment as the visual system, growing on a developmental schedule the
way an organ grows, requiring environmental input to develop but not
being \emph{caused} by that input any more than nutrition causes the
shape of the heart. The philosopher Jerry Fodor gave the idea a more
general frame, arguing that the mind contains a number of
\emph{modules} --- systems that are domain-specific (each handling one
kind of problem), fast, automatic, and \emph{informationally
encapsulated}, meaning they do their work without consulting the rest of
what you know. Language was his prime example. The reason the McGurk
effect survives your knowing about it, on this account, is that the
speech module is encapsulated: it computes ``da'' from the mismatched
inputs and hands the result to consciousness, and your knowledge that
the audio is really ``ba'' sits outside the module and cannot reach in
to correct it. The module is special, sealed, and not under your
control.

\section{Or the rules you learned from listening}

The Chomskyan picture --- innate, modular, special --- has been the most
influential idea in the modern study of language, and it has been
contested at every point by a tradition that takes the opposite view:
that language is learned from the input after all, by powerful but
general-purpose learning machinery, and that the apparent poverty of the
stimulus has been overstated because the richness of what an infant can
extract from the stimulus was underestimated.

The single most striking piece of evidence for this side arrived in 1996,
from the psychologists Jenny Saffran, Richard Aslin, and Elissa Newport.
They confronted a problem that seems, on its face, to favour the
nativists: the speech stream contains no reliable gaps between words. We
hear spaces between words because we know the words; an infant, or an
adult hearing an unfamiliar language, confronts a continuous run of
sound with no marks showing where one word stops and the next begins.
How could a learner ever break in? Saffran and colleagues exposed
eight-month-old infants to two minutes of a monotone, unbroken stream of
invented syllables, in which the only cue to word boundaries was
statistical: within a ``word,'' one syllable predicted the next
reliably, while across a boundary, the next syllable was unpredictable.
Two minutes. After it, the infants distinguished the ``words'' of this
artificial language from equally familiar syllable sequences that
straddled boundaries --- which they could only do if, in two minutes,
they had tracked the transitional probabilities between syllables and
used the dips to find the seams. Infants are statistical learners of
formidable power, extracting structure from raw input through
general-purpose machinery, with no word-specific instruction at all.

This is the empirical core of the \emph{usage-based} approach to
language, associated with the developmental psychologist Michael
Tomasello and others, which holds that children build their grammar
gradually out of the concrete utterances they hear, generalising from
stored examples toward increasingly abstract patterns, using the same
cognitive capacities --- pattern detection, analogy, the drive to read
others' intentions --- that they use everywhere else. There is no
dedicated language organ on this account, no innate universal grammar.
There is a richly structured input, a socially motivated learner, and
learning mechanisms that are powerful and general rather than special
and sealed. The fact that the language faculty looks distinct in the
brain is, to a usage-based theorist, what you would expect from any
skill practised relentlessly from infancy --- the marks of expertise,
not the fingerprints of an innate module.

\aside{One observation neither camp can wave away is the
\emph{critical period}. The capacity to acquire a first language to
native fluency appears to have a developmental window that closes. The
clearest and most painful evidence comes from cases of extreme
deprivation --- children, studied in the 1970s and after, who reached
adolescence with almost no exposure to language and who, despite
intensive subsequent teaching, acquired vocabulary but never normal
grammar. These cases are ethically fraught and scientifically
imperfect, confounded by the broader devastation of the deprivation that
produced them, and no honest account leans on them too hard. But the
softer evidence converges: second languages learned after childhood are
almost never spoken without an accent and a residue of grammatical
effort, where those learned early are. Something about the language
faculty is timed, and the timing constrains both stories.}

\section{The gene that wasn't}

There is one more episode worth telling, because it rhymes so exactly
with the mirror-neuron story of the previous chapter --- a genuine
discovery, a seductive metaphor, an overshoot, and a sober correction
--- that it has become the cautionary tale of language genetics.

In the 1990s a large London family, known in the literature as the KE
family, came to scientific attention because roughly half its members
across three generations had a severe inherited disorder of speech and
language: difficulty sequencing the mouth movements of articulation,
together with broader problems in grammar and language processing. The
inheritance pattern was clean enough to suggest a single gene, and in
2001 it was found --- a mutation in a gene named FOXP2 on chromosome 7.
The press, and some scientists, reached immediately for the obvious
phrase. FOXP2 was ``the language gene,'' the genetic basis of the human
capacity for speech, perhaps the very mutation that separated us from
our speechless ancestors.

It is nothing so tidy. FOXP2 is not a language gene in any meaningful
sense. It is a regulatory gene --- one that switches other genes on and
off --- and it is active in many tissues and many species. Versions of
it are found in mice, where it shapes the development of motor circuits,
and in songbirds, where it is involved in song learning, and across the
mammals, where it is highly conserved. What its disruption damages in
the KE family is, at the core, the fine motor sequencing that
articulation requires, with linguistic consequences downstream of that.
FOXP2 is one gene among many that must work for language to develop
normally, contributing to the construction of neural circuits that
language happens to use --- not a blueprint for grammar, not a switch
that turns language on, and certainly not the lone key to what makes us
linguistic. The grand version of the claim has been retired by the
researchers closest to the work, and what survives is real and modest:
a gene whose disruption reveals how much ordinary motor and
developmental machinery language is built on top of.

The lesson is the atlas's recurring one about the distance between a
finding and the theory draped over it. The discovery was sound. The
metaphor --- \emph{the} gene \emph{for} language --- was a category
error, smuggling in the assumption that a capacity this complex must
have a single dedicated cause. Language no more has a gene than memory
has a storehouse or attention has a spotlight. The convenient picture is
the thing to distrust.

\section{Why the maps disagree}

Return, finally, to the seam. The field's reference works carve language
in ways that contradict each other, and the contradictions cluster
exactly where this chapter began: at the boundary between hearing and
language, which no catalogue draws the same way twice.

Some of the atlases keep auditory perception strictly pre-linguistic ---
the ear's processing of sound, walled off from anything to do with words,
with language a separate territory downstream. Others define a broad
auditory ability as explicitly \emph{nonverbal}, excluding speech on
principle --- and then, in the same breath, file the perception of
phonemes, the hearing of the sound structure of words, as a kind of that
very ability, contradicting their own wall within a single entry. Still
others run the boundary directly through individual concepts, filing the
processing of speech sounds as simultaneously a part of hearing and a
part of language, refusing to choose because the phenomenon itself seems
to sit on both sides. One tradition treats the perception of a word's
sound as audition; another treats it as the first step of language; a
third treats it as both at once.

These are not filing errors, any more than the opposite-facing trees of
the Memory chapter were. They are an honest record of a question the
science has not answered: whether the perception of speech is a
perceptual act or a linguistic one, whether ``speech is special'' or
speech is just sound that language later makes use of. A catalogue that
filed the phoneme cleanly on one side would be claiming a resolution the
field does not have. The contradiction is the more faithful map.

Underneath the boundary dispute sits the larger one this chapter has
circled throughout --- the question that divides Chomsky from Tomasello,
Fodor's sealed module from the general-purpose statistical learner. Is
language a \emph{thing}, a distinct faculty with its own innate
endowment and its own dedicated tissue, the kind of entity that earns
its own heading in any honest taxonomy of the mind? Or is it a
\emph{label} for what happens when general cognition --- perception,
memory, the reading of other minds, the detection of statistical
pattern --- is pointed at the special problem of words? The aphasias and
the imaging say language has real neural specialisation; the
statistical-learning infants and the usage-based developmentalists say
that specialisation could be the signature of a deeply practised skill
rather than the fingerprint of an innate organ. The atlases that give
language its own grand heading and the atlases that distribute it among
audition, memory, and social cognition are not making clerical choices.
They are taking sides, mostly without saying so, in an argument that is
still alive.

What is not in dispute is the achievement itself, and it is worth
ending on its strangeness. Out of a disturbance of the air, in a few
hundred milliseconds, below the threshold of any awareness, you
reconstruct another person's meaning --- sorting the continuous into
phonemes, the phonemes into words, the words into a structure that
distinguishes the dog that bit from the dog that was bitten, and the
structure into an intention you grasp well enough to be addressed,
moved, informed, or deceived by it. Every conversation runs this
process continuously and in both directions at once, and it is so
reliable and so invisible that we mistake it for nothing at all ---
for simply hearing someone speak. The seam is hidden because the
machine that hides it is the best one evolution ever built.








ch9.tex

\atlaschapter{Emotion}{We treat feeling as the opposite of thinking --- a weather that happens to us, the enemy of good judgment. The science says nearly the reverse: that without emotion reason stalls, and that a feeling may be less something you receive than something you build.}

A feeling arrives the way weather arrives. Fear floods in before you have
decided anything --- the heart already going, the stomach dropping, the
attention yanked to the source --- and the experience is one of being
\emph{seized}, taken hold of from somewhere below the level of choice.
This is the opposite of how the previous chapters' subjects feel from the
inside. Reading a sentence, recalling a fact, attending to a face: these
seem like things you do. An emotion seems like something that is done to
you. It comes uninvited, refuses to leave on command, overrides what you
have rationally concluded you ought to feel. Small wonder that for most
of Western intellectual history emotion was cast as reason's adversary
--- the passions that good judgment must master, the heat that clouds
the cool work of thought.

Almost every part of that picture turns out to be wrong, and wrong in
ways that have taken neuroscience by surprise. Emotion is not the enemy
of reason; strip it away and reason does not get sharper, it collapses.
Emotion is not simply read off the body, the way the oldest theory had
it, nor is it simply a judgment in the head. And --- most unsettling, and
most contested --- the discrete emotions we name so confidently, the
fear and anger and disgust that feel like natural kinds carved at the
joints, may not be the brain's furniture at all. They may be something
closer to what the Perception chapter found behind vision: not things
received but things constructed, the brain's best reading of an
ambiguous signal, assembled so fast and so silently that we mistake the
construction for a fact of nature.

This chapter is about a field at war with itself over the most basic
question one can ask about feeling: whether emotions are things the brain
\emph{has} or things the brain \emph{makes}. The combatants agree on
almost none of the territory. What they have produced, between them, is
the clearest case in this atlas of two communities looking at the same
evidence and seeing incompatible worlds.

\section{The man who could not choose}

The discovery that emotion is necessary for reason began, like so much
in this book, with a brain that broke in an instructive place. The
patient, whom the neurologist Antonio Damasio called Elliot, had been a
successful professional --- a good husband, a good father, a capable
manager --- until a tumour grew in the midline above his eye sockets and
had to be removed, taking with it a portion of the ventromedial
prefrontal cortex, the patch of frontal lobe that sits behind the bridge
of the nose. The surgery was a success in the narrow sense. The tumour
was gone and Elliot recovered. By every standard test, his mind was
intact: his IQ remained high, his memory excellent, his language
unimpaired, his knowledge of the world complete. He could discuss
current events, reason through a logical problem, hold forth on the
options facing him in any situation with perfect lucidity.

His life, nonetheless, fell apart. He could no longer hold a job, because
he could not prioritise --- he would spend an entire afternoon agonising
over how to sort a set of documents, by date or size or relevance, never
arriving at a decision. He made a series of catastrophic choices,
investing his savings with a disreputable partner against all advice,
running through marriages and money. Asked to schedule his next
appointment, he could weigh the considerations for and against two dates
endlessly, producing an exhaustive cost-benefit analysis for each,
without ever being able to \emph{choose}. The intellect was pristine and
the life was a ruin, and the puzzle was why.

What Damasio noticed, watching Elliot describe the most appalling
episodes of his own downfall with the detachment of a man reading a
stranger's case file, was that the emotion had gone out of him. Elliot
discussed his losses without distress. He could narrate disaster without
feeling it. And Damasio's hypothesis, which became the
\emph{somatic-marker hypothesis}, was that this flatness was not a side
effect of Elliot's troubles but their cause. Normal decision-making, he
proposed, is not the cool calculation we imagine. When you contemplate
an option, your body registers a faint echo of how past options like it
turned out --- a small visceral lean toward or away, a ``gut feeling''
in the literal sense, a bodily signal that marks the option as good or
bad before you have consciously tallied anything. These somatic markers,
generated through the ventromedial prefrontal cortex, prune the
explosion of possibilities down to the few worth deliberating over.
Without them, every option stays live, equally weighted, and you drown,
as Elliot drowned, in an endless even-handed analysis that never tips
into action.

The laboratory demonstration was the Iowa Gambling Task. Players draw
cards from four decks; two pay well but carry occasional ruinous
penalties, two pay modestly but are safe in the long run. Healthy players
gradually shift to the safe decks --- and, tellingly, their skin begins
to produce a faint stress response when their hand hovers over a bad deck
\emph{before} they can articulate that the deck is bad. The body knows,
and steers, ahead of the conscious mind. Patients with Elliot's damage
produce no such anticipatory signal, and keep drawing from the ruinous
decks to the end, calmly, as their fortunes collapse.

\aside{The headline lesson --- that emotion is not optional equipment for
a reasoner but part of the machinery of choice --- has held up. The
specific mechanism has not held up so cleanly. When researchers
questioned Gambling Task players more carefully, many turned out to know
consciously which decks were bad earlier than the somatic-marker story
required, which means the ``body deciding before the mind knows'' may be
less than it first appeared. This is the atlas's standing caution in
another setting: a genuine and important finding, wrapped in a mechanism
more seductive and more specific than the evidence quite supports. That
emotion is necessary for good judgment survives. \emph{How} it does its
work is still being argued.}

The deeper point reverses the folk picture entirely. The fantasy of the
perfectly rational agent, the Vulcan who decides better for feeling
nothing, is not what vmPFC damage produces. What it produces is Elliot:
not a cold genius but a stalled one, possessed of every fact and unable
to want anything enough to act. Reason without emotion is not lucid. It
is paralysed.

\section{Does the body make the feeling?}

If emotion is bound up with the body --- the racing heart, the dropped
stomach, the anticipatory sweat --- then a question arises that is older
than neuroscience and still not closed: which comes first? Common sense
says you see the bear, you feel afraid, and the fear makes your heart
pound. In 1884 the psychologist William James proposed the reverse, and
the reversal is so counterintuitive that it still produces a jolt. You
see the bear, James argued, your body responds --- heart, lungs, legs ---
and your \emph{perception of that bodily response is the fear}. ``We
feel sorry because we cry,'' he wrote, ``angry because we strike, afraid
because we tremble,'' not the other way around. Strip out the bodily
changes, on this view, and nothing emotional remains --- only a pale,
neutral recognition that a bear is present. The Danish physician Carl
Lange proposed much the same thing at nearly the same time, and the idea
has carried both their names ever since.

The James--Lange theory was challenged within a generation by the
physiologist Walter Cannon, who raised objections that still have force.
The body's visceral responses, he noted, are slow --- too slow to
account for the immediacy of feeling --- and crude, much the same across
emotions that feel utterly different. A pounding heart attends fear,
rage, and joy alike; how can the perception of so blunt an instrument
produce the fine distinctions of emotional life? Cannon and his student
Philip Bard proposed instead that the brain generates the feeling and the
bodily response in parallel, neither causing the other. The debate
between these two pictures --- emotion as the perception of the body, or
emotion as a central event merely accompanied by the body --- set the
agenda for the century that followed.

A famous attempt to split the difference came in 1962, when Stanley
Schachter and Jerome Singer proposed that an emotion is bodily arousal
\emph{plus} a cognitive label. Inject volunteers with adrenaline, their
theory held, and the racing heart they feel will become euphoria or anger
depending entirely on the social situation you place them in --- the
arousal supplies the heat, the context supplies the name. It is a tidy
idea and it entered every textbook. It has also proved stubbornly hard to
reproduce; the original study's effects were weaker and less reliable
than its fame suggests, and the two-factor theory now sits in the
uncomfortable position of being canonical and not well replicated at the
same time.

The cleanest modern test of the bodily-feedback idea took aim at its
simplest prediction: if the body's expression feeds the feeling, then
arranging your face into a smile should make you feel happier. In 1988 a
classic study had people hold a pen either in their teeth, which forces
the mouth into a smile, or in their lips, which prevents one, and rate
cartoons for funniness. The teeth group --- secretly smiling --- found the
cartoons funnier. The result was elegant, widely taught, and apparently
strong support for James and Lange. Then, in 2016, seventeen laboratories
ran the experiment again in a coordinated replication, and the effect
largely evaporated. Later work has muddied even the failure --- some
analyses recover a small effect under some conditions --- so the honest
verdict is not that facial feedback is dead but that a result everyone
had filed as settled turned out to be fragile, and the question James
asked remains genuinely open.

\aside{The body has not been exonerated, only complicated. The most
durable descendant of James's insight is the study of
\emph{interoception} --- the brain's sensing of its own internal state,
the heartbeat and breath and gut --- which the previous chapters' patient
S.M. will shortly show can generate raw fear all on its own. James was
probably wrong that the feeling simply \emph{is} the perceived body. But
he was onto something his critics missed: that the body is not a
bystander to emotion. It is, at minimum, a major source of the signal
the brain is trying to read.}

\section{The woman who knew no fear}

No structure has carried more of the popular weight of emotion than the
amygdala, a pair of almond-shaped clusters deep in the temporal lobes,
and no structure better illustrates how a real discovery hardens into an
over-tidy slogan. The amygdala became, in textbook and headline alike,
``the fear centre'' --- the brain's alarm, the seat of dread. The story
had a near-perfect human illustration.

A woman known in the literature as S.M., afflicted by a rare genetic
condition that had destroyed the amygdala on both sides of her brain,
appeared to have lost the capacity for fear. Researchers took her to
exotic-pet stores, where she handled snakes and spiders she had claimed
to avoid, drawn to them with an unguarded curiosity. They took her
through one of the most aggressive haunted-house attractions in the
country, where she led the group, laughing, and startled the actors
rather than the reverse. They showed her films designed to terrify and
combed her life history. Across years of study she reported, and showed,
essentially no fear. She could not reliably recognise fear in other
people's faces, either --- her gaze skipped over the wide, frightened
eyes that the rest of us read instantly. Here, it seemed, was the proof:
remove the amygdala and you remove fear itself.

And then the story broke open from inside. Researchers had S.M. and two
other patients with the same bilateral amygdala damage breathe air
enriched with carbon dioxide --- a procedure that produces a sensation
of suffocation and reliably triggers panic. All three, the supposedly
fearless among them, panicked. S.M., who had not been frightened by
snakes or haunted houses or threats from the external world, was
flooded with terror by a threat arising \emph{inside her own body}. The
amygdala, it turned out, was not the seat of fear. It was a detector of
threats in the environment, routing alarming sights and sounds into a
fast protective response --- and when the threat came not from the world
but from the blood, from a body suffocating, fear arrived without it.

\aside{The neuroscientist Joseph LeDoux built much of the original
amygdala-and-fear story, mapping in exquisite detail the circuit by which
a threatening stimulus drives the body's defensive response. He has spent
recent years arguing, with some urgency, against the way his own work
gets used. The amygdala circuit, he insists, detects threat and
orchestrates the body's reaction --- freezing, the spike in heart rate,
the surge of stress hormones --- and all of this can run without any
conscious feeling at all. The \emph{feeling} of fear, the terror you are
aware of, is something else: a higher-order construction assembled
elsewhere in the cortex, when the brain becomes aware of its own
defensive state and labels it. He would prefer we stop calling the
amygdala circuit a ``fear'' circuit altogether. A researcher walking back
the tidy version of his own most famous discovery is, by now, a familiar
figure in these pages --- the mirror-neuron enthusiasts, the FOXP2
``language gene.'' The pattern is the field correcting its own
metaphors.}

LeDoux's distinction matters for everything that follows, because it
prises apart two things the slogan had fused: the machinery that
\emph{responds} to threat, which is ancient, fast, shared across the
animals, and largely unconscious; and the \emph{feeling}, the conscious
experience of being afraid, which may be a different and later
achievement, assembled rather than triggered. Once that seam is visible,
the central question of the chapter comes into focus. When you feel fear,
is a dedicated fear system delivering its characteristic product to
consciousness? Or is your brain doing something more like interpretation
--- taking a churning body and a dangerous situation and \emph{building}
the fear, the way it builds the seen world?

\section{What dopamine is for}

Before that question can be joined properly, one more tidy story needs
dismantling, because it is the chemical version of the same mistake. If
the brain has discrete emotion centres, the thinking went, perhaps it
also has discrete emotion \emph{chemicals} --- and a popular scheme
proposes exactly this, casting the brain's main neuromodulators as the
primary colours of feeling: dopamine for pleasure and joy, serotonin for
some axis of mood and aversion, noradrenaline for fear and stress, the
whole palette of emotion mixed from three or four molecular pigments. It
is a beautiful idea. It is also far more speculative than its confident
repetition suggests, and it gets the best-studied of the chemicals
flatly wrong.

Dopamine is not the pleasure molecule, though it is named that everywhere
from journalism to dietary advice. The neuroscientist Wolfram Schultz,
recording from dopamine neurons while animals learned to expect rewards,
found that the cells did not fire for pleasure. They fired for
\emph{surprise} --- specifically, for reward that was better than
predicted. A monkey given an unexpected drop of juice showed a dopamine
burst; once the juice became predictable, the burst migrated backward in
time to the cue that forecast it, and the juice itself, now expected,
produced nothing. Dopamine, that is, encodes a \emph{prediction error},
a running signal of how much better or worse the world is turning out
than the brain had bet --- the engine of learning, not the feeling of
joy. The psychologist Kent Berridge sharpened the dissociation further by
separating ``wanting'' from ``liking'': you can knock out an animal's
dopamine and it will still show every sign of \emph{enjoying} a sweet
taste, the pleasure intact, while losing all motivation to pursue it. The
dopamine drives the wanting; the liking lives elsewhere, in small opioid
hotspots scattered through the reward system. Pleasure and the pursuit of
pleasure, fused in the folk picture, come apart in the brain.

The lesson generalises, and it points where this chapter is heading. The
neuromodulators are not the atoms of discrete emotions. They are broad,
diffuse, slow-acting signals that tune the whole system at once ---
adjusting how much the brain learns, how vigorously it pursues, how
strongly it weights threat, how it colours everything currently in
play. They supply something more like the \emph{valence and intensity}
of experience --- its pleasantness or unpleasantness, its urgency ---
than its specific identity as fear or anger or grief. And that raw
dimension of pleasant-to-unpleasant and calm-to-aroused, smeared across
the brain by chemistry rather than localised in any centre, is precisely
the material from which one of the two great theories of emotion proposes
that all the rest is built.

\section{Six faces}

That theory has a formidable opponent, however, and to feel the force of
the disagreement you have to take the other side at its strongest. The
case that emotions are real, discrete, hardwired things --- not cultural
inventions, not constructions, but evolved equipment as universal as the
skeleton --- rests on one of the most influential bodies of work in
twentieth-century psychology.

In the late 1960s the psychologist Paul Ekman set out to test a claim of
Darwin's: that the facial expressions of emotion are universal,
inherited rather than learned, the same in every human population because
they evolved before the cultures diverged. The obvious worry was that
people everywhere recognise a smile because Hollywood and advertising
have taught the whole world the same code. So Ekman went to the Fore
people of the Papua New Guinea highlands, who at the time had had almost
no contact with Western media. He told them small stories --- a person
whose child has died, a person about to step on a pig, a person meeting
a friend --- and asked them to match each to a photograph of a Western
face. The Fore sorted the faces much as Westerners did. Sadness, anger,
fear, disgust, surprise, happiness: the expressions read across the
cultural gulf. From this work came the doctrine of the basic emotions ---
a small set of evolved, pan-cultural states, each with its own facial
signature, its own physiological profile, its own dedicated circuitry ---
and the meticulous Facial Action Coding System that Ekman built to
catalogue every muscle movement the face can make.

The neuroscientific complement to Ekman's psychology came from Jaak
Panksepp, who spent his career stimulating and recording from the deep
subcortical structures of other mammals and concluded that evolution had
furnished all of us, human and animal alike, with a handful of primal
emotional systems wired into the ancient core of the brain. He named
seven, in the small-capital convention he favoured to mark them as
biological primitives rather than English words: \textsc{seeking},
the engine of curiosity and pursuit; \textsc{fear}; \textsc{rage};
\textsc{lust}; \textsc{care}, the circuitry of nurturance;
\textsc{panic/grief}, the anguish of separation; and \textsc{play}, the
source of social joy. These were not, in Panksepp's view, human
constructions laid over a neutral substrate. They were the substrate ---
shared with rats and dogs, switchable on and off with an electrode,
older than cortex, older than thought. A rat separated from its mother
emits distress cries from the same \textsc{panic/grief} system that
aches in a grieving human; the continuity is the point.

\aside{The basic-emotion view has the great advantage of matching how
emotion feels. Fear and disgust really do seem like different things,
not different labels on one underlying broth --- they arrive with
different bodily signatures, different facial expressions, different
action urges, and they appear early in infancy and across the animals.
Any theory that says the discreteness is an illusion has to explain away
a great deal of apparent natural-kind structure. This is the camp's
enduring strength, and it is why the argument has not simply been
conceded.}

\section{The emotions you construct}

The opposing camp begins by looking for the basic emotions in the brain
and failing to find them. If fear is a discrete evolved module, there
should be a fear circuit --- a patch of tissue that activates for fear
and not for anger, a neural signature as specific as the feeling. The
psychologist Lisa Feldman Barrett and her colleagues went looking, in a
large meta-analysis pooling the brain-imaging evidence across hundreds of
studies, and what they found was not a set of dedicated centres. The
amygdala, the supposed fear centre, lit up for novelty and surprise and
salience of many kinds; the insula, the supposed disgust centre, lit up
for almost every feeling there is. No discrete emotion mapped cleanly
onto its own dedicated territory. The same regions kept reappearing
across fear and anger and sadness alike, in shifting, context-dependent
combinations. The brain did not seem to be organised the way the
basic-emotion theory required.

From this absence Barrett built the \emph{theory of constructed emotion},
and its debt to the Perception chapter is direct and deliberate. Recall
the central claim there: that perception is not reception but inference,
the brain's best hypothesis about what in the world produced an ambiguous
signal, assembled from prior experience and corrected by the senses ---
that to see is to run a controlled hallucination tethered by evidence.
Barrett applies exactly this machinery to the body. The brain receives a
stream of interoceptive signals --- the state of the heart, the gut, the
lungs, the diffuse chemical weather of the neuromodulators --- which on
their own are raw and ambiguous, carrying only \emph{core affect}: a
position on the two crude axes of pleasant-to-unpleasant and
calm-to-aroused that the previous section's chemistry supplies. A racing,
unpleasant arousal could be fear, or anger, or excitement, or the early
warning of illness; the bodily signal alone does not say. The brain, on
this account, does to that signal what it does to the ambiguous light on
the retina: it interprets it, predicting its most probable cause using
concepts drawn from a lifetime of experience, and the emotion you feel is
the category it settles on. Fear is not detected. It is constructed --- a
prediction the brain makes about what its churning body means, given the
situation and everything it has learned about situations like this.

If that is right, then the discrete emotions are not the brain's
furniture but its \emph{readings} --- as real as the seen world is real,
and as constructed. It would explain the meta-analytic mess (there is no
fear circuit because fear is not a circuit but a recurring interpretation
assembled from general-purpose parts). It would explain why the same
bodily state can be felt as different emotions depending on what you
believe is happening --- why the trembling before a performance can flip
between dread and exhilaration with a change of mind. And it returns,
with a vengeance, to the universality that Ekman's faces seemed to
settle.

For the cross-cultural evidence has been reopened. When later researchers
returned to remote, small-scale societies and dropped Ekman's method ---
which had handed people a short list of emotion words and asked them to
match --- the universality faded. Presented with the same Western faces
but asked to sort them freely, without the supplied labels, people of the
Himba in Namibia did not group them into the Western basic-emotion
categories. Studying the Trobriand Islanders, another team found that the
wide-eyed, gasping face that Westerners read as fear was read instead as
a threat display --- the face of someone \emph{about to attack}, not
someone afraid. The forced-choice method, the critics argued, had not
discovered universal emotions so much as imposed them, leading witnesses
to a conclusion the looser test would not support. The construction camp
takes this as evidence that the discreteness of emotion is partly a gift
of language and culture --- that the categories are learned, vary across
peoples, and are read into faces rather than off them.

\aside{It would be wrong to leave the impression that the constructionists
have won; they have not, and the basic-emotion theorists return fire
vigorously, disputing the meta-analytic methods, the cross-cultural
re-studies, and the claim that core affect plus concepts can really
generate the full specificity of emotional life. The field is genuinely
split, as split as the memory taxonomists and the memory architects were,
and for a structurally identical reason. What is no longer seriously
disputed is the narrow, deflationary point underneath the fight: there is
no tidy one-to-one map from named emotions to dedicated brain centres.
Whatever emotions are, they are not a row of labelled organs.}

\section{Why the maps disagree}

The field's reference works carve emotion in ways that cannot be
reconciled, and by now the reason should be familiar: they are not
disagreeing about the data but instantiating two different theories of
what emotion fundamentally is.

One family of catalogues lists the emotions as discrete kinds --- fear,
anger, joy, disgust, sadness --- each a distinct entry, the way a
field guide lists species. This is the basic-emotion view written into
taxonomy: emotions are natural kinds, evolved and universal, and the job
of a catalogue is to enumerate them. The other family refuses the list.
It represents emotion as a space --- the valence and arousal of core
affect --- across which particular feelings are transient regions rather
than fixed entries, and it scatters the machinery of feeling across
networks shared with perception, memory, and bodily regulation, granting
no emotion its own dedicated home. This is the constructionist view
written into taxonomy: there is nothing discrete to list, only a
continuous affective ground that concepts carve into emotions after the
fact.

These are the opposite-facing trees of the Memory chapter in yet another
costume. There, one tradition put a single \textsc{Memory} at the root
and made every kind a subtype beneath it, while another refused the root
and hung the components side by side as parts of a process. Here, one
tradition takes the discrete emotions as the primitives and treats the
underlying affect as derived, while the other takes the affect as
primitive and treats the discrete emotions as constructed. As with
memory, the dispute cannot be settled by collecting more data of the same
kind, because it is a disagreement about which level is fundamental ---
about whether ``fear'' names a thing the brain contains or a thing the
brain does. A catalogue that listed the basic emotions cleanly would be
taking Ekman's side; a catalogue that dissolved them into a valence map
would be taking Barrett's; and most catalogues, built by committees over
decades, quietly contain both, contradicting themselves in exactly the
places where the science is unresolved.

Underneath even that lies the boundary this chapter opened on --- the one
between emotion and cognition, feeling and thought. The folk picture, and
much of the philosophical tradition, drew it as a hard border with
hostile powers on each side: the passions below, reason above, the second
forever struggling to master the first. The atlas's evidence erases the
border. Elliot proved that reason without emotion does not soar but
stalls --- that the somatic lean of feeling is part of the apparatus of
choice, not an impurity in it. The constructionists go further and deny
the border exists at all: if an emotion is a prediction the brain makes
about its own body, assembled from learned concepts, then feeling is a
kind of cognition, built by the same predictive machinery that builds the
seen world, and the question of whether a given act is ``emotional'' or
``rational'' may be as ill-formed as asking whether seeing a face is
perception or memory. It is both, because it was never two things.

What survives all the disagreement is the overturning of the opening
intuition. Emotion is not the weather, the involuntary storm that good
sense must wait out. It is nearer to perception --- a reading, possibly a
construction, the brain's interpretation of the state of its own body in
the light of everything it has learned --- and it is so far from being
reason's adversary that reason cannot run without it. The fear that seems
to seize you from outside your control is, on the best current account,
something your own brain is building, in the present, out of a pounding
body and a learned idea of danger, and handing to you so seamlessly that
you take the verdict for a fact. We are, here as everywhere in this
atlas, the last to know how much of what we feel we have made.











ch10.tex

\atlaschapter{Reasoning}{We are the animal that prides itself on thinking. So why are we so reliably bad at it --- and what if the experiments measuring our failures have been asking the wrong question all along?}

Here is a problem. A bat and a ball cost a dollar and ten cents
together. The bat costs a dollar more than the ball. How much does the
ball cost?

An answer has very likely already arrived, unbidden, and it is ten
cents. It feels clean, obvious, barely worth the arithmetic. It is also
wrong. If the ball were ten cents and the bat a dollar more, the bat
would be a dollar ten and the pair would come to a dollar twenty. The
ball costs five cents; the bat, a dollar more, costs a dollar five; the
two together make a dollar ten. The correct answer requires no
mathematics beyond what a child commands. And yet a majority of people,
including a majority of students at the most selective universities in
the world, give the wrong one --- not because they cannot do the sum,
but because they never do it. An answer presented itself with such
confidence that the question seemed already settled, and the slower work
of checking it was never triggered.

That small, slightly humbling experience --- the wrong answer arriving
effortlessly and \emph{feeling right}, the correct one available only to
whoever stops to do the work --- is the doorway into the science of
reasoning, and into a controversy that has run for half a century over a
question that ought to be simple: are human beings rational? The evidence
that we are not is voluminous, ingenious, and at first overwhelming. The
case that the evidence has been rigged --- that the experiments
manufacture the irrationality they claim to discover --- is more
interesting than it sounds, and it is far from settled. This chapter is
about the cleverest demonstrations of human stupidity ever devised, and
about the serious possibility that they have been measuring the wrong
thing.

\section{The answer that feels right}

The bat-and-ball problem belongs to a tiny quiz called the Cognitive
Reflection Test, devised by the psychologist Shane Frederick, whose three
short questions share a single structure: each dangles an intuitive
answer that is wrong, and each can be solved correctly only by noticing
the pull of that answer and overriding it. The test predicts, with
surprising power, all sorts of other things about how a person reasons
--- not because its questions are hard, but because they measure a
disposition. Some people feel the wrong answer arrive and stop to check.
Most people feel it arrive and go with it.

What the test makes vivid is that thinking is not one thing happening at
one speed. There is the fast process that produced ``ten cents'' --- it
ran automatically, required no sense of effort, and delivered its verdict
with unearned confidence. And there is the slow process that produces
``five cents'' --- it has to be summoned, it feels like work, and it can
override the first only if it is engaged at all. You can sense the two of
them in the gap between the answer you wanted to give and the answer that
survived a moment's scrutiny. The whole modern science of reasoning grew
up around the relationship between those two, and around the discovery
that the fast one runs the show far more than we like to admit.

\section{Two systems, or one?}

The framework that organised this was dual-process theory, and though it
reached the public through Daniel Kahneman's account of a fast,
intuitive ``System 1'' and a slow, deliberate ``System 2,'' the idea was
built over decades by many hands --- Peter Wason and Jonathan Evans, Seymour
Epstein, Steven Sloman, Keith Stanovich among them. The two systems, in
the standard telling, have opposite characters. The first is fast,
automatic, parallel, effortless, and emotionally flavoured; it is the
source of snap judgments, gut feelings, the answer that feels right. The
second is slow, serial, effortful, and rule-governed; it is what you use
to check a sum, follow an argument, or resist a tempting conclusion. The
first runs constantly and almost for free. The second is lazy, expensive
to run, and easily left in its chair.

The framework's reach extends well past arithmetic, and nowhere more
strikingly than into moral judgment, where the psychologist Joshua Greene
turned it into a theory of why our ethical intuitions contradict
themselves. Consider the famous dilemma in two forms. A runaway trolley
will kill five people unless diverted; you can throw a switch to send it
down a side track where it will kill one. Most people say throw the
switch. Now the same arithmetic in a different staging: the only way to
stop the trolley is to push a large stranger off a footbridge into its
path, killing him to save the five. Most people recoil --- \emph{no}. The
numbers are identical; the intuition reverses. Greene's brain imaging
suggested why: the up-close, hands-on violence of the footbridge case
engages the emotional machinery, and that fast affective alarm produces
the refusal, while the impersonal switch case leaves the emotional system
quiet and lets cooler calculation reach the utilitarian answer. The two
systems, on this account, are not just fast and slow. They can deliver
\emph{opposite verdicts} on the same problem, and which one wins shapes
what we call right and wrong.

\aside{The Emotion chapter's patient returns here with new force. Greene's
account predicts that people whose emotional alarm is silenced should
find the footbridge push easier to endorse --- and patients with damage
to the ventromedial prefrontal cortex, the region that left Elliot unable
to feel his way to a decision, do exactly that. They judge the personal,
violent sacrifice more coldly acceptable than the rest of us can bear to.
The somatic marker that the Emotion chapter argued was necessary for
ordinary choice turns out to be doing moral work too: the flinch is part
of the judgment. Remove it and the arithmetic wins.}

All of which makes dual-process theory sound more settled than it is.
Its critics, and there are many, point out that the cartoon version ---
two little agents in the head, a hot one and a cool one, taking turns at
the controls --- explains everything and therefore predicts nothing.
Almost any result can be told as System 1 overruled by System 2 or System
2 captured by System 1, after the fact. The theory's own leading
defenders, Evans and Stanovich among them, have spent years trying to
rescue it from its popularity, insisting that there are not two systems
but two \emph{types} of processing --- one fast and autonomous, one slow
and demanding of working memory --- and that the neat personification is
a teaching metaphor, not a claim about the brain's architecture. The
framework everyone reaches for to describe reasoning is, like the
two-streams models and the multicomponent memory before it, a contested
map rather than a settled finding. It is useful. Whether it is true, in
the strong sense, is exactly the kind of question the rest of this chapter
keeps running into.

\section{The selection task}

The single most studied demonstration in the history of reasoning is also
one of the simplest, and like the bat-and-ball it is best met by failing
it yourself. Four cards lie on a table. Each has a letter on one side and
a number on the other. The faces you can see read: \texttt{E},
\texttt{K}, \texttt{4}, \texttt{7}. Here is a rule about these four cards:
\emph{if a card has a vowel on one side, it has an even number on the
other.} Which cards must you turn over --- the fewest that will settle it
--- to find out whether the rule holds?

Most people turn the \texttt{E}, which is right, and then turn the
\texttt{4}, which is useless. Turning the \texttt{4} can only confirm the
rule, never break it: behind an even number the rule permits anything at
all. The card that can break the rule, and which almost nobody turns, is
the \texttt{7} --- because if a vowel hides behind that odd number, the
rule is false. The logically decisive cards are the \texttt{E} and the
\texttt{7}, and when the psychologist Peter Wason first ran this in the
1960s, only about one person in ten found them. The rest reached,
reliably, for the card that could only flatter the rule and ignored the
one that could refute it. We are drawn to seek confirmation and we shy
from seeking falsification --- the precise opposite of the move that, as
the philosopher Karl Popper had argued, actually tests a claim.

It looks like proof that humans are simply bad at logic. And then comes
the twist that made the selection task immortal. Keep the logic exactly
the same, but change what the cards are about. Now they show drinkers and
their ages, and the rule is a law: \emph{if you are drinking alcohol, you
must be over eighteen.} One card says \textsc{beer}, one says
\textsc{cola}, one says \textsc{twenty-five}, one says \textsc{sixteen}.
Which do you check to catch a violation? Suddenly almost everyone is
correct --- you check the beer-drinker and you check the sixteen-year-old,
and you do not bother with the cola or the twenty-five-year-old. The
underlying structure is identical to the vowels and even numbers, a piece
of logic that nine in ten people had just failed. Dressed as a rule about
catching rule-breakers, it becomes so easy that it feels like common
sense. The same abstract form is nearly impossible in one costume and
nearly trivial in another. Whatever reasoning is, it is not a single
content-blind logic engine that grinds through any problem of the right
shape. The content changes everything.

\aside{Why should a rule about drinkers be so much easier than a rule
about vowels? One bold answer became, for a while, the standard one ---
and then a cautionary tale. The evolutionary psychologists Leda Cosmides
and John Tooby proposed that natural selection had built into us a
dedicated \emph{cheater-detection module}: a special-purpose mental organ
for policing social contracts, for spotting the person who takes a
benefit without paying the cost, because reciprocal cooperation could
never have evolved without one. We fail the abstract task and ace the
drinking-age one, on this view, because the second trips a circuit built
to catch cheats. It is a clean and powerful story. It is also, like the
``neurons that shaped civilisation'' and \emph{the} gene \emph{for}
language, almost certainly too clean. The performance boost shows up for
rules that involve no cheating at all --- permissions and obligations
generally, not specifically the detection of freeloaders --- and rival
accounts explain the effect with general-purpose machinery for reasoning
about what is permitted and required, no dedicated cheater module needed.
The real finding survives: the content and social framing of a problem
transform our ability to reason about it. The module built to explain it
was a metaphor that outran its evidence.}

\section{Are we irrational?}

The selection task is about logic. The deeper indictment of human
rationality came from the study of judgment under uncertainty, and from
the partnership of two Israeli psychologists, Amos Tversky and Daniel
Kahneman, whose work in the 1970s mapped the systematic ways that
intuition departs from the laws of probability. Their claim was not that
people make random errors. It was that we reason by a small set of mental
shortcuts --- heuristics --- that work well enough most of the time and
fail in specific, predictable directions, and that the failures reveal
the shortcuts.

One shortcut is \emph{representativeness}: we judge how likely something
is by how well it matches our mental prototype, ignoring the cold base
rates that ought to govern the estimate. Told about a person described in
terms that fit the stereotype of a librarian, people judge them more
likely to be a librarian than a salesperson, untroubled by the fact that
salespeople vastly outnumber librarians. The same heuristic produces the
notorious conjunction fallacy: given a vivid sketch of a woman who in her
youth was outspoken about social justice, people rate it \emph{more}
probable that she is now a bank teller \emph{and} a feminist than that
she is a bank teller --- which cannot be, since every feminist bank
teller is already a bank teller, and the conjunction can never be likelier
than one of its parts. The detailed story is more representative, so it
feels more probable, and feeling overrides the logic. Other shortcuts
followed: \emph{availability}, judging frequency by how easily examples
come to mind, so that vivid risks loom larger than common ones; and
\emph{anchoring}, in which an arbitrary number thrown out before an
estimate drags the estimate toward it, even when everyone knows the
number was arbitrary. The accumulated catalogue of biases was vast, and
its message was bleak: the human mind is not a faulty calculator but a
systematically miscalibrated one, running on shortcuts that misfire in
knowable ways.

This was the consensus for a generation, and then it met a determined
opponent. The psychologist Gerd Gigerenzer mounted an argument that the
heuristics-and-biases program had stacked the deck. The ``errors,'' he
contended, were largely artefacts of how the problems were posed ---
specifically, of posing them in terms of single-event probabilities, a
format that is recent, unnatural, and not the format in which evolved
minds, or working physicians, actually encounter the world. Take the
most damning result of all: the medical-test problem, in which doctors
asked to interpret a positive result on a screening test for a rare
disease wildly overestimate the chance the patient is actually sick,
neglecting the base rate so badly that the clinical implications are
frightening. Gigerenzer and his colleagues showed that when the very same
problem is rephrased in \emph{natural frequencies} --- not ``the disease
affects one per cent of the population and the test has such-and-such a
false-positive rate'' but ``ten out of every thousand people have it; of
the rest, so many test positive anyway'' --- the fog lifts. Performance,
including among physicians, improves dramatically. The Bayesian reasoning
the doctors supposedly lacked was there all along; it had been hidden by
a representation the mind is not built for. On this telling the biases are
not flaws in human reasoning but mismatches between an ecological
intelligence and a laboratory format, and the fast-and-frugal heuristics
that Tversky and Kahneman cast as sources of error are, in their natural
habitat, remarkably accurate --- smart, not stupid.

\aside{The dispute between these two camps was real and sometimes sharp,
and it is worth being clear about what it did and did not settle. The
heuristics-and-biases tradition was not refuted; many of its effects are
robust and survive reframing. But the strong story --- humans as
incorrigibly irrational, error woven into the fabric of cognition --- did
not survive intact. The same data that looked like proof of stupidity in
one representation looked like competence in another, which means the
data alone could not decide whether we are rational. What you conclude
depends on what you think reasoning is \emph{for}, and against what
standard its performance should be judged. That question, lurking under
the whole debate, is the one the chapter turns to next.}

\section{The Bayesian turn}

If neither pure logic nor classical probability is the standard the mind
was built to meet, what is? A growing body of work answers: the mind is
not a broken logician but a rather good \emph{probabilist}, an engine for
drawing reasonable inferences from uncertain evidence --- and many of its
apparent errors dissolve once you measure it against that standard rather
than the textbook one.

This is the same idea the Perception and Emotion chapters arrived at from
their own directions. There, the brain was cast as a prediction machine:
perception as the brain's best probabilistic guess about what produced an
ambiguous signal, emotion as its best guess about what a churning body
means. The ``new paradigm'' in the psychology of reasoning, associated
with Mike Oaksford and Nick Chater, extends the same logic to thought
itself. People reason about the everyday world not with the rigid
certainties of formal logic --- where one counterexample destroys a rule
--- but with the graded confidence of probability, where evidence makes
conclusions more or less likely and rules hold \emph{usually}, not
always. Seen this way, even the Wason selection task looks different. A
reasoner trying not to verify a certain truth but to gather the most
\emph{informative} evidence about an uncertain rule --- to reduce
uncertainty efficiently, the way a good experimenter chooses what to
measure --- will rationally favour exactly the cards most people choose.
The behaviour that classical logic scored as a failure can be scored, by
the standard of optimal information-gathering, as something close to
rational. The participants were not flunking logic. They were, perhaps,
solving a different and more sensible problem than the one the
experimenter thought they had set.

The Bayesian turn does not vindicate every intuition --- people are still
manifestly capable of real and costly errors --- but it reframes the
central question. The issue may never have been whether humans live up to
the rules of logic, a system invented late and deliberately precisely
because it does \emph{not} come naturally. The issue is whether we make
good use of uncertain evidence, and by that measure the mind looks far
more competent, and its shortcuts far more like sound engineering, than
the bleakest reading of the bias literature allowed.

\section{What reasoning is for}

There is a deeper possibility still, and it inverts the entire enterprise.
Every framework so far has assumed that reasoning is for getting to the
truth --- that a lone mind reasons in order to arrive at correct
conclusions, and that its biases are bugs in that solitary truth-seeking.
What if that assumption is the error? What if reasoning did not evolve to
help an individual think correctly, but to help an individual argue ---
to produce justifications that persuade others and to evaluate the
justifications others offer?

This is the argumentative theory of reasoning, proposed by the cognitive
scientists Hugo Mercier and Dan Sperber, and its explanatory payoff is
that it turns the most embarrassing fact about human reasoning into a
feature. Confirmation bias --- our relentless tendency to seek evidence
for what we already believe and to neglect evidence against it --- is a
catastrophe if reasoning is for finding truth. But it is exactly what you
would build if reasoning is for \emph{winning arguments}: a device
optimised to generate reasons for your own side is supposed to be
one-sided. And the same theory explains why we are so much better at
spotting the flaws in other people's arguments than in our own, and why
reasoning that fails in isolation so often succeeds in groups --- where
each person's biased production of arguments is checked by everyone
else's biased evaluation, and a good answer can emerge from a roomful of
partial ones. Reason, on this account, is not the solitary truth-detector
of the philosophical tradition. It is a social instrument, built for the
group, that works poorly alone precisely because it was never meant to be
used alone.

\aside{This connects reasoning back to the most social chapter in the
atlas. The capacity to model other minds --- to represent what someone
else believes and why --- is the substrate an argumentative reason would
require, and it is no accident that the same species which reads other
minds continuously is the one that reasons by producing and weighing
arguments. The moral psychologist Jonathan Haidt reached a parallel
conclusion from the study of ethics: that our moral judgments are made
first, fast, by intuition, and that the reasoning we offer for them
arrives afterward, as justification --- the rider explaining decisions the
elephant has already made. In Haidt's image the conscious reasoner is not
the president of the mind but its press secretary, tasked with defending
choices it did not make. Mercier, Sperber, and Haidt are describing the
same animal from different angles: one whose reasoning is built less to
discover what is true than to defend what it already holds and to
contest what others hold.}

\section{Why the maps disagree}

The field's reference works carve reasoning along fault lines that, by
now, have a recognisable shape --- and the deepest of them concerns where
inference even belongs in the order of cognition.

The psychometric tradition, descended from the measurement of
intelligence, treats induction and deduction as the very summit of the
mind: the core of \emph{fluid reasoning}, the most complex and most
general cognitive capacity there is, the thing that the hardest items on
an intelligence test are built to tax. To reason --- to infer a rule from
instances, to draw a conclusion from premises --- is, on this map, the
crowning achievement, the apex ability beneath which the simpler
faculties are arranged. But the educational tradition draws the opposite
picture. In the taxonomy of learning objectives that organises much of
how teaching is structured, \emph{inferring} sits near the bottom ---
filed as a humble sub-skill of mere comprehension, the second of six
rising levels, something a student does on the way to the harder work of
analysis, evaluation, and creation. The identical operation occupies the
penthouse of one hierarchy and the basement of another. One tradition
says inference is the most sophisticated thing the mind does; the other
says it is among the most basic.

Both are right, which is why neither can be reconciled into the other.
They are measuring different things and calling them by the same name.
The psychometrician means the capacity that distinguishes a brilliant
reasoner from an ordinary one --- the rare power to crack a hard novel
problem --- and that genuinely is at the ceiling of cognition. The
educator means the everyday act of reading a little more out of a text
than it literally states, the small inference any competent
comprehender makes constantly and without strain --- and that genuinely
is near the floor. The same word, ``inference,'' names both the
mountain's peak and the ground it stands on, and a catalogue that filed
it cleanly in one place would be losing half its meaning. The maps
disagree because reasoning runs the whole vertical range of cognition,
from the automatic inference that completes a sentence to the deliberate
inference that proves a theorem, and no single shelf holds both.

The same fracture runs through the chapter's larger quarrel. Is reasoning
one thing or many? A single content-blind logic engine, as the selection
task seemed at first to assume --- or a toolbox of content-specific
competences, devastating about drinkers and helpless about vowels? A
faculty for tracking truth, or an instrument for winning arguments? A
faulty approximation to logic, or a sound approximation to probability?
The reference atlases that list reasoning as a unitary ability and those
that fracture it into a dozen specific skills are not making clerical
choices; they are taking sides, mostly silently, on whether the thing has
a nature at all or is just the name we give to whatever the mind does when
a problem cannot be solved by reflex. Even the company a catalogue keeps
betrays its commitments: some file mathematical \emph{knowledge} --- what
the symbol $\pi$ means, what a theorem states --- right beside reasoning
\emph{ability}, the declarative \emph{what} shelved next to the procedural
\emph{how}, because the statistics of who is good at each place them
together, never mind that knowing a fact and drawing an inference are
different acts. Others enter \emph{intuition} --- understanding that
arrives without any conscious reasoning at all --- as reasoning's shadow
and opposite, the thing it is defined against.

What the bat and ball taught at the outset holds all the way down. The
mind runs two kinds of process, a fast one that delivers answers and a
slow one that checks them, and most of waking life is governed by the
first. The science of reasoning has spent fifty years documenting the
ways the fast process misleads us, and then a second wave arguing that
those misleadings are, in the mind's natural habitat and against the
mind's proper standard, mostly not errors at all. Both waves are part of
the truth. We are not the flawless reasoners the Enlightenment hoped for,
nor the walking bundle of fallacies the bias literature seemed to find.
We are something stranger and more interesting: a creature whose thinking
was built for speed, for survival, and for argument, on which the slow,
deliberate, truth-seeking reason we most admire sits lightly, expensively,
and late --- summoned only when something stops us, the way the bat and
ball can stop us, long enough to do the work.










ch11.tex

\atlaschapter{Executive Function}{Something in you keeps a goal alive against distraction, holds back the easy wrong move, and switches tasks when the situation turns. We call it the executive --- and the hardest problem in the science is that the word smuggles in a little person who does not exist.}

The previous chapters have each described a piece of mental machinery ---
a memory store, an attentional filter, a perceptual inference, a reasoning
process --- as though it ran on its own. But none of them runs on its own.
Something coordinates them. Something decides that right now you will
attend to this and ignore that, hold this in mind and let that go, follow
the rule you set yourself a moment ago rather than the louder pull of
habit. When you resist checking your phone to finish a paragraph, when
you stop yourself from saying the cutting thing, when you keep a goal in
view across an afternoon of interruptions, you are exercising the faculty
this chapter is about. Psychologists call it executive function, or
cognitive control, and it is the nearest thing the mind has to a manager.

It is also the place where the science of cognition meets its oldest and
most dangerous temptation. The moment you describe a part of the mind as
the manager --- the executive, the controller, the central command ---
you have quietly installed a little person inside the head to do the
managing, and a little person inside the head explains nothing, because
now you need a smaller person inside \emph{them}. The whole history of
this field is, in a sense, the history of trying to describe control
without a controller: to explain how a brain keeps itself on task without
positing an inner agent who keeps it on task. This chapter is about the
faculty that holds the rest of the mind together, about the patients who
have lost it, and about the long campaign to evict the homunculus that the
very word ``executive'' keeps trying to let back in.

\section{The iron and the legend}

The science of the frontal lobes is usually said to begin on an afternoon
in 1848, on a railway construction site in Vermont, with a man named
Phineas Gage and a tamping iron --- a metre-long, inch-thick rod for
packing blasting powder into rock. A premature explosion drove the iron
up through Gage's left cheek, behind his eye, through the front of his
brain, and out the top of his skull, landing many yards away. Gage did
not die. He did not even lose consciousness for long. Within minutes he
was sitting up and talking, and within months the wound had healed and
he was walking around a living, conversing man with a hole blown through
his frontal lobe.

What made Gage famous was not that he survived but what was said to have
happened to him afterward. According to his physician, John Harlow,
writing years later, the efficient and capable foreman his employers had
valued was gone, replaced by someone profane, capricious, unreliable,
unable to settle on a plan or stick to it --- a man whose friends said he
was ``no longer Gage.'' The intellect was intact; the person had changed.
For a century this story did exactly the work the young science needed: it
located in the frontal lobes the seat of character, planning, and
self-government, the higher faculties that make a person a functioning
agent rather than a bundle of impulses.

The trouble is that the story is mostly legend. The historian Malcolm
Macmillan, who tracked down what little hard evidence survives, found the
documentation of Gage's supposed transformation to be thin, late, and
contradictory --- a few sentences, written long after the fact, that
later authors inflated into the lurid case study every textbook repeats.
The most telling fact cuts the other way: in his later years Gage worked
for a stagecoach company in Chile, driving teams of horses over long
routes, a job demanding planning, route-finding, social negotiation, and
sustained self-control --- precisely the executive capacities he was
supposed to have lost. He seems, in other words, to have substantially
recovered, and the famous wreckage of his personality may owe more to the
needs of the science than to the man.

\aside{Gage is the founding case of this chapter and also its first
warning. The frontal lobes do govern something real; damage to them does
disrupt the management of behaviour. But the field's eagerness to read a
clean lesson off a single dramatic patient --- the eagerness this atlas
has watched produce ``the neurons that shaped civilisation'' and \emph{the}
gene \emph{for} language --- was there from the very beginning, in the
story of the man with the iron through his head. The executive is the most
seductive thing in cognition to mythologise, because we so badly want
there to be someone in charge.}

\section{The patient who passes every test}

If Gage's deficit was overdrawn, the genuine syndrome it gestured at is
real, and it has a strange and revealing signature: the patients who
suffer it can pass nearly every test you give them and still cannot run
their lives.

A person with damage to the frontal lobes may retain a high IQ, an intact
memory, fluent language, and normal perception. Tested in the clinic,
they perform within normal limits on task after task --- and then they go
home and cannot hold a job, cannot plan a meal, cannot start a task and
see it through, cannot stop themselves from following whatever impulse or
distraction the moment presents. For a long time this was a paradox that
made clinicians doubt the patients, because the formal tests came back
clean. The resolution was an insight into what the tests were missing. A
test supplies its own structure: it tells you when to start, what to do,
when to stop, and it sets the goal for you. What frontal patients have
lost is precisely the capacity to supply that structure \emph{themselves}
--- to generate a goal, hold it against competing pulls, and organise
behaviour around it over time. The clinic, by handing them the structure,
was concealing the very deficit it was meant to find. Out in a world that
hands you nothing, the deficit is catastrophic.

The most uncanny demonstration of what control is comes from watching its
total collapse. The neurologist François Lhermitte described patients with
frontal damage who exhibited what he called \emph{utilisation behaviour}:
placed in front of an object, they would compulsively use it, whether or
not it made any sense to. Set a pair of glasses before such a patient and
they put them on; set down a second pair and they put those on too, now
wearing two. Place a glass and a pitcher of water and they pour and drink,
again and again, because the objects are there. Lhermitte could lead a
patient around his apartment and find them spontaneously getting into the
bed at the sight of it, in the middle of the day, fully dressed, because a
bed affords lying down and the affordance had become a command. The
environment was driving the behaviour directly, as if the patient had
become a passive conduit from stimulus to action.

What this reveals, by its absence, is the real work of executive function:
the capacity \emph{not} to respond. A healthy brain is constantly
suppressing the obvious, prepotent, environmentally suggested action in
favour of whatever its current goal requires. You see a bed and do not
lie down, because you have somewhere to be. You see the word and name the
colour instead, because that was the instruction. Strip out the
suppression, the holding of an internal goal against the pull of the
external world, and behaviour decays into a chain of reflexes triggered by
whatever happens to be in front of you. Executive function is, at its
core, the freedom from the immediate --- the ability to act from a goal
you are holding rather than a cue you are facing.

\section{The ghost in the machine}

So far this is description, and the description keeps reaching for the same
kind of word: control, suppression, holding, governing. Each one implies a
governor. And here the science walks up to the edge of its deepest
conceptual hazard. If executive function is the work of a central
executive --- a controller that decides what to attend to, what to
inhibit, what to hold in mind --- then we have explained the intelligent
coordination of the mind by positing an intelligent coordinator inside it,
and that is no explanation at all. For the controller would itself need to
attend, decide, and remember, which means it would need its own executive,
which would need its own, and so on without end. The homunculus --- the
little person in the head --- does not solve the problem of control. It
only hides it inside a smaller skull.

The earlier chapters already felt this hazard without naming it. The
Working Memory chapter described Baddeley's central executive as the most
important component of his model and the one he understood least --- ``in
some ways a placeholder,'' a label for the coordinating work that could
not be assigned to the specialised stores. The Attention chapter found a
central executive lurking again, an attentional supervisor deciding what
gets processed. In both cases the executive was where the explanation ran
out: the box you draw around the part you cannot yet explain, and into
which, if you are not careful, a little person quietly moves.

The decisive move against the homunculus came from a model proposed by
the neuroscientists Earl Miller and Jonathan Cohen, and its elegance is
that it gets control without a controller. The prefrontal cortex, on their
account, does not house an agent who issues commands. It maintains
\emph{representations of goals} --- sustained patterns of activity that
encode what you are currently trying to do, the rule currently in force ---
and these representations do their work not by commanding anything but by
\emph{biasing} the rest of the brain. Recall the Attention chapter's
picture of perception as a biased competition, in which many neural
representations contend and attention tips the scales. Control, in the
Miller--Cohen model, is the same mechanism turned to the service of goals:
the prefrontal goal-representation feeds a bias into the competing
pathways, tilting the contest toward whatever serves the current aim.

The Stroop task is the worked example, and you can feel it operate. Shown
the word \textsc{red} printed in blue ink and asked to name the ink
colour, you must say ``blue'' --- and it is hard, because reading is so
practised that the word-naming pathway fires fast and strong, shouting
``red'' before you can stop it. To succeed, the prefrontal cortex must
hold the goal --- \emph{name the ink, not the word} --- actively and
continuously, and that held goal biases the weak colour-naming pathway
enough to let it win against the dominant reading habit. There is no
homunculus reading the word and deciding to override it. There is only a
goal, maintained as a pattern of activity, lending its weight to the
pathway that would otherwise lose. The controller dissolves into a
representation held against distraction, and the mystery of control
becomes the more tractable mystery of how a goal is kept alive and made to
bias a competition. The little person is gone. What remains is a goal that
does its work by tilting the odds.

\section{One thing, or many?}

With the homunculus evicted, a different question comes into focus, and it
is the one the field has argued over most productively: is executive
function a single capacity or a collection of separate ones? Intuition
pulls both ways. On one hand, ``self-control'' or ``willpower'' feels like
one thing, a single resource you have more or less of, deployed equally
against the snooze button and the second slice of cake. On the other hand,
the specific operations look distinct: suppressing a prepotent response
(inhibition) seems a different act from holding and revising the contents
of mind (updating), which seems different again from switching between
tasks or rules (shifting). Is the executive one muscle or a toolbox?

The most influential answer came from a study by Akira Miyake, Naomi
Friedman, and their colleagues, who used a statistical technique --- 
latent-variable analysis --- to ask whether the common control tasks
measure one underlying ability or several. They found, in a phrase that
became the title of the work, \emph{unity and diversity}. The three
core components --- inhibition, updating, and shifting --- were
\emph{separable}: a person's skill at one did not perfectly predict their
skill at the others, and they could be teased apart as distinct factors.
But they were also \emph{correlated}: the three leaned on something shared,
a common factor that ran through all of them. Executive function was
neither one thing nor three independent things. It was a family --- 
distinct members held together by a common inheritance. The lumpers who
called it a single faculty and the splitters who called it a list of
separate abilities were each half right, and each half wrong, because the
construct genuinely has both structures at once.

\aside{Later work sharpened the picture in a way that is humbling for the
component that intuition most wanted to be central. When Friedman and
Miyake refined the analysis, the common factor and the updating-specific
and shifting-specific factors held up --- but the \emph{inhibition}
factor dissolved. Once the shared executive variance was accounted for,
there was nothing distinctly inhibitory left over: the ability to suppress
a prepotent response was apparently just the common executive factor doing
its work, not a faculty of its own. This is a striking result, because
inhibition --- the capacity \emph{not} to respond that the utilisation-
behaviour patients had so dramatically lost --- had seemed the very heart
of control. It may instead be the most visible face of the common factor
rather than a separate organ. What looks, from the outside, like the
single most executive of acts may not be a distinct thing at all.}

The toolkit that probes these components is by now a century in the
making: the Stroop task and the stop-signal for inhibition; the n-back and
similar tasks for updating, in which you must continuously hold and
refresh the last few items of a stream; the Wisconsin Card Sorting Test
for shifting, in which the sorting rule changes without warning and
frontal patients keep doggedly applying the old one --- perseverating, in
the term of art, unable to release a rule that has stopped working. Each
task isolates, imperfectly, one strand of the executive. None isolates it
cleanly, because the strands are braided. That impurity is not a flaw in
the tasks. It is the unity-and-diversity structure showing through: you
cannot find a pure measure of one component because the common factor is
present in all of them.

\section{The marshmallow}

No strand of executive function has been more celebrated, or more oversold,
than the inhibitory one --- the capacity to delay gratification, to forgo
a smaller reward now for a larger one later --- and its story is the
chapter's second cautionary tale.

In a series of studies beginning in the late 1960s, the psychologist
Walter Mischel sat preschool children alone in a room with a marshmallow
(or a cookie, or a pretzel) and a deal: eat this one now, or wait until I
come back and have two. The children's struggles --- covering their eyes,
turning away, singing to themselves, and often cracking within seconds ---
became one of the most charming bodies of footage in psychology. What made
the marshmallow test famous, though, was the follow-up. Years later, the
children who had waited longer were reported to have higher test scores,
better social competence, healthier outcomes across a range of measures.
The conclusion wrote its own headline: that a single early-emerging trait,
the willpower to delay, was a master key to success in life, more
predictive even than intelligence.

It was too good, and it did not hold. When a larger and more diverse
sample was studied with proper statistical controls, in work published in
2018, the predictive power of the marshmallow test shrank dramatically.
Much of what had looked like the long arm of willpower turned out to be
the arm of something else: a child's family background, the security and
resources of their environment, their general cognitive ability. A child
who waits may be one who has learned that promises are reliably kept ---
which is a fact about their world as much as their self-control --- and
who comes from circumstances that independently predict the later
outcomes. Control the background, and the marshmallow's magic mostly
evaporates. Delay of gratification is real, and executive function is
real, but the dream of a single clean trait that determines a life was a
projection, and the replication era retired it.

\aside{The deflation rhymes with the dissolving inhibition factor of the
previous section. Both say the same uncomfortable thing: that the most
intuitively central, most morally resonant piece of executive function ---
the willpower to hold back --- may be less of a free-standing, measurable,
destiny-shaping thing than we wanted. Executive function matters
enormously. But it is entangled, all the way down, with intelligence, with
environment, with circumstance, and the attempt to extract it as a pure
and powerful single variable keeps failing in the same direction. This is
the atlas's standing lesson once more: the clean, separable, master
variable is the thing to distrust.}

\section{The lobes that grow last}

Two facts about the prefrontal cortex frame everything executive function
does, and they are the same fact running forward and backward in time: it
is the last part of the brain to mature, and among the first to decline.

Executive function is largely absent in early childhood and arrives
slowly. An infant of eight months shows its lack in a tiny, reliable
error studied by the developmental psychologist Adele Diamond: hide a toy
under one cloth, let the baby find it a few times, then hide it under a
\emph{second} cloth in plain view, and the baby reaches back to the first
--- unable to hold the new location in mind and inhibit the now-practised
reach, a failure that tracks the immaturity of the prefrontal cortex.
Across childhood the capacity grows; through adolescence it is still
visibly under construction, which is the unglamorous neuroscience behind
the impulsivity, the susceptibility to immediate reward, and the
difficulty with long-horizon planning that mark the teenage years. The
prefrontal cortex is not structurally mature until the mid-twenties --- 
the brain's last region to finish building, and the marshmallow test,
seen in this light, is partly a measure of how far along that slow
construction a given four-year-old happens to be. At the other end of
life the same region leads the decline: executive functions are among the
earliest casualties of normal aging, which is why the older mind, often
richer in knowledge and vocabulary, can struggle more with juggling tasks,
suppressing distraction, and switching flexibly.

The breadth of what these slow-growing lobes govern is easy to
underestimate, because we file executive function as cold cognition --- 
planning, switching, the Stroop task --- and forget that the same control
machinery governs feeling. The Emotion chapter noted in passing that
reappraising an emotion, deliberately reframing a situation to change its
emotional force, recruits the lateral prefrontal cortex to damp down the
amygdala's response. That is executive function operating on emotion: the
held goal --- \emph{see this differently} --- biasing the affective
machinery the way it biases the colour-naming pathway in Stroop. The
control of attention, the control of action, and the control of feeling
are not three separate systems. They are the same prefrontal capacity to
hold a goal against a prepotent pull, applied to three different kinds of
pull. The teenager who cannot resist the phone, the gambler who cannot
leave the table, and the grieving person who cannot stop the spiral of
rumination are all, in part, contending with the limits of one slow-built,
early-fading faculty.

\section{Why the maps disagree}

The field's reference works carve executive function in incompatible ways,
and by now the disagreement is legible as a record of the science rather
than a failure of it.

Some catalogues grant executive function a single grand heading --- a
central executive, one faculty of control --- and arrange the specific
abilities beneath it. Others refuse the unifying heading and list the
components side by side: inhibition, updating, shifting, planning, each a
separate entry, as though the executive were a department of loosely
related clerks rather than a single manager. And a third kind declines to
give executive function its own territory at all, folding it into
attention --- where it appears as attentional control --- or into working
memory, as the central executive that the Working Memory chapter found
sitting at the heart of Baddeley's model. One catalogue's grand faculty is
another's list of parts is a third's mere aspect of attention.

The unity-and-diversity finding is, almost uniquely in this atlas, a
genuine \emph{resolution} of a map disagreement rather than a deepening of
it. The lumpers who give executive function one heading and the splitters
who give it several are not making opposite errors; they are each
capturing one true half of a structure that is really, demonstrably, both
unified and diverse. A catalogue that lists a single executive is right
that the components share a common factor. A catalogue that lists separate
components is right that they are separable. The only catalogue that is
wrong is one that insists on \emph{only} unity or \emph{only} diversity,
and the contradiction that appears when a reference work contains both is
not sloppiness but accuracy: the construct itself contains both.

Underneath that resolvable disagreement, though, sits a harder one that
the chapter has circled from its first paragraph, and it concerns whether
``executive function'' names a real thing at all or merely marks the
boundary of our ignorance. There is a real danger that the executive is a
\emph{residual category} --- defined not positively, by what it is, but
negatively, as whatever the frontal lobes do that we cannot yet explain by
simpler means. Define a faculty as the part of cognition you have not
reduced to mechanism, and it will shrink as the science advances,
dissolving piece by piece into specific, explicable processes --- the
maintenance of a goal, the biasing of a competition, the updating of a
buffer --- until the unitary ``executive'' has gone the way of the
homunculus it always threatened to become. Something like this may already
be happening: the inhibition factor that dissolved into the common factor,
the central executive that the Miller--Cohen model replaced with held
goal-representations. Each advance turns a piece of the mysterious
executive into ordinary, controller-free mechanism, and the box marked
``executive function'' gets a little emptier.

That, in the end, is the chapter's quiet argument, and the resolution of
the thread that began with the iron through Gage's head. The reason the
word ``executive'' is so hard to define is that it was never the name of a
thing in the brain. It was the name of our intuition that there must be
someone in charge in there --- a manager, a self, a little person holding
the controls --- and the whole achievement of the science has been to show
how a brain governs itself with no one at the helm. There is no executive.
There is a goal, held as a pattern of activity against the pull of
distraction and habit, lending its weight to the pathways that serve it
and starving the ones that do not, kept alive by frontal lobes that take
twenty-five years to finish growing and begin to fail before the rest. When
those lobes are damaged, what is lost is not a ruler but a capacity: the
ability to keep meaning to do something long enough to actually do it. The
manager we imagined was the last and most flattering of the mind's
homunculi, and the strangest thing the science has found is that we run
remarkably well without one.












ch12.tex

\atlaschapter{Learning}{Every previous chapter has been about learning without quite saying so. Named at last, it yields two surprises: that we learn not from what happens to us but from what we failed to predict, and that the studying which feels most like learning is mostly the kind that isn't.}

You can ride a bicycle, probably, and you cannot say how. Asked to
explain what you do when the bike begins to tip to the left, you will
likely say you lean right, which is almost the reverse of the truth ---
you steer briefly into the fall, to the left, sliding the wheels back
underneath your centre of mass, a correction your body makes many times
a second and your conscious mind has no access to at all. The skill is
in there, indelibly; a person who has not cycled for thirty years climbs
on and rides. But the knowledge is mute. It cannot be told, only shown.
Something was learned, durably and well, and the learning left no account
of itself that you can read.

Now recall the night before an examination you crammed for. The
material, gone over and over in a single long session, came to feel
fluent, familiar, mastered --- and then, days later, it had largely
evaporated, far more of it gone than you expected from how thoroughly
known it had felt at midnight. The fluency was real and the learning was
not, or not much of it. The sensation of having learned and the fact of
having learned had come apart, and the sensation had lied.

These two ordinary experiences frame the chapter, because between them
they overturn the folk picture of learning as thoroughly as the folk
picture of memory was overturned by the dark room with shelves. Learning
feels like the accumulation of told knowledge through effortful study,
the conscious filling of a vessel. It is mostly neither conscious nor
told nor effortful in the way that feels like effort, and the conditions
under which it feels most successful are frequently the ones under which
it succeeds least. As with every faculty in this atlas, the learner is
the last to know what the learning is doing.

This is also the chapter where the atlas turns and looks at itself.
Almost everything in the preceding pages was secretly about learning. The
engram that Memory located and the reconsolidation that made it editable
were learning, viewed from the side of what is retained. The priors that
Perception sends down the visual hierarchy are learned regularities ---
predictions are memories, the chapter said, which is to say they are the
residue of learning. The dopamine signal that Emotion found coding reward
``better than predicted'' was, though the chapter only hinted at it, a
learning signal. The infants in Language who extracted words from a
two-minute stream of nonsense were learning, by a mechanism so general it
needed no words at all. Learning is the connective tissue of cognition,
and this chapter makes the tissue the subject --- which means its job is
less to introduce new territory than to find the process that has been
running underneath all the others, and to ask the question the field has
fought over for a century: whether that process is one thing or many.

\section{The surprise machine}

Begin with the deepest and least intuitive fact, the one that does for
learning what reconstruction did for memory. We assume we learn what goes
together. Things that co-occur --- the bell and the food, the cue and the
consequence --- get associated, and the more reliably they co-occur, the
stronger the bond. This is the oldest idea in the psychology of learning,
older than psychology, and it is the natural picture: learning is the
registering of contiguity, the mind noting what tends to accompany what.

It is wrong, or at least radically incomplete, and the experiment that
showed it wrong is one of the quiet landmarks of the field. In the late
1960s the psychologist Leon Kamin ran a procedure now called
\emph{blocking}. First he taught an animal that a tone predicts a mild
shock: tone, then shock, until the tone reliably produces fear. Then he
presented the tone \emph{together with} a new cue --- a light --- and
followed the pair with the same shock, many times. The light now
accompanies the shock perfectly; every time the light appears, the shock
follows. By the contiguity picture the animal should learn to fear the
light. It does not. Tested on the light alone, the animal is almost
indifferent to it. The prior learning about the tone has \emph{blocked}
the learning about the light.

Why? Because by the time the light is introduced, the tone already
predicts the shock completely. The shock, when it arrives, is fully
expected --- and an expected outcome teaches nothing. The light is
perfectly correlated with the shock but \emph{redundant}; it carries no
news. What drives learning, blocking reveals, is not contiguity but
\emph{prediction error}: the gap between what was expected and what
occurred. When the world turns out as predicted, the synapses do not
change. Only surprise --- an outcome better or worse than forecast ---
moves the machinery. You learn, precisely and only, to the extent that
you were wrong.

In 1972 Robert Rescorla and Allan Wagner turned this into a model of
almost brutal simplicity, an equation stating that the change in an
association on any trial is proportional to the prediction error on that
trial --- the difference between the outcome that occurred and the sum of
what all present cues already predicted. When the error is large, learning
is fast; as predictions improve, the error shrinks, and learning slows and
stops. The model explained blocking exactly (the tone has already driven
the error to zero, so the light learns nothing), and it explained a dozen
other effects that the contiguity picture could not, and it remains, half
a century later, one of the most successful theories psychology has
produced. The lesson it carries is the same one this atlas keeps
arriving at from every direction: the brain is not a recorder of what
happens but a prediction machine, and learning is what it does when the
predictions fail.

\aside{The Rescorla--Wagner rule has a second life its authors did not
foresee. The error-driven update at its heart --- nudge your prediction
toward reality in proportion to how wrong you were --- is, in all but
notation, the learning rule behind modern artificial intelligence. The
``reinforcement learning'' that taught machines to play games and the
error-correction that trains neural networks are descendants of an
equation written to explain why a rat ignores a redundant light. The
mind's most ancient learning mechanism and the newest engineering turn
out to be running, at the core, the same arithmetic of surprise.}

The most satisfying part of the story is that the prediction error is not
an abstraction. It has a physical signature, and the atlas has already met
it. The Emotion chapter described Wolfram Schultz recording from the
brain's dopamine neurons and finding, to general surprise, that they did
not fire for pleasure. They fired for reward that was \emph{better than
predicted} --- a burst for the unexpected drop of juice, nothing for the
juice once it became expected, and a dip below baseline when an expected
reward failed to come. That firing pattern is the Rescorla--Wagner error
term, rendered in neurons. Dopamine is not the chemical of pleasure but
the chemical of surprise, broadcasting through the brain a running
estimate of how much better or worse the world is turning out than the
brain had bet --- and wherever that signal lands, it adjusts the
predictions that generated it. The behavioural rule from the rat and the
cellular signal from the monkey are the same thing seen at two altitudes.
Learning, at its root, is the correction of a forecast by its error, and
the brain has a dedicated system for computing the error and a dedicated
chemical for shipping it where the correction must happen.

\section{Stamping in}

That root mechanism --- association adjusted by error --- is the heart of
what is called the \emph{associative} tradition, and to feel the force of
the century-long argument that organizes this chapter you have to take
that tradition at its considerable strength. Its founders built it on
animals doing simple things, very carefully measured, and what they
established has never been undone.

Edward Thorndike, working in the 1890s with cats in what he called puzzle
boxes --- crates a cat could escape by pressing a lever or pulling a loop
--- watched the animals learn and concluded that there was no flash of
understanding involved at all. A cat placed in the box struggled, clawed,
and thrashed at random until it happened to hit the release, and then, on
the next trial, struggled and thrashed and hit it slightly sooner. The
escape time fell gradually, trial after trial, in a smooth curve with no
sudden drop --- no moment where the cat ``saw'' the solution. Thorndike's
\emph{law of effect} stated the mechanism: a response followed by a
satisfying consequence is gradually stamped in, strengthened as a
connection between the situation and the action, while responses followed
by nothing fade. No insight, no reasoning, no representation of the
problem --- only the blind strengthening of whatever happened to work.
Learning was selection among behaviours, a Darwinism inside the animal's
own lifetime, and the smooth curve was its fingerprint.

The other founder needs no introduction and is universally
misremembered. Ivan Pavlov's dogs, salivating at the bell that
preceded their food, gave the world its image of conditioning as a kind
of reflex transfer: the response to the food simply slides over onto the
bell, stimulus substituting for stimulus. The modern understanding,
argued most forcefully by Rescorla, is that this picture is almost
entirely wrong about what the animal is doing. The dog is not
mechanically swapping one trigger for another. It is learning a
\emph{predictive relationship} --- that the bell carries information about
the food --- and the proof is that mere pairing is not enough. In
Rescorla's decisive control, if the food is just as likely to arrive
\emph{without} the bell as with it, no conditioning occurs no matter how
often bell and food coincide, because the bell predicts nothing it
contributes. What the animal tracks is not contiguity but
\emph{contingency}: whether the cue genuinely changes the odds of the
outcome. Even Pavlov's dog, the very emblem of mindless association, turns
out to be running a small statistical inference about the structure of its
world.

\aside{It is worth pausing on how far this already bends the associative
tradition toward its rival. Thorndike's cat strengthens what works;
Pavlov's dog, on the modern reading, learns what predicts what. The first
is closer to blind selection, the second to the building of a little
predictive model. The seam that the rest of the chapter runs along ---
between learning as the mechanical adjustment of bonds and learning as the
construction of knowledge about the world --- is already visible inside
the founding experiments of the school that was supposed to have no use
for knowledge at all.}

\section{The promise that swallowed psychology}

Out of the law of effect and the conditioned reflex grew the most
imperial idea in the history of the science of mind, and --- this being
the atlas --- its most instructive overreach. If all learning is the
strengthening of associations by their consequences, then perhaps
\emph{everything} a human being becomes is learned that way: every
preference, every fear, every skill, every trait of what we vainly call
personality, all of it conditioned into a malleable organism by its
history of reward and punishment. This was \emph{behaviourism}, and for
the better part of fifty years it did not merely study the mind; it
proposed to replace the study of the mind with the study of behaviour,
banishing the inner world as unobservable and therefore unscientific.

Its founder, John Watson, made the promise in its most naked form in
1924, in a boast that has echoed ever since: give him a dozen healthy
infants, he wrote, and his own specified world to raise them in, and he
would guarantee to train any one of them to become any type of specialist
he might select --- doctor, lawyer, artist, beggar, thief --- regardless
of the child's talents, tendencies, abilities, or ancestry. The human
being was raw material; conditioning was the whole of the shaping. Watson
had earlier supplied a chilling demonstration of the principle's reach
into the emotions, conditioning an infant known as Little Albert to fear a
white rat by pairing its appearance with a sudden frightening noise, until
the child cried at the sight of the animal he had previously reached for
--- a study whose ethics were as questionable as its controls, and which
later scholarship has called into doubt on nearly every factual point, but
which fixed the behaviourist creed in the public mind: emotion itself was
just conditioning, installable on demand.

B.~F. Skinner built the mature cathedral. With his operant chambers and
his exquisite manipulation of \emph{reinforcement schedules} --- the
discovery that rewarding a behaviour intermittently and unpredictably
makes it far more persistent than rewarding it every time, the principle
that keeps a gambler at a slot machine --- Skinner mapped the laws of
operant learning with a rigour no one has surpassed. And then he
extrapolated them to everything, most fatefully to language, arguing in
his 1957 book \emph{Verbal Behavior} that a child acquires speech as a
rat acquires lever-pressing, by reinforcement of successive
approximations. It was the high-water mark of the associative empire, and
it broke, as the Language chapter recounted, on a single review. Noam
Chomsky's 1959 demolition argued that no amount of reinforcement could
account for a child's ability to produce and understand sentences never
heard before, that the input was too poor and the achievement too
structured, and that behaviourism had mistaken the discardable surface of
behaviour for the generative system beneath it. The review did not just
wound Skinner's book; it helped end the era and inaugurate the cognitive
one.

\aside{Behaviourism is this chapter's entry in a catalogue the atlas has
been keeping --- the genuine discovery inflated by a grand metaphor until
it overshoots and has to be hauled back. The mirror neurons that
``shaped civilisation,'' the gene \emph{for} language, the amygdala as the
seat of fear: each was a real finding wrapped in a claim too clean for the
evidence. Behaviourism is the largest specimen of the type. Conditioning
is real, its laws are real, and intermittent reinforcement really does
grip a gambler --- but the leap from there to ``a human is nothing but its
conditioning history'' was a category error of enormous ambition, and the
correction, when it came, did not restore a little person to the centre of
the mind so much as restore the mind itself as a thing worth studying.}

What survives the collapse is precise and durable. The laws of
conditioning hold; they describe a real and ancient layer of learning;
intermittent reinforcement, shaping, extinction, and the rest are
genuine and useful. What does not survive is the imperial claim that this
layer is all there is --- that there is nothing in the learner but the
bonds its history has stamped in. The evidence against that claim had in
fact been sitting in the literature since before Skinner's book, gathered
by a man working with the same rats in the same mazes, who kept finding
that they knew more than reinforcement could explain.

\section{The map in the head}

Edward Tolman was a behaviourist by training and a heretic by
temperament, and the experiments he ran in the 1930s and 40s are the
cleanest historical counterpunch to the idea that learning is just
reinforced association. His rats ran mazes, like everyone's rats, but
Tolman kept noticing that they behaved as though they possessed something
the theory forbade them: a \emph{map}.

The decisive result is called \emph{latent learning}. Take three groups
of rats in a complex maze. One group is rewarded with food at the goal
every day, and learns to run the maze efficiently, errors falling on the
usual curve. A second group is never rewarded, and wanders, showing little
improvement. The third group is rewarded only after the tenth day --- for
the first ten days they explore the maze with nothing waiting at the end.
By the associative account this third group should begin learning only
when the reward appears, and should then descend the learning curve slowly
from the top, like beginners. They do nothing of the kind. The day after
the reward is introduced, their performance plunges almost to the level of
the rats that had been rewarded all along. In a single rewarded trial they
reveal a knowledge of the maze they must have been building, unrewarded
and invisibly, throughout the ten days of mere wandering. The reward did
not teach them the maze. It only gave them a reason to show what they had
already learned.

This is the distinction, drawn here for the first time and crucial to
everything after, between \emph{learning} and \emph{performance}. The
rats had learned without performing; the reward changed the performance
without adding to the learning. And what they had learned was not a chain
of reinforced turns but, Tolman argued, a \emph{cognitive map} --- an
internal representation of the maze's layout, its spatial structure, the
relations among its places --- which they could then exploit in flexible
ways no stamped-in habit could produce, taking novel shortcuts and
detours when the maze was altered. Learning, on this account, is not the
strengthening of bonds between stimulus and response. It is the
construction of knowledge about the world, built spontaneously, in the
absence of reward, by an animal that explores because exploring is what
brains are for.

\aside{Tolman's cognitive map was a psychological inference for thirty
years before the brain confirmed it almost literally. In 1971 John
O'Keefe, recording from the hippocampus of moving rats, found neurons that
fired only when the animal occupied a particular place in its environment
--- \emph{place cells}, each tuned to a location, together forming a
neural chart of the space. Later work uncovered \emph{grid cells} laying a
regular coordinate mesh over the world. The hippocampus, the very
structure whose loss left the Memory chapter's H.M. unable to form new
declarative memories, turns out to build and store exactly the kind of
spatial map Tolman's rats behaved as though they had. The cognitive map
was not a metaphor. It was a discovery waiting for an electrode, and it
earned a Nobel Prize.}

The cognitive pole has a more dramatic and more contested exhibit as
well. In the 1920s Wolfgang Köhler, studying chimpanzees, set problems
that could not be solved by trial and error --- a banana hung out of
reach, with boxes scattered about the enclosure. His chimps did not
gradually thrash toward the solution in the manner of Thorndike's cats.
One, after a period of apparent inactivity, would suddenly stack the boxes
and climb, or fit two short sticks together to make a long one, arriving at
the answer whole, as if it had been \emph{seen} rather than stumbled upon.
Köhler called it \emph{insight}, and offered it as proof that higher
animals learn by restructuring their understanding of a problem, not by
stamping in successful movements. The cleanness of the demonstration has
been argued over ever since --- the chimps had relevant prior experience,
the ``sudden'' solutions were less sudden on close inspection --- and
insight has proved far harder to pin down than latent learning. But the
basic contrast stands: some learning looks like the slow selection of what
works, and some looks like the abrupt assembly of a solution from an
understanding of how the parts relate, and no single mechanism has ever
comfortably covered both.

\section{Two ways to learn}

By now the chapter has, in effect, split learning the way the Memory
chapter split memory --- into something declarable and something merely
demonstrable. And the parallel is not loose. It is the same split, in the
same tissue, proved by the same patient.

The Memory chapter described Brenda Milner asking H.M. --- the man who
could form no new conscious memories after the loss of his hippocampi ---
to trace a star while watching only its mirror image, a maddening task at
which people improve over days. H.M. improved exactly as normal people do,
his hand acquiring the skill day by day, while he insisted at each session,
sincerely, that he had never attempted the task before. There, the point
was about memory: two memory systems, dissociated. Seen from this chapter,
it is equally a fact about learning. H.M. could \emph{learn} --- durably,
measurably --- while being unable to form any conscious, declarable record
of having learned. The learning that built the skill and the learning that
would have built a memory of acquiring it were different systems, and one
was destroyed while the other ran on untouched.

This is the distinction between \emph{explicit} learning --- conscious,
declarable, the kind you know you are doing and can describe afterward ---
and \emph{implicit} learning, which proceeds without awareness and yields
knowledge you cannot put into words. The cyclist's balance is implicit.
So, the Language chapter showed, is the infant's first grip on the
boundaries between words: Jenny Saffran's eight-month-olds, exposed for two
minutes to an unbroken stream of invented syllables, extracted the hidden
``words'' by tracking which syllables predicted which --- a feat of
statistical learning performed with no instruction, no reward, and no
possibility of conscious awareness in a pre-verbal infant. Implicit
statistical learning of this kind appears to run constantly and
automatically, in infants and adults alike, soaking up the regularities of
whatever stream of experience the senses deliver, building knowledge the
possessor cannot report and often does not suspect.

The two systems can be dissociated as cleanly as the two memory systems,
and by an experiment of real elegance. Barbara Knowlton and her colleagues
used a \emph{probabilistic classification} task --- a kind of weather
prediction game in which cues are associated with outcomes only
imperfectly, so that no single trial reveals the rule and the structure can
be absorbed only gradually, across many trials, below the level of
articulable knowledge. Amnesic patients with hippocampal damage, of the
kind that left H.M. unable to learn facts and events, learned this task at
a normal rate --- their performance climbing trial by trial --- while being
unable to remember the training sessions at all, sometimes denying they had
ever seen the task before. Patients with Parkinson's disease, whose damage
falls instead on the basal ganglia, showed the mirror image: they could
remember the training perfectly well and yet failed to learn the task. One
system, dependent on the hippocampus, learns facts you can declare; another,
dependent on the basal ganglia and its dopamine signal, learns
statistical regularities you cannot. Damage them separately and you knock
out one kind of learning while sparing the other. Learning, like memory
before it, is not one thing. It is a committee of systems, running on
different tissue, obeying different rules, and dissociable by injury ---
and the version of it we are aware of doing is only the part that reaches
awareness.

\section{The molecule}

If learning is the installation of a change, the change must, in the end,
be physical --- something altered in the tissue, the way the Memory chapter's
engram was the physical trace of a remembering. And the most complete
account we have of what that something is comes not from a human brain or
a rat's but from a sea slug, in the long work of Eric Kandel.

Kandel made a deliberate bet that the basic mechanisms of learning would
be ancient and universal, conserved across the animals, and that they
could therefore be cracked in a creature simple enough to study one cell
at a time. He chose \emph{Aplysia}, a marine snail with a nervous system
of some twenty thousand neurons, many of them large enough to see and
individually identifiable from animal to animal --- and a simple defensive
reflex, the withdrawal of its gill when its siphon is touched. That reflex
can learn. Touch the siphon harmlessly again and again and the withdrawal
weakens --- \emph{habituation}, the simplest learning there is, the
fading of a response to a stimulus that proves to mean nothing. Pair the
touch with a shock to the tail and the withdrawal grows brisk and
exaggerated --- \emph{sensitization}. The slug can even be classically
conditioned, Pavlov's dog rendered in a few traceable cells.

Because Kandel could record from the identified neurons on both sides of
the relevant synapses, he could watch learning happen as a physical event.
Habituation, he found, is a weakening of transmission at specific synapses
--- less neurotransmitter released, the connection turned down. Sensitization
is the reverse, a strengthening, driven by a chemical cascade inside the
cell. And here the story reaches the boundary the Memory chapter drew
between the fleeting and the durable. Short-term learning, lasting minutes,
is a temporary change in how much transmitter the existing synapses
release --- the machinery turned up or down. Long-term learning, lasting
days and longer, requires something more: the activation of genes in the
neuron's nucleus, the synthesis of new proteins, and the physical
\emph{growth of new synaptic connections}. The memory of the training is
built, literally, into new structure. This is why, as the Memory chapter
noted, a consolidating memory can be blocked by a drug that halts protein
synthesis: the durable trace is a construction project, and you can stop it
by withholding the materials.

\aside{The complement to Kandel's slug is a phenomenon in the mammalian
hippocampus called \emph{long-term potentiation} --- the finding, by Bliss
and Lømo in 1973, that a burst of activity across a synapse can leave it
strengthened for hours or days, more responsive to the same input
afterward than before. It is widely taken to be a building block of
learning in the brain, and it obeys a rule the psychologist Donald Hebb had
proposed in 1949 from pure reasoning: neurons that fire together wire
together. A cue and an outcome that activate their neurons at the same time
strengthen the synapses between them --- which is the cellular face of
association itself, the law of effect written in the language of ions and
receptors. The slug and the hippocampus together supply what the
behaviourists lacked and the cognitivists rarely sought: a mechanism, at
the level of the synapse, for the change that learning leaves behind.}

\section{The studying that feels like learning}

Return now to the night of cramming, because the gap it opened --- between
the feeling of having learned and the fact of it --- has been mapped by
research precise enough to change how anyone ought to study, and the
findings are almost perfectly counterintuitive.

Start with the most robust result in the entire field, old enough that
Hermann Ebbinghaus noticed it on himself in the 1880s: the \emph{spacing
effect}. The same amount of study produces dramatically more durable
learning when it is spread across time than when it is massed into a single
session. Two hours divided among four days beats two hours in one evening,
and not by a little. Cramming works in the narrow sense that it gets you
through tomorrow's exam; it fails in the sense that matters, the retention
of anything past it. And the cruel part is that massed study \emph{feels}
better while you do it --- the material, gone over and over in one sitting,
grows fluent and familiar, and that fluency is read by the mind as
mastery. It is not mastery. It is the fluency of the moment, which fades.
The feeling of learning and the fact of learning point in different
directions, and the studier trusts the feeling.

Then the \emph{testing effect}, demonstrated cleanly by Henry Roediger and
Jeffrey Karpicke: retrieving information from memory --- testing yourself,
struggling to produce the answer --- produces far more durable learning
than studying the same material for the same time. Reading a passage twice
feels more productive than reading it once and then trying to recall it,
and it is less productive, often by a wide margin, when retention is
measured a week later. The act of retrieval is not a neutral readout of
what you know; it is itself a learning event, one that strengthens the
very trace it draws on --- which is the reconsolidation insight of the
Memory chapter wearing work clothes. Pulling a memory out changes it, and
in this case changes it for the better, anchoring it more firmly than any
amount of passive review.

The psychologist Robert Bjork gathered these and others under a phrase
that captures the whole paradox: \emph{desirable difficulties}. Conditions
that make learning feel harder and slower while it is happening ---
spacing it out, testing instead of rereading, interleaving different kinds
of problem rather than blocking them together, varying the circumstances of
practice --- tend to produce learning that is more durable and more
flexible, more readily transferred to new situations. Conditions that make
learning feel fluent and fast --- massing, rereading, practising one type
of problem in a long uninterrupted block --- tend to produce learning that
looks good immediately and decays quickly. The difficulty is desirable
precisely because the struggle is the learning; the error you make and
correct, the retrieval you strain for, the prediction that fails and is
updated, is the prediction error of the chapter's first section, the thing
that actually moves the machinery. Smooth, easy, error-free study
generates no prediction error and therefore little learning, however
satisfying it feels.

\aside{Here is the same caution the Working Memory chapter raised, in a
new place. There, the dream was a brain-training drill that would expand
your raw mental capacity, and the finding was that practice transfers
narrowly: you get better at the drill and little else. The lesson rhymes.
What feels like learning --- the rising fluency of massed practice, the
rising score on a trained task --- is a poor guide to what has durably and
generally been learned. The performance you can see during training is not
the learning, and Tolman's rats made the same point ninety years ago: the
two come apart, and the visible one is the less trustworthy. We are, in
study as in so much else, the last to know what we have actually built.}

The interleaving result is the most surprising of all, because it pits the
truth against the learner's confidence most directly. Asked to learn to
recognize the styles of several painters, students who study one painter's
works in a block before moving to the next feel they are learning more ---
and, tested afterward on new paintings, do worse than students who saw the
painters' works shuffled together. The mixing forces a harder
discrimination, a sharper attention to what distinguishes one style from
another, and it produces better learning while feeling worse. When the
students are asked which method taught them more, most pick blocking, the
method that taught them less. The metacognitive illusion is nearly
total: the mind cannot feel its own learning accurately, and trusts the
sensation of ease over the fact of acquisition, every time.

\section{Why the maps disagree}

The field's reference works carve learning along fault lines that, by now,
read like a record of the disputes this chapter has walked through --- and
the sharpest of them concerns where learning sits in the order of
cognition, whether it is the floor or the whole building.

The educational tradition, in the taxonomy of objectives that structures
much of how teaching is organized, places learning's most basic act ---
\emph{remembering}, the retention and retrieval of what was taught --- at
the very bottom of its hierarchy, the foundational first of six rising
levels, beneath understanding, applying, analyzing, evaluating, and
creating. On this map, to have learned in the sense of retained is the
ground floor, the prerequisite the higher cognitive work is built upon ---
a humble base, the least of the things a mind does with knowledge. But the
psychometric and neuroscientific traditions draw the opposite picture.
For them, learning is not a floor beneath the interesting cognition; it is
the fundamental adaptive process \emph{itself}, the thing the nervous
system most essentially does, shared all the way down to a sea slug
withdrawing its gill --- the deepest and most general capacity there is,
on which everything else is a specialization. The identical phenomenon
sits at the bottom of one tradition's hierarchy as a mere precondition and
at the very root of another's as the master process. One tradition's floor
is the other tradition's foundation of all cognition.

This is, exactly, the shape the Reasoning chapter found when it set the
educational taxonomy's placement of \emph{inferring} --- filed near the
bottom, as a sub-skill of comprehension --- against the psychometric
tradition's placement of inference at the apex of fluid intelligence. The
same word names the penthouse in one hierarchy and the basement in
another, because the two traditions are measuring different things and
calling them by one name: the everyday automatic act and the rare
sophisticated one. Learning fractures the same way. The educator means the
retention of a taught fact; the neuroscientist means the synaptic
adaptation underlying all behavioural change; and these are genuinely
different referents wearing a single word, so the catalogues that file
learning cleanly in one place are losing half its meaning.

Beneath that disagreement sits the deeper one the chapter has circled from
the start --- the century-long quarrel between the associative and the
cognitive. Is learning, at bottom, one process: the adjustment of
associations by prediction error, the mechanism that runs in the slug and
the rat and the dopamine system, mechanical and universal and blind? Or is
it the building of structured internal models --- Tolman's maps, the
infant's grammar, the insight that assembles a solution whole --- a
constructive, representational achievement that the strengthening of bonds
can never capture? The reference works take sides, mostly without saying
so. A catalogue that lists learning as a single basic process,
co-equal with perception and memory and shared across the animal kingdom,
is siding with the associationists. A catalogue that splits it into
explicit and implicit, declarative and procedural, statistical and
insightful --- a family of distinct systems on distinct tissue --- is
siding, in effect, with the proposition that there is no one thing called
learning at all, only a set of different ways the brain installs different
kinds of change.

And the honest resolution, as so often in this atlas, is that the dispute
cannot be settled by collecting more data of the same kind, because the
two camps are not disagreeing about the facts but about which level is
fundamental. The associative mechanism is real --- prediction error
genuinely drives synaptic change, in the slug and in us, and the
Rescorla--Wagner arithmetic genuinely describes a great deal of what
learning does. The cognitive achievement is equally real --- the rat
genuinely builds a map it was never rewarded for building, the infant
genuinely extracts a structure no reinforcement could have stamped in.
These are not rival theories of one process; they are faithful
descriptions of different processes that the single word ``learning'' has
always been forced to cover. The maps disagree because learning, like
memory and attention and language before it, was never one thing, and the
contradiction in the catalogues is the more accurate record of that.

What is not in dispute is the strangeness the chapter opened on, and it is
worth ending there. You learn from your errors and not from your
successes, from the predictions that fail and not from the world that
confirms them --- which means the moments of being wrong, that feel like
failures of learning, are the only moments learning occurs. And you cannot
feel yourself doing it: the studying that feels like mastery is the
shallow kind, the practice that feels like floundering is the deep kind,
and the skill that is most thoroughly yours --- the balance that keeps you
upright on the bicycle --- is the one you can least put into words. We are,
in learning as in everything this atlas has touched, built to do the thing
and not to know that we are doing it, and the last to learn how we learn.













ch13-self and other.tex

\atlaschapter{Self and Other}{We treat the self as the one fact we cannot be wrong about, and empathy as the self reaching out toward others. The science nearly reverses both: the self is a model the brain builds, the reach toward others borrows the same machinery, and the boundary between them is drawn rather than given.}

Watch someone catch their hand in a slamming car door and you will wince
--- a small involuntary flinch, a tightening, perhaps a sympathetic ache
in your own hand, all of it arriving before you have decided anything. You
did not reason your way to the wince. It happened to you, the way the fear
in an earlier chapter happened, from somewhere below choice, and it
happened on behalf of a hand that is not yours, attached to a person whose
pain you cannot in any direct sense feel. Something in you reached across
the gap between two bodies and produced, faintly, in your own nervous
system, a version of an event occurring in someone else's.

That wince is the subject of this chapter, and it sits on top of two
assumptions so natural that stating them feels redundant. The first is
that there is a secure, bounded \emph{you} doing the wincing --- a self,
the one thing in all of experience you cannot be mistaken about, the
fixed point from which everything else is observed. The second is that
empathy is this self reaching outward: a stable inner observer extending
sympathy across a clear border toward a separate person on the far side.
Self first, secure and given; then the reach toward the other, secondary
and effortful.

The research of the last few decades has worked steadily against both
assumptions, and against the order they come in. Empathy turns out not to
be one reach but several, running on different circuitry, dissociable so
sharply that a person can have one and lack another --- and the most
visceral of them, the felt sharing of another's state, works by borrowing
the very machinery the brain uses to represent its own. Which raises the
question the chapter is really about. If feeling-with-others is the self's
machinery turned outward, what is that self? And the answer the science
has been reluctantly arriving at is that the self is no more a secure
given than the seen world or the felt emotion were: it is a model, built
by the brain, useful and imperfect, and the boundary it draws between you
and everyone else is something the brain constructs and can be made, in
the laboratory, to redraw. The previous chapter on reading other minds
asked how we infer what others believe. This one asks something stranger
and more intimate: how the brain builds the border between self and other
at all, and why that border is so much more porous than it feels.

\section{The empathy that isn't one thing}

Start where the last chapter left a gap. ``Reading other minds'' was
about \emph{knowing} --- inferring what someone believes, wants, intends,
the cool cognitive work of the mentalizing network. But knowing what
someone feels and \emph{feeling} it are not the same act, and the word
``empathy'' has long been forced to cover both, along with a third thing
that is neither. The first job of the science was to pull them apart.

The cleanest division runs three ways. There is \emph{cognitive empathy}
--- understanding another's emotional state, working out that they are
grieving or anxious or relieved. This is theory of mind pointed at
feelings rather than beliefs, and it is the capacity the previous chapter
anatomized. There is \emph{affective empathy} --- actually sharing the
state, feeling a version of the distress you perceive, the wince at the
slammed door. And there is something better called \emph{compassion} or
\emph{empathic concern} --- caring about the other's welfare, being moved
to help, which is a motivation rather than a perception or a shared
feeling and, as we will see, runs on a different system from either. You
can have any of these without the others, and the dissociations are not
subtle.

The mechanism of affective empathy is where the chapter's thesis first
shows itself, and it was made visible by Tania Singer in a study that
deserves its fame. She put people in a brain scanner and delivered, to
each of them, a brief painful stimulus, mapping the network of regions
that lit up --- the anterior insula and the anterior cingulate cortex
prominent among them, the parts of the pain system that, as the Pain
chapter argued, carry not the raw sensory signal but the felt,
affective, this-is-bad dimension of hurting. Then she did the
revealing thing. With the person still in the scanner, she gave the same
painful stimulus to their romantic partner, seated nearby, and signalled
to the scanned person that their loved one was now in pain. The affective
pain network activated again --- the same anterior insula, the same
cingulate --- in the brain of the person who was not being hurt. Watching
a loved one suffer, the brain ran a partial copy of suffering, using the
machinery it would use for its own. Affective empathy is not a message
about another's pain arriving from outside. It is the brain constructing,
in its own pain circuitry, a version of the other's state --- the
self's apparatus, borrowed.

\aside{The Pain chapter argued that pain is not a signal read off the body
but a construction the brain assembles, which is why it can be modulated,
faked by a phantom limb, or produced with no injury at all. Affective
empathy is the same constructive machinery aimed at someone else's body.
When you ache at another's injury, you are not receiving their pain; you
are building one, in your own circuits, from the evidence that theirs is
occurring. The wince at the car door is a controlled hallucination of
someone else's hurt --- which is exactly what you would expect of a brain
that builds its own pain the same way.}

That the three empathies are genuinely separate systems is proved, as so
often in this atlas, by the people in whom they come apart. Psychopathy is
the textbook dissociation: individuals who read other people's mental
states perfectly well --- often \emph{better} than average, the skill
turned to manipulation and deception --- while the affective sharing and
the concern are blunted or absent. The cognitive empathy is intact; one
can know precisely what another feels and use it as a lever, without the
felt resonance or the care that, in most people, comes attached. This is
the inverse of the profile too often assumed for autism, and here the
previous chapter's caution must be carried forward. The old story that
autism is a deficit of empathy conflated the three: the difficulty, where
it exists, is substantially in the \emph{cognitive} reading of
neurotypical social signals --- the effortful, less automatic mentalizing
the last chapter described --- and is frequently accompanied by affective
empathy that is intact or, by many accounts, heightened to the point of
being overwhelming. The double-empathy point returns with force: a
mismatch in cognitive reading between two differently wired minds is not
an absence of the capacity to feel with others. The two can vary
independently, because they are different systems, and a person may feel
another's distress acutely while struggling to infer its cause, or infer
the cause flawlessly while feeling nothing at all.

\section{When feeling-with goes wrong}

Here the chapter takes the turn each chapter of this atlas seems to
require --- the point at which a faculty everyone assumes is wholly good
turns out to be something more complicated. Affective empathy, the felt
sharing that we praise as the root of human decency, is not reliably good,
not reliably prosocial, and on one serious argument is a poor foundation
for morality altogether.

The first problem is internal to the mechanism. If feeling-with-others
works by constructing a version of their distress in your own affective
system, then empathizing with suffering means generating suffering in
yourself --- and a brain flooded with secondhand distress does not
reliably turn toward the sufferer. It often turns away. This is the
distinction, well documented in the helping literature, between
\emph{empathic concern}, which is other-directed and motivates approach
and aid, and \emph{personal distress}, which is self-directed --- the
aversive arousal of catching someone's pain --- and motivates escape from
the source of it. The nurse, the aid worker, the carer who burns out is
not deficient in empathy; they may be drowning in it, the shared distress
accumulating until withdrawal becomes self-protection. Feeling another's
pain, past a point, makes you flee the person in pain. The mechanism that
is supposed to bind us can, run too hard, drive us apart.

The second problem is that affective empathy is a spotlight, and
spotlights are biased by where they happen to point. This is the core of
the psychologist Paul Bloom's argument, made under the deliberately
provocative banner of being \emph{against empathy}: that empathy in the
felt, affective sense is a terrible guide to moral action precisely
because it is concrete, vivid, and innumerate. It fastens onto the single
identified victim and goes dark for the statistical thousands. It flows
more readily toward the near, the familiar, the similar, the
in-group, and dries up across distance and difference --- so that a
morality built on empathy is a morality warped toward whoever the
spotlight already favours. Bloom's positive proposal is that the felt
sharing should be demoted in favour of \emph{rational compassion} ---
a deliberate, principled concern for others' welfare that can be extended
evenhandedly to people the spotlight would never reach, that can do
arithmetic, that does not require you to manufacture distress in yourself
in order to act well. The argument is contested and not everyone accepts
it. But it lands a real blow against the sentimental view, and it does so
without trespassing on the moral-reasoning machinery an earlier chapter
already examined --- the trolley dilemmas and the intuitions that arrive
before their justifications. The question here is not how moral judgments
are made but whether the warm feeling we most trust as our moral compass
deserves the trust. Bloom's answer is that it does not, and the
neuroscience of the next section suggests he is at least pointing at
something real.

\aside{Empathy is this chapter's specimen of the atlas's recurring
pattern --- the genuine good inflated into a cure-all and then quietly
walked back. The popular literature made empathy the solution to cruelty,
prejudice, and conflict, the civilizational glue, the thing the world
simply needs more of, in the same register that made mirror neurons
``the neurons that shaped civilization'' and one gene \emph{the} gene
for language. The correction here is gentler, because empathy is real
and often does real good. But the strong claim --- that more felt empathy
reliably makes people kinder and the world better --- does not survive
contact with a mechanism that is biased toward the near, exhausting to
run, and as available to the manipulator as to the saint. The thing to
distrust, once again, is the clean and flattering version.}

\section{Empathy and compassion are different things}

If affective empathy can curdle into distress and bias, the obvious
question is whether there is anything in the social repertoire that does
the work we hoped empathy would do --- that moves us to help without
burning us down or playing favourites. The answer, from a striking line of
work by Singer and her colleague Olga Klimecki, is yes, and it is a
\emph{different system}, which is why empathy's failures are not the last
word.

They trained people, over weeks, in two different mental practices and
watched what each did to brain and behaviour. One was, in effect, empathy
training --- attending to and resonating with others' suffering. It
increased participants' negative affect and strengthened activity in
exactly the shared-distress network the first section described, the
anterior insula and anterior cingulate. Sharing more suffering made people
feel worse, as the mechanism predicts. The other practice was
\emph{compassion} training, in the tradition of loving-kindness
meditation --- cultivating warmth and concern and the wish for others to
be well, rather than the mirroring of their pain. This produced a quite
different signature. Negative affect did not rise; \emph{positive} affect
did, even when participants were contemplating others' distress. And the
brain activity moved to a different network altogether --- regions
associated with reward, affiliation, and care, the medial orbitofrontal
cortex and the ventral striatum and their neighbours, the circuitry of
warmth and approach rather than of shared pain. Crucially, the two could
be trained separately. Compassion was not more empathy; it was another
capacity, with its own neural home and its own opposite emotional tone,
and it could be strengthened on its own.

This is the resolution to the previous section's problem, and it reframes
the whole sentimental picture. The warmth that moves a person to help,
sustainably and without drowning, is not the felt copy of another's pain.
It is a distinct, trainable system of care and concern that can take the
shared distress as its input but need not be flooded by it --- the
difference, roughly, between a clinician who is wrecked by each patient's
suffering and one whose steady warmth lets them act all day without
collapse. Bloom's rational compassion and the meditator's loving-kindness
turn out to point at the same underlying fact: there is a route to caring
for others that does not run through manufacturing their pain in yourself,
and it has a different address in the brain. The reach toward the other
is not one thing. It is at least three --- knowing, feeling-with, and
caring --- and they can be told apart not only in behaviour and in the
clinic but in the tissue.

\section{The self that gets borrowed}

Return now to the thread the first section pulled and left hanging.
Affective empathy works by running another's state on the self's own
machinery --- your pain circuits firing for their injury, your affective
systems lending themselves to a model of someone else. But this presumes a
self whose machinery is available to be borrowed, a representation of
\emph{you} that the representation of \emph{them} can piggyback on. Pull
that thread and the chapter turns from the other to the self, and finds
that the self is no more a simple given than anything else in this atlas.

The most famous attempt to catch self-awareness in the act is the
\emph{mark test}, devised by Gordon Gallup. Place a mark --- a spot of
odourless dye --- somewhere on an animal's body that it can see only in a
mirror, then give it a mirror, and watch. An animal that treats the
reflection as another creature learns nothing about the mark. An animal
that touches the mark \emph{on its own body}, using the reflection to
guide its hand, has apparently understood that the thing in the glass is
itself --- that it has a self, locatable, inspectable from the outside.
Human children begin to pass at around eighteen months to two years. The
great apes pass; most other species do not, though a handful --- some
dolphins, an elephant, perhaps a magpie --- appear to, in results that are
still argued over. The mark test became, for a generation, the operational
definition of self-recognition.

\aside{The mark test is worth a caution of the kind this atlas keeps
having to give. Like the marshmallow that was supposed to measure
willpower and the false-belief task that was supposed to measure theory of
mind whole, a single clean behavioural test gets loaded with more than it
can carry. Passing the mark test shows that an animal can recognize its own
body in a mirror --- a fact about \emph{visual} self-representation, no
more. It is not the same as having a sense of one's own thoughts, a
narrative identity, or an inner life, and children in some non-Western
settings pass it late or inconsistently in ways that suggest the test is
partly cultural rather than a clean readout of a faculty. The self is not
one thing any more than empathy or memory was, and no mirror catches all
of it. The bodily self the mirror tests is only the ground floor.}

Above that ground floor, the self has turned out to have something like a
characteristic signature in the brain, discovered almost by accident. When
researchers looked at what the brain does during \emph{rest} --- when a
person in a scanner is given nothing to do --- they found not quiet but a
distinct network humming away, most active precisely when the mind is
\emph{least} engaged with the outside world. This \emph{default mode
network}, anchored in the medial prefrontal cortex and the midline
structures running back toward the parietal lobe, is the system that comes
on when attention turns inward: during mind-wandering, during the
recollection of one's own past, during thinking about oneself --- and,
tellingly, during thinking about \emph{other people}. For the default mode
network overlaps heavily with the mentalizing network of the previous
chapter. The brain regions most active when you reflect on yourself are
substantially the same ones most active when you model someone else. Self
and other do not occupy separate quarters of the brain. They share an
address.

That overlap is the neural form of the chapter's thesis, and it shows up
in the smaller findings too. The medial prefrontal cortex responds when
you judge whether a trait applies to you, and there is a long-known quirk
of memory --- the \emph{self-reference effect} --- whereby information
processed in relation to yourself is remembered better than the same
information processed any other way, as though the self were an especially
sticky filing category. The Memory chapter's reconstructive past is, in
large part, a past organized around a self that did and felt and witnessed
things. The self is not a separate module sitting apart from the rest of
cognition; it is woven through memory, through the modeling of others,
through the brain's idling default --- a way of organizing experience as
much as a thing that has experiences.

\section{The self is a model}

The Consciousness chapter proposed something that, extended one step,
delivers this chapter's deepest claim. There, Michael Graziano argued that
awareness is the content of a model the brain builds of its own attention
--- an \emph{attention schema}, a simplified self-description that the
control system uses to regulate itself, and that we mistake, from the
inside, for a direct acquaintance with experience. The brain models its
own processes, the argument went, the way it models its body, and the
model's content is what we take for an inner fact. Apply the same logic
one level up. The brain also builds a model of itself as a whole --- as a
bounded agent, a locus of experience, a thing with a body and a history
and a point of view --- and that model is the self. Not a soul-pearl at
the centre, not the secure given we take it for, but a construction: a
useful, mostly accurate, sometimes mistaken representation that the brain
assembles and then inhabits.

The proof that the self is a constructed model, and not a fact the brain
simply registers, is that you can change it in the laboratory in under a
minute. Sit a person at a table with one hand hidden behind a screen and a
rubber hand placed where their real hand would be. Stroke both the hidden
real hand and the visible rubber one, in synchrony, with two soft brushes.
Within a minute or two, most people begin to feel the rubber hand as
\emph{theirs} --- to feel the brushstrokes as occurring on the rubber
hand, to flinch when it is threatened, to mislocate their own hand toward
it. The brain, presented with a seen hand being stroked exactly when a
felt hand is being stroked, draws the obvious inference: this seen thing is
my hand. The \emph{rubber hand illusion} is the bodily self behaving
exactly as the Perception chapter said vision behaves --- as an inference
from the available evidence, the brain's best hypothesis about what belongs
to it, revisable the moment the evidence changes. The felt boundary of the
self is not a fixed fact of the body. It is a continuously computed guess,
and you can fool it with two paintbrushes.

\aside{The rubber hand is the self-chapter's hollow mask. There, a face
prior so dominated the visual evidence that you saw a concave mask pop
convex and could not unsee it, because perception is the brain's
hypothesis and the hypothesis won. Here, multisensory correlation
overrides the body's own map and you feel a rubber hand as your own flesh,
because bodily selfhood is the brain's hypothesis and the same logic
applies. Both illusions are windows onto the same machinery: a brain that
does not receive the world or the body but infers them, and that hands the
inference to you as a finished, seamless fact. The self is the most
intimate of those inferences, and the rubber hand shows it is an inference
all the same.}

The illusion scales up alarmingly. Using a head-mounted display fed by a
camera behind them, researchers have given people the experience of
\emph{owning a whole body seen from the outside}, or of floating behind
their own physical body --- laboratory-induced out-of-body experiences,
the felt location of the self pried loose from the flesh and parked a
metre away. And the brain region most implicated in this --- the
temporoparietal junction, where electrical stimulation can by itself
trigger the sensation of leaving one's body --- is, once again, a node of
the very network the previous chapter identified for representing
\emph{other} minds. The same patch of cortex that helps you model where
another person's perspective sits is involved in locating your own self in
your own body. Disturb it and the self comes unmoored. This is the
chapter's thesis in its hardest form: the machinery that places you inside
yourself and the machinery that lets you take the standpoint of another are
not merely similar, not merely neighbours, but overlapping --- so that
self-location and other-modeling are two operations of one system,
and the border between you and everyone else is drawn by a brain using the
same pen for both sides.

\section{Why the maps disagree}

The field's reference works handle the social and self-directed mind worse
than almost any other region of cognition, and by now the reasons are
predictable. The disagreements are not clerical. They record a domain in
which the field cannot decide how many things it is even looking at.

Empathy is the milder case. Some catalogues enter it as a single faculty
--- one heading, one capacity, the ability to share and understand others'
states. Others split it along exactly the seams this chapter drew, into
cognitive and affective and motivational components that dissociate in the
clinic and localize to different circuits. Still others decline to give
empathy its own territory at all, folding it into emotion, or into the
emotional side of theory of mind, or into the regulation of one's own
affect. A reference work that lists empathy once is making, without
saying so, the claim that it is one thing; a reference work that splits it
three ways is taking the opposite side; and the evidence of psychopaths
who read without feeling, carers who feel until they flee, and meditators
who can be taught to care without sharing pain says the splitters have the
better of it. Empathy names a committee, and the catalogues that treat it
as an individual are mapping a consensus the phenomenon does not support.

The self is the harder case, and it breaks the catalogues more deeply,
because the field cannot agree on whether the self is the kind of thing
that belongs in a catalogue at all. Some reference works give it no entry,
as though a taxonomy of cognition could be complete with no heading for the
one who does the cognizing --- and there is a logic to the omission, since
the self is implicated everywhere and isolable nowhere. Others file
\emph{self-awareness} as a subtype of consciousness, or \emph{self-knowledge}
as a province of memory and metacognition, or self-recognition as a
perceptual achievement caught by a mirror. The same self appears as a
mode of awareness in one map, a kind of knowledge in another, a
developmental milestone in a third, and a node of a resting brain network
in a fourth, and none of these is wrong and none contains the others.

Beneath all of it sits the fracture this atlas has met in every chapter,
arriving here at its most vertiginous object. Is the self a \emph{thing}
--- an entity, a faculty, a centre, the secure given we feel it to be, the
kind of item a catalogue could in principle list? Or is it a
\emph{process}, a \emph{construction} --- a model the brain builds and
inhabits, no more a discrete object than ``the seen world'' or ``the felt
emotion'' are discrete objects, and therefore not a thing to be entered in
a taxonomy but an activity that taxonomy keeps trying and failing to pin?
This is the had-versus-made disagreement that divided the memory
taxonomists from the memory architects, the basic-emotion theorists from
the constructionists, the consciousness-as-threshold camp from the
consciousness-as-property camp --- and it is sharpest here, because the
thing in dispute is not memory or fear or awareness but the one who
remembers and fears and is aware. A catalogue that grants the self a
heading is betting it is a thing. A catalogue that dissolves it among
memory and consciousness and social cognition is betting it is made. And
the rubber hand, the out-of-body display, the default network shared with
other minds, the pain circuits lent out to a stranger's injury, all point
the same way the rest of this atlas has pointed: toward made.

Which returns us, finally, to the wince at the slammed door, and to how
much it contained. There is no secure inner self reaching across a clear
border toward a separate other. There is a brain that builds a model of a
body and calls it ``me,'' builds models of other bodies and other minds on
the same machinery and in overlapping tissue, and draws the boundary
between them as a working hypothesis rather than a given fact --- a
boundary it can be made to redraw with two paintbrushes, or move a metre
into empty air, or extend, in the half-second of a wince, into the body of
someone whose pain it has just constructed in its own circuits as if it
were its own. We take the border between self and other to be the most
basic fact there is, the one line in experience that draws itself. It is
the last and most intimate of the things the brain builds without telling
us, and, as everywhere in this atlas, we are the last to know that we are
the ones who drew it.














ch14-the minds eye.tex

\atlaschapter{The Mind's Eye}{Picture an apple, and a private theatre lights up --- or, for some readers, it does not. Mental imagery is perception run backward, the eye's machinery fired without the eye. The oldest dispute is whether the picture it makes does real work or is a story the mind tells about itself; the newest is how wildly, and how invisibly, the picture varies from one person to the next.}

How many windows are there in the home you grew up in? Almost nobody knows
the number as a fact. To answer, you do something peculiar: you go back
there, in your head, and you \emph{look} --- moving through the rooms,
turning to the walls, counting the windows as you find them, as though the
house were standing in front of you and you had only to inspect it. The
number was never stored. It is generated, on demand, by building a kind of
model of the house and examining it with something that behaves
remarkably like sight. Ask instead whether a German shepherd's ears stand
up or flop down, or which way a galloping horse's legs are arranged at
full stretch, and you will likely catch yourself doing the same thing:
calling up an image and reading the answer off it, consulting a picture
that is, in some sense, not there.

That act --- the summoning of a quasi-perceptual experience in the absence
of the thing perceived --- is mental imagery, and it is one of the most
useful and most slippery capacities the mind has. Useful, because it lets
you rehearse a route before walking it, try a sofa against a wall before
carrying it across the room, design a bridge that does not yet exist,
relive a face that is gone. Slippery, because the moment you try to say
what a mental image actually \emph{is} --- a picture? in what? seen by
whom? --- the whole thing threatens to dissolve into paradox. For most of
the twentieth century the slipperiness won, and serious psychology refused
to study imagery at all. The behaviourists of the Learning chapter had
banished it as the very type of the unobservable, unscientific
inner event --- you cannot measure another person's mental picture, so it
was ruled out of bounds, despite Aristotle's old conviction that the soul
never thinks without one. It took a single, beautifully simple experiment
to bring imagery back, and the way it did so points at everything this
chapter is about.

Because here is the destabilizing part, the part that turns a tidy topic
into one of the strangest in the atlas. Whatever just happened in your
head when you counted those windows --- whether a vivid picture bloomed,
full of colour and detail, or something fainter and more diagrammatic, or
nothing visual at all, just a sense of the answer arriving --- you have
almost certainly assumed your entire life that it happens, more or less
the same way, in everyone. You are almost certainly wrong. The private
theatre behind your eyes runs a production that may be radically unlike
the one running behind the eyes of the person next to you, and because no
one can climb inside anyone else's, the difference can go a whole lifetime
undetected. The mind's eye is the most private organ there is, and the
science of it has had to find ways to measure a picture that only its
owner can see.

\section{The turning figures}

The experiment that made imagery respectable again, in 1971, asked people
to do something so concrete that even a behaviourist could measure it.
Roger Shepard and Jacqueline Metzler showed their subjects pairs of
drawings of three-dimensional figures --- little structures of cubes
stuck together into bent, armlike shapes --- and asked a single question:
are these two drawings the same object, just rotated, or are they mirror
images, two different shapes that can never be turned into one another?
The task is easy in the sense that people get it right; the interesting
thing is how long they take.

What Shepard and Metzler found is one of the cleanest results in
psychology. The time it took to decide rose in an almost perfectly
\emph{straight line} with the angle between the two figures. If one drawing
was rotated forty degrees from the other, people answered quickly; at
eighty degrees they took about twice as long; at a hundred and twenty,
three times; and so on, the reaction time climbing in lockstep with the
angular difference, as though people were mentally rotating one figure to
bring it into alignment with the other at a roughly constant speed, and
the further they had to turn it, the longer it took. The linearity is the
whole point. A symbolic, look-it-up process would not have to pass through
every intermediate angle; it could in principle jump straight to the
answer. But a process that crawls through the orientations in between ---
forty-five degrees, then fifty, then fifty-five, the figure swinging
around in some internal space --- would take time proportional to the
distance travelled, which is exactly what the data showed. People seemed
to be turning something in their heads, and the something behaved like a
rigid object obeying the geometry of real rotation.

This was the first hard evidence for what came to be called an
\emph{analog} or \emph{depictive} representation --- a mental image that
preserves the spatial structure of what it depicts, the way a picture
does, rather than merely describing it in symbols, the way a sentence
does. And Stephen Kosslyn, more than anyone, built the case out. In one
experiment he had people memorize a map of a fictional island dotted with
landmarks --- a hut, a tree, a well --- and then, with the map gone, asked
them to imagine a black speck zipping from one landmark to another across
their mental image. The time it took scaled with the actual distance
between the landmarks on the original map: people took longer to scan
across the image than across a short stretch of it, as though the mental
map preserved the real map's distances and the speck had to traverse them.
In another, he found that people answer questions about the details of an
imagined animal faster when they picture it large than when they picture
it small --- as if, to check whether a mouse has whiskers, you must first
zoom in, the details of a small image being harder to see, just as they
would be in a real scene viewed from too far away.

\aside{The natural home for all of this, in the architecture of an earlier
chapter, is Baddeley's \emph{visuospatial sketchpad} --- the ``inner eye''
of the Working Memory chapter, the component that holds and manipulates
visual and spatial material. Mental rotation is the sketchpad doing work
you can time; image-scanning is the sketchpad holding a layout you can
traverse. The imagery debate this chapter is about is, in one sense, an
argument over what the sketchpad is actually made of --- whether the inner
eye looks at something genuinely picture-like, or whether ``picture'' is
just the most natural word for a process that has nothing pictorial in it
at all.}

From these results Kosslyn drew a strong conclusion: that mental images are
represented in a spatial medium --- a ``visual buffer'' --- in which they
are laid out and can be scanned, zoomed, and rotated, and that this medium
is, quite literally, picture-like, depicting rather than describing. The
mind's eye, on this account, looks at something that has parts arranged in
space the way the parts of a real scene are arranged in space. It was a
bold and concrete claim, and it drew an equally concrete attack.

\section{No pictures in the head}

The objection, pressed for decades and with great force by Zenon Pylyshyn,
begins with a question that sounds naive and is in fact lethal: if there is
a picture in your head, who is looking at it? A picture is only a picture
to an eye that sees it. To explain how you understand a scene by saying
there is a picture of it inside you merely relocates the problem, because
now something inside you has to look at the inner picture and understand
\emph{that} --- and so you have smuggled a little viewer into the head, a
homunculus with its own little eye, and you must now explain how
\emph{it} sees, presumably with a still smaller picture and a still
smaller viewer, on and on forever.

This is precisely the hazard the Executive Function chapter spent its
length escaping --- the homunculus, the little person installed to do the
work, who explains nothing because he needs his own homunculus inside him.
There the ghost was a controller; here it is a viewer. And Pylyshyn's
charge was that the whole depictive picture of imagery quietly depends on
him. The feeling that there is a picture before the mind's eye is vivid and
universal, but a feeling is not evidence about the format of a
representation, and the underlying representation, he argued, is not
pictorial at all. It is \emph{propositional} --- a structure of
symbols and relations, description rather than depiction, more like the
sentence ``the hut is northeast of the tree, a long way off'' than like a
map. There is no picture; there is a body of knowledge, stored in some
abstract symbolic code, and the sense of seeing is a kind of user
illusion laid over it.

But what, then, of Kosslyn's beautiful scaling results --- the scan times
that track distance, the rotation times that track angle? Here Pylyshyn
made his sharpest move. Those results, he argued, do not reveal the
structure of the image; they reveal what the subject \emph{knows} about how
the corresponding real situation would unfold. You scan your mental map at
a rate proportional to distance because you know, tacitly, that scanning a
real scene takes longer over longer distances, and you unconsciously
reproduce that timing --- not because anything in your head is actually
travelling across an actual space. The phenomenon is \emph{cognitively
penetrable}: change what the subject believes the task is about, or what
they expect should happen, and the ``image'' obligingly behaves
differently. A genuine picture would not care what you believed; its
geometry would be fixed. The fact that the imagined results bend to the
subject's expectations suggested, to Pylyshyn, that there was no fixed
geometry there at all --- only a person simulating, from general
knowledge, how a real scene would behave, and experiencing that simulation
as a picture. On this account the depictive quality of imagery is
\emph{epiphenomenal}: real as an experience, but doing none of the
computational work, a coat of subjective paint over a propositional engine.

The debate that followed ran for thirty years and generated some of the
most ingenious experiments in cognitive science, each side parrying the
other, and for a long time it had the frustrating character of a dispute
that more data could not settle --- because the two theories, cleverly
enough specified, could each accommodate almost any behavioural result.
What was needed was a way to look past behaviour, into the machinery
itself, and ask the question the armchair could not: when a person forms a
mental image, what is the brain actually \emph{doing}, and does it look
more like a picture or more like a sentence?

\section{Perception running backward}

The answer, when the brain-imaging tools of the 1990s and after could
finally supply one, came down substantially --- though not cleanly --- on
Kosslyn's side, and in doing so it revealed something deeper than either
party had quite framed. When you form a visual mental image, the parts of
your brain that switch on are, to a striking degree, the same parts that
switch on when you actually \emph{see}.

Imagine a face and the fusiform face area --- the patch of visual cortex
that responds to faces during ordinary perception --- becomes active.
Imagine a place, a building or a landscape, and the parahippocampal place
area, the perception system's scene-detector, lights up instead. The brain
recruits, for imagining a thing, the very specialized machinery it uses for
seeing that kind of thing. And the result goes deeper than these
high-level, category-specific regions. Imagery activates \emph{early}
visual cortex, including the primary visual area at the very back of the
brain where signals from the eye first arrive --- and it does so in a way
that respects that region's most pictorial property. Early visual cortex
is laid out \emph{retinotopically}: it is a literal map, with neighbouring
patches of the visual field handled by neighbouring patches of cortex, and
larger images occupying more of the map. When people imagine larger
objects, more of this retinotopic cortex is engaged, in proportion to the
imagined size --- the mental image is splashed across the cortical map at
a scale that matches what you are picturing. This is about as close to a
literal picture in the head as neuroscience could hope to find: not a
picture for an inner eye to see, but a pattern of activity, in
topographically organized visual tissue, that preserves the spatial
structure of the imagined scene. And when researchers temporarily
disrupted early visual cortex with magnetic stimulation, imagery was
impaired --- evidence that this depictive activity is not a useless
by-product but is doing genuine work.

\aside{The honest caveat, in the atlas's habit, is that this did not
deliver a knockout. The involvement of the earliest visual areas in
imagery is real but variable --- stronger in some people and some tasks
than others --- and a handful of patients have lost early visual cortex
yet kept their imagery, while others have lost imagery from damage
elsewhere, in the temporal and parietal lobes. The format question, a
determined propositionalist can still insist, is not strictly settled by
showing that imagery uses visual machinery. But the weight of the
evidence shifted the burden decisively. After the imaging work, ``there is
no picture-like representation anywhere in imagery'' became a much harder
position to hold than ``the depictive story is roughly right, with
propositional knowledge doing real work alongside it.''}

The deeper thing the imaging revealed is the unifying idea of this
chapter, and it reaches straight back to the Perception chapter. There,
perception was recast as inference --- not a passive reception of the
world but the brain's active, top-down hypothesis about what is out there,
predictions flowing \emph{down} from higher areas to meet and explain the
sensory evidence flowing up. Anil Seth's phrase for it was
\emph{controlled hallucination}: a perception is a hallucination the
incoming signal happens to constrain. Mental imagery, the brain data
suggest, is that same generative machinery running in the other direction
--- the top-down predictions fired off \emph{without} the bottom-up
sensory evidence to rein them in. Perception is the generative model
checked against the world; imagery is the generative model let off the
leash, producing the perceptual representation from the inside, with the
eyes contributing nothing. This is why imagery and seeing share so much
tissue: they are not two systems but one system in two modes, with and
without the correcting torrent from the senses. To imagine is to perceive
with the world subtracted.

That single reframing tidies up a surprising amount. \emph{Dreaming} is
imagery at its most unconstrained --- the generative model running in full
flood while the gates to sensory input are closed and the corrective
signal cannot get in, which is why dreams are so vivid and so unmoored
from reality. \emph{Hallucination}, in psychosis or in sensory
deprivation, is imagery that has slipped its label --- a top-down
construction mistaken for a bottom-up perception, the brain's own
generated image misfiled as a report from the world, exactly the
reality-monitoring failure the Perception and Consciousness chapters
circled. The same machine that lets you helpfully picture an absent apple
will, when its outputs go uncorrected or unlabelled, dream you a world at
night and, in illness, conjure voices and visions by day. The mind's eye
and the dreaming brain and the hallucinating one are running the same
program at different settings.

\section{Remembering and imagining}

If imagery is the machinery of perception turned to the absent, then its
reach extends far past the imagined apple, and one of its most important
uses was named, without quite being recognized as imagery, back in the
Memory chapter. There, Endel Tulving's account of episodic memory carried a
striking coda: that recalling a specific past event is a kind of
\emph{mental time travel}, a re-inhabiting of a moment from the inside ---
and that the very same faculty might run forward as well as back, the
capacity to remember a particular past and the capacity to imagine a
particular future being, perhaps, two uses of one machine. This chapter is
where that suggestion comes home, because the machine in question is the
mind's eye.

To remember sitting on a particular beach and to imagine sitting on a
beach you have never visited are, it turns out, deeply similar acts, drawing
on overlapping brain systems and the same constructive, image-building
process. The psychologist Daniel Schacter, whose catalogue of memory's
distortions the Memory chapter drew on, proposed that this is no accident:
that memory is reconstructive --- assembling each recollection afresh from
stored fragments rather than replaying a record --- precisely \emph{because}
its deeper job is not faithful playback but flexible simulation. A system
that can recombine the pieces of past experience into a coherent imagined
scene is a system that can picture futures, weigh options, rehearse the
not-yet-real. The reconstructive memory that makes us unreliable witnesses
is the very same capacity that lets us imagine at all; the bug and the
feature are one mechanism.

The proof came, as it so often does in this atlas, from people who had lost
the part in question. Patients with damage to the hippocampus --- the
structure whose loss left the Memory chapter's H.M. unable to lay down new
recollections --- turn out to be impaired not only at remembering their
pasts but at \emph{imagining} novel experiences. Asked to describe a scene
they might encounter --- a sandy beach, a busy market --- such patients
produce something thin and fragmentary, lacking the rich spatial coherence
that healthy people supply without effort, unable to bind the elements into
a unified imagined place. The faculty that reconstructs the remembered past
and the faculty that constructs the imagined future failed together,
because they were never two faculties. The mind's eye does not only look
at what is absent in space, the apple on the table that is not there. It
looks at what is absent in \emph{time} --- the beach of last summer, the
meeting tomorrow, the life not lived --- and it builds those, too, from
the same depictive parts.

\aside{This is also where ``imagery'' and ``imagination'' --- words the
reference catalogues constantly tangle --- come briefly apart. The
quasi-perceptual re-experiencing of something (the apple, the remembered
beach) is one thing; the \emph{generative} act of combining stored
elements into something genuinely new (a market you have never seen, a
creature that does not exist, a bridge not yet built) is another, and the
spatial \emph{manipulation} of imagined material (rotating the block
figure, planning a route) is arguably a third. They share the visual
machinery and they share a vocabulary, and they are not the same capacity
--- a distinction that will return when the maps are laid side by side.}

\section{The blind mind's eye}

Everything so far has assumed that you, the reader, have a mind's eye to
turn this way and that. For a portion of people --- something like one in
twenty, by current estimates, though the figure is soft --- that assumption
is simply false, and the discovery that it is false is among the most
quietly astonishing findings in the modern science of the mind.

The condition now called \emph{aphantasia} is the absence of voluntary
visual imagery. Asked to picture an apple, a person with aphantasia
pictures nothing --- not a faint apple, not a dim one, but no visual
experience at all. They know perfectly well what an apple looks like; they
can tell you it is round and red and has a stem; they can recognize one
instantly. But the act the rest of us perform when we ``see it in the
mind's eye'' produces, for them, no image whatsoever, only the knowledge
without the picture. At the other extreme are people with
\emph{hyperphantasia}, whose voluntary imagery is as vivid and detailed as
actual seeing, a mental picture barely distinguishable from the real thing.
And between these poles runs a vast, continuous range --- captured, crudely,
by a questionnaire that asks people to rate the vividness of imagined
scenes, on which the population spreads itself from ``perfectly clear, as
vivid as normal vision'' to ``no image at all, you only know you are
thinking of the object.'' This is not a small difference of degree around a
common centre. It is one of the widest documented ranges of any feature of
human experience, and most people occupying it have no idea where on it
they sit, because they have never been able to compare.

The thing was glimpsed long ago and then lost. Francis Galton, in 1880,
sent around a questionnaire asking people to picture their breakfast table
and rate the image, and was startled to find that some respondents ---
including, to his bafflement, several distinguished scientists --- reported
no imagery at all and were puzzled that anyone took the language of mental
pictures literally, while others described images of photographic richness.
Galton had stumbled onto the entire range. But the finding sank, the
behaviourist freeze descended, and it was not until 2015 that the
neurologist Adam Zeman gave the blind end of the spectrum its name and its
modern study, after a patient who had lost his vivid imagery following a
medical procedure led Zeman to discover the many people who had simply
never had it. Naming a thing is how science gets to study it, and once
``aphantasia'' existed, the people who had it --- many of whom experienced
a small shock on first learning that ``picture it in your mind'' was not a
metaphor for everyone else --- came forward in numbers.

\aside{It is tempting, and it has been popular, to leap from the reality of
imagery variation to the idea of fixed ``visual learners'' who should be
taught with pictures and ``verbal learners'' with words --- the
\emph{learning styles} doctrine that saturates education. The leap is not
supported. People do differ, sometimes enormously, in the vividness of
their imagery; but study after study has failed to find that matching
instruction to a person's supposed style improves their learning, and the
Learning chapter's deeper point applies --- what feels like it should help
is a poor guide to what does. Real variation in the mind's eye is one of
the firmest findings in this chapter. The educational prescription built on
top of it is one of the era's most durable myths. The variation is true;
the tidy use everyone wants to make of it is not.}

What makes aphantasia matter for this chapter's central debate, and not
merely as a curiosity, is what aphantasic people can and cannot do. By and
large they live unremarkable cognitive lives. Many work in visual and
spatial professions; many perform normally on tests of spatial reasoning,
and a fair number do the mental-rotation task at ordinary speed --- getting
the right answer, apparently by some route that does not pass through a
conscious picture. Most striking of all, many people with no voluntary
imagery nonetheless \emph{dream} in vivid images, the involuntary,
bottom-up flood intact while the voluntary, top-down summoning is absent.
And the absence has measurable downstream consequences: when people are
made to read a frightening written scenario, ordinary imagers show a rise
in the body's stress response --- the skin's conductance climbing as they
picture the threat --- while aphantasic readers, who cannot picture it, show
a markedly flatter response. The image, it turns out, is part of how a mere
description reaches into the body and frightens it; remove the image and
the words lose some of their grip.

These dissociations cut to the heart of the old debate. They suggest that
the vivid, conscious \emph{picture} --- the thing the depictive theory
made central and the propositional theory dismissed --- is at least partly
separable from the spatial \emph{computation} that imagery tasks require.
An aphantasic person can rotate the block figure correctly while
experiencing no picture of it turning, which means the picture and the
problem-solving are not the same thing. The felt image may be the
\emph{experience} of a process rather than the process itself --- which
grants Pylyshyn something real about the picture's dispensability, while
granting Kosslyn the reality and the cortical machinery of the picture
where it does occur. Both were partly right, and aphantasia is the
condition that shows how, by prying apart what the debate had assumed was a
single thing.

\section{Why the maps disagree}

The field's reference works carve the imagining mind in ways that, by now,
have a thoroughly familiar shape --- and the disagreements track exactly
the fault lines this chapter has walked.

The first disagreement is the one the asides kept flagging: whether imagery
is one capacity or several. Some catalogues, built to \emph{measure}
abilities that differ between people, carefully separate the strands --- 
the generative power to produce novel images filed as one narrow ability,
the capacity to manipulate visual material in space, to rotate and scan and
navigate, filed as a distinct one, both tucked under a broad visual-
processing heading. Others, built from the bottom up out of described
concepts, name the quasi-perceptual experience of imagery several times
over in nearly identical terms and quietly fold the generative act in with
it, treating imagining-a-thing-you've-seen and imagining-a-thing-that-never-
was as one entry. One tradition splits imagery from imagination from
spatial transformation; another runs them together; and the dissociations
--- the aphantasic who cannot picture but can rotate, the patient who can
perceive but cannot generate a scene --- say the splitters are tracking
something real that the lumpers miss.

Beneath that lies the format question itself, the thirty-year war over
whether the image is a picture or a description --- and the catalogues take
sides without announcing it, in something as simple as where they file
imagery at all. To place mental imagery under \emph{visual processing},
beside perception and visual memory, is already to bet that it is
depictive, picture-like, a creature of the visual system. To place it
instead among the abstract reasoning faculties, or to treat it as a mode of
thought that merely happens to feel visual, is to lean propositional. The
shelf a reference work chooses encodes a theory of what an image is made
of, usually without the makers having meant to commit to one.

And beneath \emph{that} sits the deepest fracture, the one this atlas has
met under different names in every chapter. Is a mental image a
\emph{thing} the mind contains --- an inner picture, an object before the
mind's eye, the kind of item a catalogue could list and a viewer could
inspect? Or is it a \emph{process} the mind runs --- the perceptual system
operating in generate-mode, a pattern of activity with no picture and no
viewer, the felt image being merely what that process is like from the
inside? This is the had-versus-made disagreement that split the memory
taxonomists from the memory architects and the constructed self from the
given one, arriving here in the inner theatre. The depictive theory leans
toward the thing; the propositional theory and the perception-run-backward
reframing lean toward the process; and aphantasia, once again, is the
evidence that tips the balance --- for if the vivid picture can vanish
entirely while the spatial work goes on, then the picture was never the
thing doing the work. It was the experience of the process, separable from
it, present in most of us and absent in some, and either way not the
engine but the engine's glow.

Which returns us, last of all, to the window you counted in a house that is
not there, and to how much that small act contained. There is no picture in
your head, in the sense the word first suggests --- no photograph, no
viewer, no little eye. There is a brain that can run its perceptual
machinery without the world plugged in, generating from the inside a
structured pattern that carries the spatial bones of a scene, and handing
the result to awareness as something that feels, to most people, like
seeing. With it you revisit the past and rehearse the future, try the
sofa against the wall, frighten yourself with a written scene, design the
bridge before the first girder is cast. And whether that handing-over
arrives as a vivid private film or as a bare, imageless knowing is one of
the largest differences there is between two human minds --- a difference
that runs straight through the centre of everyone's experience and that,
because the theatre is sealed and admits no second viewer, most of us will
live our whole lives without ever knowing we are on one side of or the
other. The mind's eye is the organ we are most sure we share, and the one
we can least confirm we do.
